\documentclass[11pt,a4paper]{article}

\usepackage[margin=2.3cm]{geometry}
\usepackage{fontspec}
\usepackage{amsmath,amssymb}
\usepackage{graphicx}
\usepackage{booktabs}
\usepackage{tabularx}
\usepackage{array}
\usepackage{siunitx}
\usepackage[table]{xcolor}
\usepackage{caption}
\usepackage{subcaption}
\usepackage{enumitem}
\usepackage{float}
\usepackage{adjustbox}
\usepackage[hidelinks]{hyperref}
\usepackage[capitalise,noabbrev]{cleveref}

\captionsetup{font=small,labelfont=bf,margin=0.5em}
\setlist{itemsep=2pt,topsep=4pt}
\setlength{\emergencystretch}{3em}
\renewcommand{\topfraction}{0.9}
\renewcommand{\bottomfraction}{0.8}
\renewcommand{\textfraction}{0.07}
\renewcommand{\floatpagefraction}{0.75}
\setcounter{topnumber}{3}
\setcounter{totalnumber}{5}
\renewcommand{\arraystretch}{1.12}
\newcolumntype{Y}{>{\raggedright\arraybackslash}X}

\definecolor{okgreen}{HTML}{1B7F3B}
\definecolor{limamber}{HTML}{B36B00}
\definecolor{failred}{HTML}{B3261E}
\newcommand{\works}{\textcolor{okgreen}{\textbf{works}}}
\newcommand{\limited}{\textcolor{limamber}{\textbf{limited}}}
\newcommand{\fails}{\textcolor{failred}{\textbf{fails}}}

\newcommand{\claim}{\par\smallskip\noindent\textbf{Claim.}~}
\newcommand{\evidence}{\par\smallskip\noindent\textbf{Evidence.}~}
\newcommand{\reading}{\par\smallskip\noindent\textbf{Interpretation.}~}
\newcommand{\caveat}{\par\smallskip\noindent\textbf{Caveat.}~}

\newcommand{\code}[1]{\texttt{#1}}
\newcommand{\fig}[1]{figures/#1.pdf}
\makeatletter\let\tablebody\@@input\makeatother % expandable input for tabular bodies

% Generated by stress/numerics/report/make_numbers.py -- do not edit.
% mode: full; generated 2026-10-07 18:01:42
% sources: e01=full(8), e02=full(33), e03=full(22), e04=full(28), e05=full(8), e06=full(98), e07=full(92), e08=full(24), e09=full(26), e10=full(10), e11=full(53)

% source: e01_map_throughput crossover 'raw': GPU median never below CPU median up to hi
\newcommand{\NumGpuCrossoverNRaw}{none up to \ensuremath{N = 10^{7}}}
% source: e01_map_throughput crossover 'map_batch': bisection in log N of median GPU/CPU time
\newcommand{\NumGpuCrossoverNMapBatch}{\ensuremath{4.1\times10^{6}}}
% source: e01_map_throughput sweep numpy vs cupy_e2e: log-log interpolated ratio = 1
\newcommand{\NumGpuCrossoverNSweepRaw}{n/a}
% source: e01_map_throughput sweep map_batch_cpu vs map_batch_gpu: log-log interpolated ratio = 1
\newcommand{\NumGpuCrossoverNSweepMapBatch}{\ensuremath{6.2\times10^{5}}}
% source: e01_map_throughput sweep numpy: max over N of N / median call time (points/s)
\newcommand{\NumMapPeakNumpy}{\ensuremath{9.2\times10^{8}}}
% source: e01_map_throughput sweep cupy_e2e: max over N of N / median call time (points/s)
\newcommand{\NumMapPeakCupyEndToEnd}{\ensuremath{1.7\times10^{8}}}
% source: e01_map_throughput sweep cupy_compute: max over N of N / median call time (points/s)
\newcommand{\NumMapPeakCupyCompute}{\ensuremath{1.6\times10^{9}}}
% source: e01_map_throughput sweep map_batch_cpu: max over N of N / median call time (points/s)
\newcommand{\NumMapPeakMapBatchCpu}{\ensuremath{7.2\times10^{8}}}
% source: e01_map_throughput sweep map_batch_gpu: max over N of N / median call time (points/s)
\newcommand{\NumMapPeakMapBatchGpu}{\ensuremath{1.6\times10^{8}}}
% source: e01_map_throughput cupy_transfer / cupy_e2e median at the largest N = 10000000
\newcommand{\NumGpuTransferFractionAtMaxN}{88\%}
% source: e02_gpu_end_to_end timed: max over growth steps (CPU median >= 10 ms) of median CPU / median GPU step time
\newcommand{\NumGpuStepSpeedupBestKTwoEight}{\ensuremath{1.007\times}}
% source: e02_gpu_end_to_end timed: median over repeats of CPU / GPU summed stage time (growth through iterate table)
\newcommand{\NumGpuTotalSpeedupMedianKTwoEight}{\ensuremath{0.999\times}}
% source: e02_gpu_end_to_end timed: max over repeats of CPU / GPU summed stage time
\newcommand{\NumGpuTotalSpeedupBestKTwoEight}{\ensuremath{1.004\times}}
% source: e02_gpu_end_to_end profile, CPU: map-category tottime / all tottime
\newcommand{\NumMapFractionKTwoEight}{0.0171\%}
% source: e02_gpu_end_to_end profile, CPU: 1 / (1 - f_map)
\newcommand{\NumAmdahlBoundKTwoEight}{\ensuremath{1.000\times}}
% source: e02_gpu_end_to_end timed: max over growth steps (CPU median >= 10 ms) of median CPU / median GPU step time
\newcommand{\NumGpuStepSpeedupBestPThree}{\ensuremath{0.953\times}}
% source: e02_gpu_end_to_end timed: median over repeats of CPU / GPU summed stage time (growth through iterate table)
\newcommand{\NumGpuTotalSpeedupMedianPThree}{\ensuremath{0.975\times}}
% source: e02_gpu_end_to_end timed: max over repeats of CPU / GPU summed stage time
\newcommand{\NumGpuTotalSpeedupBestPThree}{\ensuremath{0.991\times}}
% source: e02_gpu_end_to_end profile, CPU: map-category tottime / all tottime
\newcommand{\NumMapFractionPThree}{0.505\%}
% source: e02_gpu_end_to_end profile, CPU: 1 / (1 - f_map)
\newcommand{\NumAmdahlBoundPThree}{\ensuremath{1.005\times}}
% source: e02_gpu_end_to_end timed: max over growth steps (CPU median >= 10 ms) of median CPU / median GPU step time
\newcommand{\NumGpuStepSpeedupBestKTen}{\ensuremath{0.916\times}}
% source: e02_gpu_end_to_end timed: median over repeats of CPU / GPU summed stage time (growth through iterate table)
\newcommand{\NumGpuTotalSpeedupMedianKTen}{\ensuremath{0.958\times}}
% source: e02_gpu_end_to_end timed: max over repeats of CPU / GPU summed stage time
\newcommand{\NumGpuTotalSpeedupBestKTen}{\ensuremath{1.002\times}}
% source: e02_gpu_end_to_end profile, CPU: map-category tottime / all tottime
\newcommand{\NumMapFractionKTen}{0.84\%}
% source: e02_gpu_end_to_end profile, CPU: 1 / (1 - f_map)
\newcommand{\NumAmdahlBoundKTen}{\ensuremath{1.008\times}}
% source: e02_gpu_end_to_end profile, CPU: max f_map over cases
\newcommand{\NumMapFractionMax}{0.84\%}
% source: e02_gpu_end_to_end profile, CPU: max 1/(1-f_map)
\newcommand{\NumAmdahlBoundMax}{\ensuremath{1.008\times}}
% source: e02_gpu_end_to_end verify: max over cases of the final max |CPU - GPU| unstable coordinate
\newcommand{\NumGpuDeviationMax}{\ensuremath{4.5\times10^{-13}}}
% source: e02_gpu_end_to_end verify: CPU and GPU point counts equal at the last step in every case
\newcommand{\NumGpuCountsEqual}{no}
% source: e02_gpu_end_to_end mbp: min_batch_points with the smallest k28 unstable growth time (1e9 = never on GPU)
\newcommand{\NumMinBatchBest}{1\,024}
% source: e02_gpu_end_to_end mbp: growth time at 'never' / at the best setting
\newcommand{\NumMinBatchBestGain}{\ensuremath{1.011\times}}
% source: e03_scaling growth k=2.8 cutoff 1e-07: log-log slope of step time vs unstable points (>= 1e3 points)
\newcommand{\NumStepTimeExponentKTwoEight}{1.07}
% source: e03_scaling growth k=2.8 cutoff 1e-07: slope of (RSS - baseline) vs points over the top two decades
\newcommand{\NumBytesPerPointKTwoEight}{1\,069}
% source: e03_scaling growth k=2.8 cutoff 1e-07: first step with max|coord| > 100
\newcommand{\NumEscapeStepKTwoEight}{11}
% source: e03_scaling growth k=2.8 cutoff 1e-05: log-log slope of step time vs unstable points (>= 1e3 points)
\newcommand{\NumStepTimeExponentKTwoEightCoarse}{1.08}
% source: e03_scaling growth k=2.8 cutoff 1e-05: slope of (RSS - baseline) vs points over the top two decades
\newcommand{\NumBytesPerPointKTwoEightCoarse}{1\,054}
% source: e03_scaling growth k=2.8 cutoff 1e-05: first step with max|coord| > 100
\newcommand{\NumEscapeStepKTwoEightCoarse}{11}
% source: e03_scaling growth k=4 cutoff 1e-07: log-log slope of step time vs unstable points (>= 1e3 points)
\newcommand{\NumStepTimeExponentKFour}{1.06}
% source: e03_scaling growth k=4 cutoff 1e-07: slope of (RSS - baseline) vs points over the top two decades
\newcommand{\NumBytesPerPointKFour}{1\,053}
% source: e03_scaling growth k=4 cutoff 1e-07: first step with max|coord| > 100
\newcommand{\NumEscapeStepKFour}{10}
% source: e03_scaling growth k=6 cutoff 1e-07: log-log slope of step time vs unstable points (>= 1e3 points)
\newcommand{\NumStepTimeExponentKSix}{1.05}
% source: e03_scaling growth k=6 cutoff 1e-07: slope of (RSS - baseline) vs points over the top two decades
\newcommand{\NumBytesPerPointKSix}{1\,075}
% source: e03_scaling growth k=6 cutoff 1e-07: first step with max|coord| > 100
\newcommand{\NumEscapeStepKSix}{10}
% source: e03_scaling growth k=10 cutoff 1e-07: log-log slope of step time vs unstable points (>= 1e3 points)
\newcommand{\NumStepTimeExponentKTen}{1.06}
% source: e03_scaling growth k=10 cutoff 1e-07: slope of (RSS - baseline) vs points over the top two decades
\newcommand{\NumBytesPerPointKTen}{1\,058}
% source: e03_scaling growth k=10 cutoff 1e-07: first step with max|coord| > 100
\newcommand{\NumEscapeStepKTen}{9}
% source: e03_scaling growth: median over (k, cutoff) of the step-time slope
\newcommand{\NumStepTimeExponent}{1.06}
% source: e03_scaling growth: median over (k, cutoff) of the RSS slope (bytes/point)
\newcommand{\NumBytesPerPoint}{1\,058}
% source: e03_scaling growth: max over (k, cutoff) of the RSS slope (bytes/point)
\newcommand{\NumBytesPerPointMax}{1\,075}
% source: e03_scaling stages: log-log slope of compute_intersections time vs segments, pooled over k (>= 1e3 segments)
\newcommand{\NumIntersectionsExponent}{1.07}
% source: e03_scaling stages: log-log slope of trim_and_bridges time vs segments, pooled over k (>= 1e3 segments)
\newcommand{\NumTrimBridgesExponent}{0.958}
% source: e03_scaling stages: log-log slope of infer_iterate_table time vs segments, pooled over k (>= 1e3 segments)
\newcommand{\NumIterateTableExponent}{0.0391}
% source: e04_area_cutoff k28 10 steps, cutoff 1e-07: max over branches of the p50 invariance error
\newcommand{\NumInvPFiftyKTwoEightCutoffSeven}{\ensuremath{1.8\times10^{-6}}}
% source: e04_area_cutoff k28 10 steps, cutoff 1e-07: max over branches of the p99 invariance error
\newcommand{\NumInvPNinetyNineKTwoEightCutoffSeven}{\ensuremath{6.9\times10^{-6}}}
% source: e04_area_cutoff k28 10 steps, cutoff 1e-07: summed pipeline stages (no measurement stages), s
\newcommand{\NumTimeKTwoEightCutoffSeven}{0.107}
% source: e04_area_cutoff k28 10 steps, cutoff 1e-07: unstable + stable points
\newcommand{\NumPointsKTwoEightCutoffSeven}{4\,757}
% source: e04_area_cutoff k28 10 steps, cutoff 1e-07: median NN distance of the finest run's crossings to this run's
\newcommand{\NumCrossingDevMedianKTwoEightCutoffSeven}{\ensuremath{2.3\times10^{-6}}}
% source: e04_area_cutoff k28 10 steps, cutoff 0.0001: max over branches of the p50 invariance error
\newcommand{\NumInvPFiftyKTwoEightCutoffFour}{\ensuremath{1.3\times10^{-4}}}
% source: e04_area_cutoff k28 10 steps, cutoff 0.0001: max over branches of the p99 invariance error
\newcommand{\NumInvPNinetyNineKTwoEightCutoffFour}{\ensuremath{9.6\times10^{-4}}}
% source: e04_area_cutoff k28 10 steps, cutoff 0.0001: summed pipeline stages (no measurement stages), s
\newcommand{\NumTimeKTwoEightCutoffFour}{0.0148}
% source: e04_area_cutoff k28 10 steps, cutoff 0.0001: unstable + stable points
\newcommand{\NumPointsKTwoEightCutoffFour}{450}
% source: e04_area_cutoff k28 10 steps, cutoff 0.0001: median NN distance of the finest run's crossings to this run's
\newcommand{\NumCrossingDevMedianKTwoEightCutoffFour}{\ensuremath{2.2\times10^{-4}}}
% source: e04_area_cutoff k28 canonical depth: log-log slope of points vs cutoff
\newcommand{\NumPointsCutoffExponentKTwoEight}{-0.332}
% source: e04_area_cutoff p3 13 steps, cutoff 1e-07: max over branches of the p50 invariance error
\newcommand{\NumInvPFiftyPThreeCutoffSeven}{\ensuremath{8.9\times10^{-6}}}
% source: e04_area_cutoff p3 13 steps, cutoff 1e-07: max over branches of the p99 invariance error
\newcommand{\NumInvPNinetyNinePThreeCutoffSeven}{\ensuremath{1.9\times10^{-5}}}
% source: e04_area_cutoff p3 13 steps, cutoff 1e-07: summed pipeline stages (no measurement stages), s
\newcommand{\NumTimePThreeCutoffSeven}{0.0462}
% source: e04_area_cutoff p3 13 steps, cutoff 1e-07: unstable + stable points
\newcommand{\NumPointsPThreeCutoffSeven}{1\,127}
% source: e04_area_cutoff p3 13 steps, cutoff 1e-07: median NN distance of the finest run's crossings to this run's
\newcommand{\NumCrossingDevMedianPThreeCutoffSeven}{\ensuremath{1.2\times10^{-5}}}
% source: e04_area_cutoff p3 13 steps, cutoff 0.0001: max over branches of the p50 invariance error
\newcommand{\NumInvPFiftyPThreeCutoffFour}{\ensuremath{1.1\times10^{-4}}}
% source: e04_area_cutoff p3 13 steps, cutoff 0.0001: max over branches of the p99 invariance error
\newcommand{\NumInvPNinetyNinePThreeCutoffFour}{\ensuremath{1.4\times10^{-3}}}
% source: e04_area_cutoff p3 13 steps, cutoff 0.0001: summed pipeline stages (no measurement stages), s
\newcommand{\NumTimePThreeCutoffFour}{0.0179}
% source: e04_area_cutoff p3 13 steps, cutoff 0.0001: unstable + stable points
\newcommand{\NumPointsPThreeCutoffFour}{207}
% source: e04_area_cutoff p3 13 steps, cutoff 0.0001: median NN distance of the finest run's crossings to this run's
\newcommand{\NumCrossingDevMedianPThreeCutoffFour}{\ensuremath{6.5\times10^{-4}}}
% source: e04_area_cutoff p3 canonical depth: log-log slope of points vs cutoff
\newcommand{\NumPointsCutoffExponentPThree}{-0.272}
% source: e05_invariants series iter_links (<= 2000 per case): median |f^n(x) - x_target|
\newcommand{\NumIterLinkErrMedian}{\ensuremath{6.4\times10^{-6}}}
% source: e05_invariants metrics iter_n<n>_max: max over cases and n
\newcommand{\NumIterLinkErrMax}{\ensuremath{1.5\times10^{-3}}}
% source: e05_invariants metrics iter_links: total checked
\newcommand{\NumIterLinks}{70}
% source: e05_invariants metrics lobe_ratio_max: max |A_img/A - 1|
\newcommand{\NumLobeRatioMax}{\ensuremath{5.0\times10^{-3}}}
% source: e05_invariants metrics lobe_pairs: total
\newcommand{\NumLobePairs}{57}
% source: e05_invariants metrics lobe_sign_flips: total
\newcommand{\NumLobeSignFlips}{0}
% source: e05_invariants metrics fp: max |f^p(x*) - x*|
\newcommand{\NumFpResidualMax}{\ensuremath{8.9\times10^{-16}}}
% source: e05_invariants metrics fp: max ||lu ls| - 1|
\newcommand{\NumDetDefectMax}{\ensuremath{1.2\times10^{-11}}}
% source: e05_invariants metrics fp: max |J v - lambda v| / |v|, both eigenvectors
\newcommand{\NumEvecResidualMax}{\ensuremath{1.8\times10^{-15}}}
% source: e05_invariants metrics fp: max |stored - analytic| cycle Jacobian entry
\newcommand{\NumJacobianDeviationMax}{\ensuremath{1.4\times10^{-12}}}
% source: e05_invariants metrics inv_*_max: largest invariance error over cases and branches
\newcommand{\NumInvarianceMax}{\ensuremath{1.1\times10^{-3}}}
% source: e05_invariants metrics inv_*_beyond: images past the image branch's tail (expected 0)
\newcommand{\NumInvarianceBeyond}{0}
% source: e05_invariants metrics cdist_manifolds.decreasing: total over cases
\newcommand{\NumCdistDecreasingManifolds}{0}
% source: e05_invariants metrics cdist_manifolds.ties: total over cases
\newcommand{\NumCdistTiesManifolds}{0}
% source: e05_invariants metrics cdist_manifolds.steps: total over cases
\newcommand{\NumCdistStepsManifolds}{77\,328}
% source: e05_invariants metrics cdist_bridges.decreasing: total over cases
\newcommand{\NumCdistDecreasingBridges}{0}
% source: e05_invariants metrics cdist_bridges.ties: total over cases
\newcommand{\NumCdistTiesBridges}{0}
% source: e05_invariants metrics cdist_bridges.steps: total over cases
\newcommand{\NumCdistStepsBridges}{317\,502}
% source: e05_invariants metrics arr_image_found: total
\newcommand{\NumArrImageFound}{47}
% source: e05_invariants metrics arr_image_missing: total
\newcommand{\NumArrImageMissing}{39}
% source: e05_invariants gpu_k28: max polyline Hausdorff distance CPU vs GPU over manifolds
\newcommand{\NumGpuHausdorffMax}{\ensuremath{2.1\times10^{-14}}}
% source: e05_invariants gpu_k28: every manifold has the same point count on CPU and GPU
\newcommand{\NumGpuSameLengths}{yes}
% source: e06_parameter_sweep configs run
\newcommand{\NumSweepConfigs}{98}
% source: e06_parameter_sweep configs ok through trellis_pips
\newcommand{\NumSweepComplete}{71}
% source: e06_parameter_sweep configs ok through trellis_pips / all configs
\newcommand{\NumSweepSuccessFraction}{72.4\%}
% source: e06_parameter_sweep b = 1 saddle configs ok through trellis_pips / b = 1 saddle configs
\newcommand{\NumSaddleSuccessFraction}{83.3\%}
% source: e06_parameter_sweep smallest k whose b = 1 saddle completes
\newcommand{\NumSaddleMinKComplete}{0.582}
% source: e06_parameter_sweep configs whose last completed stage is none
\newcommand{\NumSweepStageNone}{14}
% source: e06_parameter_sweep configs whose last completed stage is fixed_point
\newcommand{\NumSweepStageFixedPoint}{0}
% source: e06_parameter_sweep configs whose last completed stage is eigen
\newcommand{\NumSweepStageEigen}{0}
% source: e06_parameter_sweep configs whose last completed stage is seeded
\newcommand{\NumSweepStageSeeded}{3}
% source: e06_parameter_sweep configs whose last completed stage is grown_unstable
\newcommand{\NumSweepStageGrownUnstable}{10}
% source: e06_parameter_sweep configs whose last completed stage is grown
\newcommand{\NumSweepStageGrown}{0}
% source: e06_parameter_sweep configs whose last completed stage is intersected
\newcommand{\NumSweepStageIntersected}{0}
% source: e06_parameter_sweep configs whose last completed stage is bridges
\newcommand{\NumSweepStageBridges}{0}
% source: e06_parameter_sweep configs whose last completed stage is trellis_pips
\newcommand{\NumSweepStageTrellisPips}{71}
% source: e06_parameter_sweep configs with outcome ok
\newcommand{\NumSweepOutcomeOk}{71}
% source: e06_parameter_sweep configs with outcome exception
\newcommand{\NumSweepOutcomeException}{24}
% source: e06_parameter_sweep configs with outcome timeout
\newcommand{\NumSweepOutcomeTimeout}{0}
% source: e06_parameter_sweep configs with outcome oom
\newcommand{\NumSweepOutcomeOom}{3}
% source: e06_parameter_sweep configs with outcome invariant_assert
\newcommand{\NumSweepOutcomeInvariantAssert}{0}
% source: e06_parameter_sweep configs with outcome crash
\newcommand{\NumSweepOutcomeCrash}{0}
% source: e07_breaking escape fixed: record outcome
\newcommand{\NumEscapeOutcome}{oom}
% source: e07_breaking escape fixed: steps completed
\newcommand{\NumEscapeStepsDone}{9}
% source: e07_breaking escape fixed: first step with max|coord| > 100
\newcommand{\NumEscapeBlowupStep}{9}
% source: e07_breaking escape fixed: max unstable points over steps
\newcommand{\NumEscapePeakPoints}{9\,881}
% source: e07_breaking escape fixed: peak RSS (MB), max of per-step samples and the watchdog's
\newcommand{\NumEscapePeakRssMb}{20\,496}
% source: e07_breaking escape fixed: largest max|coord|
\newcommand{\NumEscapeMaxCoord}{\ensuremath{2.0\times10^{3}}}
% source: e07_breaking escape arclength 0.5: final unstable points (ok)
\newcommand{\NumEscapeStopArclengthPoints}{478}
% source: e07_breaking escape arclength 3.0: final unstable points (ok)
\newcommand{\NumEscapeStopArclengthThreePoints}{9\,881}
% source: e07_breaking escape predicate 2.5: final unstable points (ok)
\newcommand{\NumEscapeStopPredicatePoints}{478}
% source: e07_breaking blasts k28: blasts completed of 10
\newcommand{\NumMaxBlastsKTwoEight}{10}
% source: e07_breaking blasts k28: record outcome
\newcommand{\NumBlastOutcomeKTwoEight}{ok}
% source: e07_breaking blasts k28: topology stack on the final state
\newcommand{\NumBlastTopologyKTwoEight}{ok (provisional)}
% source: e07_breaking blasts p3: blasts completed of 10
\newcommand{\NumMaxBlastsPThree}{10}
% source: e07_breaking blasts p3: record outcome
\newcommand{\NumBlastOutcomePThree}{ok}
% source: e07_breaking blasts p3: topology stack on the final state
\newcommand{\NumBlastTopologyPThree}{ok (provisional)}
% source: e07_breaking blasts nested: blasts completed of 8
\newcommand{\NumMaxBlastsNested}{8}
% source: e07_breaking blasts nested: record outcome
\newcommand{\NumBlastOutcomeNested}{ok}
% source: e07_breaking blasts nested: topology stack on the final state
\newcommand{\NumBlastTopologyNested}{ok}
% source: e07_breaking depth p3: deepest run through bridges
\newcommand{\NumDepthNumericsMaxPThree}{17}
% source: e07_breaking depth p3: deepest run with topology ok
\newcommand{\NumDepthTopologyMaxPThree}{17}
% source: e07_breaking depth inversion: deepest run through bridges
\newcommand{\NumDepthNumericsMaxInversion}{6}
% source: e07_breaking depth inversion: deepest run with topology ok
\newcommand{\NumDepthTopologyMaxInversion}{6}
% source: e07_breaking tangency: smallest |sin| at any crossing
\newcommand{\NumTangencyMinSin}{\ensuremath{1.3\times10^{-7}}}
% source: e07_breaking tangency: crossings with |sin| < 1e-6, total
\newcommand{\NumTangencySinBelowSix}{2}
% source: e07_breaking tangency: WARNING records mentioning 'parallel'
\newcommand{\NumTangencyNearParallelWarnings}{0}
% source: e08_fixed_point_solver solves (seeds x 2 variants x configs)
\newcommand{\NumSolverSolves}{9\,600}
% source: e08_fixed_point_solver solves where compute_fixed_point returned / all
\newcommand{\NumSolverConvergedRate}{88\%}
% source: e08_fixed_point_solver saddle of the requested minimal period / all
\newcommand{\NumSolverSaddleRate}{49.7\%}
% source: e08_fixed_point_solver converged with minimal period < requested / converged
\newcommand{\NumSolverWrongPeriodRate}{34.8\%}
% source: e08_fixed_point_solver fsolve did not converge / all
\newcommand{\NumSolverNoConvergenceRate}{12\%}
% source: e08_fixed_point_solver max |f(x_{i-1}) - x_i| of converged orbits
\newcommand{\NumSolverMaxResidual}{\ensuremath{5.8\times10^{-11}}}
% source: e08_fixed_point_solver k = 10: periods whose distinct orbits found equal the full 2-shift count
\newcommand{\NumHorseshoeMatched}{6/6}
% source: e08_fixed_point_solver k = 10: distinct orbits found, periods 1-6
\newcommand{\NumHorseshoeFound}{23}
% source: e08_fixed_point_solver full 2-shift orbits of minimal period 1-6
\newcommand{\NumHorseshoeExpected}{23}
% source: e09_intersections ok configs
\newcommand{\NumIsectConfigs}{26}
% source: e09_intersections brute-force proper crossings, total
\newcommand{\NumIsectBruteTotal}{154}
% source: e09_intersections brute crossings matched by the registry, total
\newcommand{\NumIsectMatchedTotal}{154}
% source: e09_intersections missed (all classes), total
\newcommand{\NumIsectMissedTotal}{0}
% source: e09_intersections missed, unexplained, total
\newcommand{\NumIsectMissedUnexplained}{0}
% source: e09_intersections extra registry crossings (no anchors), total
\newcommand{\NumIsectExtraTotal}{0}
% source: e09_intersections extra, unexplained, total
\newcommand{\NumIsectExtraUnexplained}{0}
% source: e09_intersections u x u proper crossings (non-adjacent segments), total
\newcommand{\NumIsectUUTotal}{4}
% source: e09_intersections configs above UU_MAX_SEGMENTS (u x u not checked)
\newcommand{\NumIsectUUSkipped}{2}
% source: e09_intersections log-log slope of library_s vs total segments
\newcommand{\NumIsectLibraryExponent}{1.04}
% source: e09_intersections log-log slope of brute_s vs total segments
\newcommand{\NumIsectBruteExponent}{1.32}
% source: e09_intersections max over configs of brute time / library time
\newcommand{\NumIsectBruteOverLibraryMax}{\ensuremath{1.932\times}}
% source: e10_determinism cpu: max over stages of the stage-time CV over repeats
\newcommand{\NumTimingCvCpuMax}{5.41\%}
% source: e10_determinism cpu: CV of the summed stage time over repeats
\newcommand{\NumTimingCvCpuTotal}{1.82\%}
% source: e10_determinism cpu: ok repeats
\newcommand{\NumRepeatsCpu}{5}
% source: e10_determinism gpu: max over stages of the stage-time CV over repeats
\newcommand{\NumTimingCvGpuMax}{4.56\%}
% source: e10_determinism gpu: CV of the summed stage time over repeats
\newcommand{\NumTimingCvGpuTotal}{1.21\%}
% source: e10_determinism gpu: ok repeats
\newcommand{\NumRepeatsGpu}{5}
% source: e10_determinism exact sha256 of all five arrays agree CPU-CPU (n/a with fewer than two runs)
\newcommand{\NumBitIdentCpuCpu}{identical}
% source: e10_determinism sha256 of the arrays rounded to 1e-12 agree CPU-CPU
\newcommand{\NumBitIdentCpuCpuRounded}{identical}
% source: e10_determinism exact sha256 of all five arrays agree GPU-GPU (n/a with fewer than two runs)
\newcommand{\NumBitIdentGpuGpu}{identical}
% source: e10_determinism sha256 of the arrays rounded to 1e-12 agree GPU-GPU
\newcommand{\NumBitIdentGpuGpuRounded}{identical}
% source: e10_determinism exact sha256 of all five arrays agree CPU-GPU (n/a with fewer than two runs)
\newcommand{\NumBitIdentCpuGpu}{differs in 5 of 5 hashes}
% source: e10_determinism sha256 of the arrays rounded to 1e-12 agree CPU-GPU
\newcommand{\NumBitIdentCpuGpuRounded}{differs in 3 of 5 hashes}
% source: e10_determinism max |coordinate| difference of the sorted crossing lists, every run vs the first CPU run
\newcommand{\NumCrossingDevCpuGpu}{\ensuremath{2.2\times10^{-16}}}
% source: e11_blast_depth configs
\newcommand{\NumDeepConfigs}{53}
% source: e11_blast_depth configs with outcome ok
\newcommand{\NumDeepConfigsOk}{53}
% source: e11_blast_depth max_blasts per config
\newcommand{\NumDeepMaxBlasts}{20}
% source: e11_blast_depth blasts run, all configs
\newcommand{\NumDeepBlasts}{886}
% source: e11_blast_depth topology snapshots (blast 0 included), all configs
\newcommand{\NumDeepSnapshots}{939}
% source: e11_blast_depth topology snapshots that completed
\newcommand{\NumDeepSnapshotsOk}{802}
% source: e11_blast_depth configs whose stop_reason is max_blasts
\newcommand{\NumDeepStopMaxBlasts}{22}
% source: e11_blast_depth configs whose stop_reason is budget
\newcommand{\NumDeepStopBudget}{8}
% source: e11_blast_depth configs whose stop_reason is memory_budget
\newcommand{\NumDeepStopMemoryBudget}{23}
% source: e11_blast_depth wall clock from the first start to the last end (min)
\newcommand{\NumDeepWallMin}{24}
% source: e11_blast_depth same-stability crossings summed over every snapshot
\newcommand{\NumDeepSameStability}{0}
% source: e11_blast_depth decreasing cdist steps on new bridges, summed
\newcommand{\NumDeepCdistDecreasing}{0}
% source: e11_blast_depth metrics law_violations summed
\newcommand{\NumDeepLawViolations}{0}
% source: e11_blast_depth lobe / preimage-lobe polygon pairs compared
\newcommand{\NumDeepLobeCompared}{5\,276}
% source: e11_blast_depth median over blasts of the per-blast median |area ratio - 1|
\newcommand{\NumDeepLobeRatioMedian}{\ensuremath{1.2\times10^{-5}}}
% source: e11_blast_depth largest per-blast |area ratio - 1| (polygonal probe, folds can self-intersect)
\newcommand{\NumDeepLobeRatioMax}{2.29}
% source: e11_blast_depth configs whose largest lobe ratio equals the overall largest
\newcommand{\NumDeepLobeRatioMaxWhere}{nested inner 2 iterations / blast, p3 2 iterations / blast}
% source: e11_blast_depth largest per-blast |area ratio - 1| over the reference cells
\newcommand{\NumDeepLobeRatioMaxReference}{0.0514}
% source: e11_blast_depth lobe pairs whose polygon areas have opposite signs
\newcommand{\NumDeepLobeSignFlips}{4}
% source: e11_blast_depth snapshots whose canonical words changed
\newcommand{\NumDeepStructureSteps}{490}
% source: e11_blast_depth snapshots with a change flag but the same canonical words
\newcommand{\NumDeepRelabelSteps}{3}
% source: e11_blast_depth snapshots after a blast with no change flag
\newcommand{\NumDeepNoChangeSteps}{256}
% source: e11_blast_depth snapshots whose topology raised
\newcommand{\NumDeepErrorSnapshots}{137}
% source: e11_blast_depth configs with at least one topology error
\newcommand{\NumDeepErrorConfigs}{23}
% source: e11_blast_depth distinct topology error messages (classes)
\newcommand{\NumDeepErrorClasses}{4}
% source: e11_blast_depth snapshots that raised an AssertionError
\newcommand{\NumDeepErrorAsserts}{136}
% source: e11_blast_depth snapshots that raised another exception
\newcommand{\NumDeepErrorExceptions}{1}
% source: e11_blast_depth error 'zero-width iterated element': snapshots
\newcommand{\NumDeepErrZeroWidthSnapshots}{94}
% source: e11_blast_depth error 'zero-width iterated element': configs
\newcommand{\NumDeepErrZeroWidthConfigs}{17}
% source: e11_blast_depth error 'zero-width iterated element': blast range over every config
\newcommand{\NumDeepErrZeroWidthBlasts}{5--14}
% source: e11_blast_depth error 'several outer faces': snapshots
\newcommand{\NumDeepErrOuterFacesSnapshots}{1}
% source: e11_blast_depth error 'several outer faces': configs
\newcommand{\NumDeepErrOuterFacesConfigs}{1}
% source: e11_blast_depth error 'several outer faces': blast range over every config
\newcommand{\NumDeepErrOuterFacesBlasts}{16}
% source: e11_blast_depth error 'holes disagree on bridge side': snapshots
\newcommand{\NumDeepErrHoleSideSnapshots}{3}
% source: e11_blast_depth error 'holes disagree on bridge side': configs
\newcommand{\NumDeepErrHoleSideConfigs}{2}
% source: e11_blast_depth error 'holes disagree on bridge side': blast range over every config
\newcommand{\NumDeepErrHoleSideBlasts}{16--17}
% source: e11_blast_depth error 'bridge rows disagree': snapshots
\newcommand{\NumDeepErrBridgeRowsSnapshots}{39}
% source: e11_blast_depth error 'bridge rows disagree': configs
\newcommand{\NumDeepErrBridgeRowsConfigs}{4}
% source: e11_blast_depth error 'bridge rows disagree': blast range over every config
\newcommand{\NumDeepErrBridgeRowsBlasts}{6--20}
% source: e11_blast_depth k28 reference: blasts done
\newcommand{\NumDeepBlastsKTwoEight}{15}
% source: e11_blast_depth k28 reference: stop reason
\newcommand{\NumDeepStopKTwoEight}{memory budget}
% source: e11_blast_depth k28 reference: crossings before blasting
\newcommand{\NumDeepCrossingsStartKTwoEight}{4}
% source: e11_blast_depth k28 reference: crossings after the last blast
\newcommand{\NumDeepCrossingsEndKTwoEight}{7\,792}
% source: e11_blast_depth k28 reference: bridge classes before blasting
\newcommand{\NumDeepClassesStartKTwoEight}{3}
% source: e11_blast_depth k28 reference: classes at the last completed snapshot
\newcommand{\NumDeepClassesEndKTwoEight}{57}
% source: e11_blast_depth k28 reference: blast of the last completed snapshot
\newcommand{\NumDeepClassesEndBlastKTwoEight}{15}
% source: e11_blast_depth k28 reference: active classes there
\newcommand{\NumDeepActiveEndKTwoEight}{52}
% source: e11_blast_depth k28 reference: inert classes there
\newcommand{\NumDeepInertEndKTwoEight}{5}
% source: e11_blast_depth k28 reference: first blast whose canonical words changed
\newcommand{\NumDeepFirstStructureKTwoEight}{1}
% source: e11_blast_depth k28 reference: blasts whose canonical words changed
\newcommand{\NumDeepStructureChangesKTwoEight}{4}
% source: e11_blast_depth k28 reference: blast of the largest class-count jump
\newcommand{\NumDeepJumpBlastKTwoEight}{8}
% source: e11_blast_depth k28 reference: classes before that jump (blast 7)
\newcommand{\NumDeepJumpBeforeKTwoEight}{4}
% source: e11_blast_depth k28 reference: classes after it
\newcommand{\NumDeepJumpAfterKTwoEight}{20}
% source: e11_blast_depth k28 reference: geometric-mean crossings factor per blast over the last 5 blasts
\newcommand{\NumDeepCrossingGrowthKTwoEight}{1.61}
% source: e11_blast_depth k28 reference: geometric-mean bridge_points factor per blast over the last 5 blasts
\newcommand{\NumDeepPointGrowthKTwoEight}{1.65}
% source: e11_blast_depth k28 reference: geometric-mean blast_s factor per blast over the last 5 blasts
\newcommand{\NumDeepTimeGrowthKTwoEight}{1.69}
% source: e11_blast_depth k28 reference: time of the last blast (s)
\newcommand{\NumDeepLastBlastSKTwoEight}{74}
% source: e11_blast_depth k28 reference: peak RSS (GB)
\newcommand{\NumDeepPeakRssGbKTwoEight}{9.0}
% source: e11_blast_depth k28 reference: n_unresolved at the last completed snapshot
\newcommand{\NumDeepUnresolvedEndKTwoEight}{12}
% source: e11_blast_depth k28 reference: n_virtual at the last completed snapshot
\newcommand{\NumDeepVirtualEndKTwoEight}{1}
% source: e11_blast_depth k28 reference: n_ambiguous at the last completed snapshot
\newcommand{\NumDeepAmbiguousEndKTwoEight}{3}
% source: e11_blast_depth k28 reference: is_reliable at the last completed snapshot
\newcommand{\NumDeepReliableEndKTwoEight}{no}
% source: e11_blast_depth k28 reference: provisional before blasting
\newcommand{\NumDeepProvisionalStartKTwoEight}{yes}
% source: e11_blast_depth k28 reference: first blast (> 0) whose symbolic result is provisional
\newcommand{\NumDeepFirstProvisionalKTwoEight}{8}
% source: e11_blast_depth k28 reference: snapshots whose topology raised
\newcommand{\NumDeepErrorsKTwoEight}{6}
% source: e11_blast_depth k28 reference: blasts whose topology raised
\newcommand{\NumDeepErrorBlastsKTwoEight}{9--14}
% source: e11_blast_depth k28: configs with a topology error
\newcommand{\NumDeepErrorCellsKTwoEight}{11}
% source: e11_blast_depth k28: configs
\newcommand{\NumDeepCellsKTwoEight}{12}
% source: e11_blast_depth k28: separation cells identical to the reference at every blast (blasts, crossings, bridge points, topology outcome and exact topology hash)
\newcommand{\NumDeepSeparationIdenticalKTwoEight}{5/5}
% source: e11_blast_depth k28: repeat cells identical to the reference at every blast (blasts, crossings, bridge points, topology outcome and exact topology hash)
\newcommand{\NumDeepRepeatIdenticalKTwoEight}{2/2}
% source: e11_blast_depth k28 cutoff 1e-08: first blast whose canonical_hash differs from the reference, compared over blasts where both snapshots completed; 'never' if none
\newcommand{\NumDeepCutoffEightWordsKTwoEight}{never}
% source: e11_blast_depth k28 cutoff 1e-08: first blast whose crossings differs from the reference, compared over every blast both ran; 'never' if none
\newcommand{\NumDeepCutoffEightCrossingsKTwoEight}{9}
% source: e11_blast_depth k28 cutoff 1e-08: first blast whose topology outcome (completed / raised) differs from the reference; 'never' if none
\newcommand{\NumDeepCutoffEightOutcomeKTwoEight}{9}
% source: e11_blast_depth k28 cutoff 1e-08: classes at its last completed snapshot
\newcommand{\NumDeepCutoffEightClassesKTwoEight}{39}
% source: e11_blast_depth k28 cutoff 1e-08: blasts done
\newcommand{\NumDeepCutoffEightBlastsKTwoEight}{13}
% source: e11_blast_depth k28 cutoff 1e-06: first blast whose canonical_hash differs from the reference, compared over blasts where both snapshots completed; 'never' if none
\newcommand{\NumDeepCutoffSixWordsKTwoEight}{15}
% source: e11_blast_depth k28 cutoff 1e-06: first blast whose crossings differs from the reference, compared over every blast both ran; 'never' if none
\newcommand{\NumDeepCutoffSixCrossingsKTwoEight}{10}
% source: e11_blast_depth k28 cutoff 1e-06: first blast whose topology outcome (completed / raised) differs from the reference; 'never' if none
\newcommand{\NumDeepCutoffSixOutcomeKTwoEight}{9}
% source: e11_blast_depth k28 cutoff 1e-06: classes at its last completed snapshot
\newcommand{\NumDeepCutoffSixClassesKTwoEight}{57}
% source: e11_blast_depth k28 cutoff 1e-06: blasts done
\newcommand{\NumDeepCutoffSixBlastsKTwoEight}{17}
% source: e11_blast_depth k28 cutoff 1e-05: first blast whose canonical_hash differs from the reference, compared over blasts where both snapshots completed; 'never' if none
\newcommand{\NumDeepCutoffFiveWordsKTwoEight}{7}
% source: e11_blast_depth k28 cutoff 1e-05: first blast whose crossings differs from the reference, compared over every blast both ran; 'never' if none
\newcommand{\NumDeepCutoffFiveCrossingsKTwoEight}{11}
% source: e11_blast_depth k28 cutoff 1e-05: first blast whose topology outcome (completed / raised) differs from the reference; 'never' if none
\newcommand{\NumDeepCutoffFiveOutcomeKTwoEight}{9}
% source: e11_blast_depth k28 cutoff 1e-05: classes at its last completed snapshot
\newcommand{\NumDeepCutoffFiveClassesKTwoEight}{69}
% source: e11_blast_depth k28 cutoff 1e-05: blasts done
\newcommand{\NumDeepCutoffFiveBlastsKTwoEight}{17}
% source: e11_blast_depth k28 2 iterations: last blast n with crossings equal to single blast 2n for every n' <= n
\newcommand{\NumDeepDoubleCrossingsAgreeKTwoEight}{4}
% source: e11_blast_depth k28 2 iterations: blasts n >= 1 whose canonical words equal single blast 2n / pairs n >= 1 both completed (n = 0 is the shared pre-blast build and is left out)
\newcommand{\NumDeepDoubleWordsSameKTwoEight}{4/4}
% source: e11_blast_depth k28 2 iterations: blasts n >= 1 whose canonical words differ from single blast 2n
\newcommand{\NumDeepDoubleWordsDifferKTwoEight}{none}
% source: e11_blast_depth k28 2 iterations: pairs (n, 2n) compared
\newcommand{\NumDeepDoubleComparedKTwoEight}{8}
% source: e11_blast_depth k28 2 iterations: blasts done
\newcommand{\NumDeepDoubleBlastsKTwoEight}{9}
% source: e11_blast_depth k32 reference: blasts done
\newcommand{\NumDeepBlastsKThreeTwo}{14}
% source: e11_blast_depth k32 reference: stop reason
\newcommand{\NumDeepStopKThreeTwo}{memory budget}
% source: e11_blast_depth k32 reference: crossings before blasting
\newcommand{\NumDeepCrossingsStartKThreeTwo}{4}
% source: e11_blast_depth k32 reference: crossings after the last blast
\newcommand{\NumDeepCrossingsEndKThreeTwo}{6\,757}
% source: e11_blast_depth k32 reference: bridge classes before blasting
\newcommand{\NumDeepClassesStartKThreeTwo}{3}
% source: e11_blast_depth k32 reference: classes at the last completed snapshot
\newcommand{\NumDeepClassesEndKThreeTwo}{37}
% source: e11_blast_depth k32 reference: blast of the last completed snapshot
\newcommand{\NumDeepClassesEndBlastKThreeTwo}{8}
% source: e11_blast_depth k32 reference: active classes there
\newcommand{\NumDeepActiveEndKThreeTwo}{32}
% source: e11_blast_depth k32 reference: inert classes there
\newcommand{\NumDeepInertEndKThreeTwo}{5}
% source: e11_blast_depth k32 reference: first blast whose canonical words changed
\newcommand{\NumDeepFirstStructureKThreeTwo}{1}
% source: e11_blast_depth k32 reference: blasts whose canonical words changed
\newcommand{\NumDeepStructureChangesKThreeTwo}{8}
% source: e11_blast_depth k32 reference: blast of the largest class-count jump
\newcommand{\NumDeepJumpBlastKThreeTwo}{8}
% source: e11_blast_depth k32 reference: classes before that jump (blast 7)
\newcommand{\NumDeepJumpBeforeKThreeTwo}{15}
% source: e11_blast_depth k32 reference: classes after it
\newcommand{\NumDeepJumpAfterKThreeTwo}{37}
% source: e11_blast_depth k32 reference: geometric-mean crossings factor per blast over the last 5 blasts
\newcommand{\NumDeepCrossingGrowthKThreeTwo}{1.67}
% source: e11_blast_depth k32 reference: geometric-mean bridge_points factor per blast over the last 5 blasts
\newcommand{\NumDeepPointGrowthKThreeTwo}{1.70}
% source: e11_blast_depth k32 reference: geometric-mean blast_s factor per blast over the last 5 blasts
\newcommand{\NumDeepTimeGrowthKThreeTwo}{1.75}
% source: e11_blast_depth k32 reference: time of the last blast (s)
\newcommand{\NumDeepLastBlastSKThreeTwo}{72}
% source: e11_blast_depth k32 reference: peak RSS (GB)
\newcommand{\NumDeepPeakRssGbKThreeTwo}{7.6}
% source: e11_blast_depth k32 reference: n_unresolved at the last completed snapshot
\newcommand{\NumDeepUnresolvedEndKThreeTwo}{2}
% source: e11_blast_depth k32 reference: n_virtual at the last completed snapshot
\newcommand{\NumDeepVirtualEndKThreeTwo}{1}
% source: e11_blast_depth k32 reference: n_ambiguous at the last completed snapshot
\newcommand{\NumDeepAmbiguousEndKThreeTwo}{0}
% source: e11_blast_depth k32 reference: is_reliable at the last completed snapshot
\newcommand{\NumDeepReliableEndKThreeTwo}{no}
% source: e11_blast_depth k32 reference: provisional before blasting
\newcommand{\NumDeepProvisionalStartKThreeTwo}{yes}
% source: e11_blast_depth k32 reference: first blast (> 0) whose symbolic result is provisional
\newcommand{\NumDeepFirstProvisionalKThreeTwo}{6}
% source: e11_blast_depth k32 reference: snapshots whose topology raised
\newcommand{\NumDeepErrorsKThreeTwo}{6}
% source: e11_blast_depth k32 reference: blasts whose topology raised
\newcommand{\NumDeepErrorBlastsKThreeTwo}{9--14}
% source: e11_blast_depth k32: configs with a topology error
\newcommand{\NumDeepErrorCellsKThreeTwo}{8}
% source: e11_blast_depth k32: configs
\newcommand{\NumDeepCellsKThreeTwo}{10}
% source: e11_blast_depth k32: separation cells identical to the reference at every blast (blasts, crossings, bridge points, topology outcome and exact topology hash)
\newcommand{\NumDeepSeparationIdenticalKThreeTwo}{5/5}
% source: e11_blast_depth k32 cutoff 1e-08: first blast whose canonical_hash differs from the reference, compared over blasts where both snapshots completed; 'never' if none
\newcommand{\NumDeepCutoffEightWordsKThreeTwo}{6}
% source: e11_blast_depth k32 cutoff 1e-08: first blast whose crossings differs from the reference, compared over every blast both ran; 'never' if none
\newcommand{\NumDeepCutoffEightCrossingsKThreeTwo}{9}
% source: e11_blast_depth k32 cutoff 1e-08: first blast whose topology outcome (completed / raised) differs from the reference; 'never' if none
\newcommand{\NumDeepCutoffEightOutcomeKThreeTwo}{9}
% source: e11_blast_depth k32 cutoff 1e-08: classes at its last completed snapshot
\newcommand{\NumDeepCutoffEightClassesKThreeTwo}{27}
% source: e11_blast_depth k32 cutoff 1e-08: blasts done
\newcommand{\NumDeepCutoffEightBlastsKThreeTwo}{13}
% source: e11_blast_depth k32 cutoff 1e-06: first blast whose canonical_hash differs from the reference, compared over blasts where both snapshots completed; 'never' if none
\newcommand{\NumDeepCutoffSixWordsKThreeTwo}{6}
% source: e11_blast_depth k32 cutoff 1e-06: first blast whose crossings differs from the reference, compared over every blast both ran; 'never' if none
\newcommand{\NumDeepCutoffSixCrossingsKThreeTwo}{10}
% source: e11_blast_depth k32 cutoff 1e-06: first blast whose topology outcome (completed / raised) differs from the reference; 'never' if none
\newcommand{\NumDeepCutoffSixOutcomeKThreeTwo}{9}
% source: e11_blast_depth k32 cutoff 1e-06: classes at its last completed snapshot
\newcommand{\NumDeepCutoffSixClassesKThreeTwo}{92}
% source: e11_blast_depth k32 cutoff 1e-06: blasts done
\newcommand{\NumDeepCutoffSixBlastsKThreeTwo}{16}
% source: e11_blast_depth k32 cutoff 1e-05: first blast whose canonical_hash differs from the reference, compared over blasts where both snapshots completed; 'never' if none
\newcommand{\NumDeepCutoffFiveWordsKThreeTwo}{6}
% source: e11_blast_depth k32 cutoff 1e-05: first blast whose crossings differs from the reference, compared over every blast both ran; 'never' if none
\newcommand{\NumDeepCutoffFiveCrossingsKThreeTwo}{10}
% source: e11_blast_depth k32 cutoff 1e-05: first blast whose topology outcome (completed / raised) differs from the reference; 'never' if none
\newcommand{\NumDeepCutoffFiveOutcomeKThreeTwo}{9}
% source: e11_blast_depth k32 cutoff 1e-05: classes at its last completed snapshot
\newcommand{\NumDeepCutoffFiveClassesKThreeTwo}{89}
% source: e11_blast_depth k32 cutoff 1e-05: blasts done
\newcommand{\NumDeepCutoffFiveBlastsKThreeTwo}{16}
% source: e11_blast_depth k32 2 iterations: last blast n with crossings equal to single blast 2n for every n' <= n
\newcommand{\NumDeepDoubleCrossingsAgreeKThreeTwo}{4}
% source: e11_blast_depth k32 2 iterations: blasts n >= 1 whose canonical words equal single blast 2n / pairs n >= 1 both completed (n = 0 is the shared pre-blast build and is left out)
\newcommand{\NumDeepDoubleWordsSameKThreeTwo}{4/4}
% source: e11_blast_depth k32 2 iterations: blasts n >= 1 whose canonical words differ from single blast 2n
\newcommand{\NumDeepDoubleWordsDifferKThreeTwo}{none}
% source: e11_blast_depth k32 2 iterations: pairs (n, 2n) compared
\newcommand{\NumDeepDoubleComparedKThreeTwo}{8}
% source: e11_blast_depth k32 2 iterations: blasts done
\newcommand{\NumDeepDoubleBlastsKThreeTwo}{9}
% source: e11_blast_depth p3 reference: blasts done
\newcommand{\NumDeepBlastsPThree}{20}
% source: e11_blast_depth p3 reference: stop reason
\newcommand{\NumDeepStopPThree}{all blasts}
% source: e11_blast_depth p3 reference: crossings before blasting
\newcommand{\NumDeepCrossingsStartPThree}{30}
% source: e11_blast_depth p3 reference: crossings after the last blast
\newcommand{\NumDeepCrossingsEndPThree}{198}
% source: e11_blast_depth p3 reference: bridge classes before blasting
\newcommand{\NumDeepClassesStartPThree}{9}
% source: e11_blast_depth p3 reference: classes at the last completed snapshot
\newcommand{\NumDeepClassesEndPThree}{51}
% source: e11_blast_depth p3 reference: blast of the last completed snapshot
\newcommand{\NumDeepClassesEndBlastPThree}{20}
% source: e11_blast_depth p3 reference: active classes there
\newcommand{\NumDeepActiveEndPThree}{38}
% source: e11_blast_depth p3 reference: inert classes there
\newcommand{\NumDeepInertEndPThree}{13}
% source: e11_blast_depth p3 reference: first blast whose canonical words changed
\newcommand{\NumDeepFirstStructurePThree}{6}
% source: e11_blast_depth p3 reference: blasts whose canonical words changed
\newcommand{\NumDeepStructureChangesPThree}{15}
% source: e11_blast_depth p3 reference: blast of the largest class-count jump
\newcommand{\NumDeepJumpBlastPThree}{6}
% source: e11_blast_depth p3 reference: classes before that jump (blast 5)
\newcommand{\NumDeepJumpBeforePThree}{9}
% source: e11_blast_depth p3 reference: classes after it
\newcommand{\NumDeepJumpAfterPThree}{23}
% source: e11_blast_depth p3 reference: classes added per blast over the last 5 blasts (steady)
\newcommand{\NumDeepClassStepPThree}{2}
% source: e11_blast_depth p3 reference: geometric-mean crossings factor per blast over the last 5 blasts
\newcommand{\NumDeepCrossingGrowthPThree}{1.08}
% source: e11_blast_depth p3 reference: geometric-mean bridge_points factor per blast over the last 5 blasts
\newcommand{\NumDeepPointGrowthPThree}{1.06}
% source: e11_blast_depth p3 reference: geometric-mean blast_s factor per blast over the last 5 blasts
\newcommand{\NumDeepTimeGrowthPThree}{0.80}
% source: e11_blast_depth p3 reference: time of the last blast (s)
\newcommand{\NumDeepLastBlastSPThree}{0.011}
% source: e11_blast_depth p3 reference: peak RSS (GB)
\newcommand{\NumDeepPeakRssGbPThree}{0.4}
% source: e11_blast_depth p3 reference: n_unresolved at the last completed snapshot
\newcommand{\NumDeepUnresolvedEndPThree}{1}
% source: e11_blast_depth p3 reference: n_virtual at the last completed snapshot
\newcommand{\NumDeepVirtualEndPThree}{1}
% source: e11_blast_depth p3 reference: n_ambiguous at the last completed snapshot
\newcommand{\NumDeepAmbiguousEndPThree}{0}
% source: e11_blast_depth p3 reference: is_reliable at the last completed snapshot
\newcommand{\NumDeepReliableEndPThree}{no}
% source: e11_blast_depth p3 reference: provisional before blasting
\newcommand{\NumDeepProvisionalStartPThree}{no}
% source: e11_blast_depth p3 reference: first blast (> 0) whose symbolic result is provisional
\newcommand{\NumDeepFirstProvisionalPThree}{6}
% source: e11_blast_depth p3 reference: snapshots whose topology raised
\newcommand{\NumDeepErrorsPThree}{0}
% source: e11_blast_depth p3 reference: blasts whose topology raised
\newcommand{\NumDeepErrorBlastsPThree}{none}
% source: e11_blast_depth p3: configs with a topology error
\newcommand{\NumDeepErrorCellsPThree}{2}
% source: e11_blast_depth p3: configs
\newcommand{\NumDeepCellsPThree}{11}
% source: e11_blast_depth p3: separation cells identical to the reference at every blast (blasts, crossings, bridge points, topology outcome and exact topology hash)
\newcommand{\NumDeepSeparationIdenticalPThree}{4/4}
% source: e11_blast_depth p3: repeat cells identical to the reference at every blast (blasts, crossings, bridge points, topology outcome and exact topology hash)
\newcommand{\NumDeepRepeatIdenticalPThree}{2/2}
% source: e11_blast_depth p3 cutoff 1e-08: first blast whose canonical_hash differs from the reference, compared over blasts where both snapshots completed; 'never' if none
\newcommand{\NumDeepCutoffEightWordsPThree}{6}
% source: e11_blast_depth p3 cutoff 1e-08: first blast whose crossings differs from the reference, compared over every blast both ran; 'never' if none
\newcommand{\NumDeepCutoffEightCrossingsPThree}{9}
% source: e11_blast_depth p3 cutoff 1e-08: first blast whose topology outcome (completed / raised) differs from the reference; 'never' if none
\newcommand{\NumDeepCutoffEightOutcomePThree}{never}
% source: e11_blast_depth p3 cutoff 1e-08: classes at its last completed snapshot
\newcommand{\NumDeepCutoffEightClassesPThree}{51}
% source: e11_blast_depth p3 cutoff 1e-08: blasts done
\newcommand{\NumDeepCutoffEightBlastsPThree}{20}
% source: e11_blast_depth p3 cutoff 1e-06: first blast whose canonical_hash differs from the reference, compared over blasts where both snapshots completed; 'never' if none
\newcommand{\NumDeepCutoffSixWordsPThree}{4}
% source: e11_blast_depth p3 cutoff 1e-06: first blast whose crossings differs from the reference, compared over every blast both ran; 'never' if none
\newcommand{\NumDeepCutoffSixCrossingsPThree}{7}
% source: e11_blast_depth p3 cutoff 1e-06: first blast whose topology outcome (completed / raised) differs from the reference; 'never' if none
\newcommand{\NumDeepCutoffSixOutcomePThree}{never}
% source: e11_blast_depth p3 cutoff 1e-06: classes at its last completed snapshot
\newcommand{\NumDeepCutoffSixClassesPThree}{41}
% source: e11_blast_depth p3 cutoff 1e-06: blasts done
\newcommand{\NumDeepCutoffSixBlastsPThree}{20}
% source: e11_blast_depth p3 cutoff 1e-05: first blast whose canonical_hash differs from the reference, compared over blasts where both snapshots completed; 'never' if none
\newcommand{\NumDeepCutoffFiveWordsPThree}{3}
% source: e11_blast_depth p3 cutoff 1e-05: first blast whose crossings differs from the reference, compared over every blast both ran; 'never' if none
\newcommand{\NumDeepCutoffFiveCrossingsPThree}{6}
% source: e11_blast_depth p3 cutoff 1e-05: first blast whose topology outcome (completed / raised) differs from the reference; 'never' if none
\newcommand{\NumDeepCutoffFiveOutcomePThree}{6}
% source: e11_blast_depth p3 cutoff 1e-05: classes at its last completed snapshot
\newcommand{\NumDeepCutoffFiveClassesPThree}{50}
% source: e11_blast_depth p3 cutoff 1e-05: blasts done
\newcommand{\NumDeepCutoffFiveBlastsPThree}{20}
% source: e11_blast_depth p3 2 iterations: last blast n with crossings equal to single blast 2n for every n' <= n
\newcommand{\NumDeepDoubleCrossingsAgreePThree}{10}
% source: e11_blast_depth p3 2 iterations: blasts n >= 1 whose canonical words equal single blast 2n / pairs n >= 1 both completed (n = 0 is the shared pre-blast build and is left out)
\newcommand{\NumDeepDoubleWordsSamePThree}{10/10}
% source: e11_blast_depth p3 2 iterations: blasts n >= 1 whose canonical words differ from single blast 2n
\newcommand{\NumDeepDoubleWordsDifferPThree}{none}
% source: e11_blast_depth p3 2 iterations: pairs (n, 2n) compared
\newcommand{\NumDeepDoubleComparedPThree}{11}
% source: e11_blast_depth p3 2 iterations: blasts done
\newcommand{\NumDeepDoubleBlastsPThree}{20}
% source: e11_blast_depth nested_outer reference: blasts done
\newcommand{\NumDeepBlastsNestedOuter}{15}
% source: e11_blast_depth nested_outer reference: stop reason
\newcommand{\NumDeepStopNestedOuter}{memory budget}
% source: e11_blast_depth nested_outer reference: crossings before blasting
\newcommand{\NumDeepCrossingsStartNestedOuter}{34}
% source: e11_blast_depth nested_outer reference: crossings after the last blast
\newcommand{\NumDeepCrossingsEndNestedOuter}{9\,170}
% source: e11_blast_depth nested_outer reference: bridge classes before blasting
\newcommand{\NumDeepClassesStartNestedOuter}{12}
% source: e11_blast_depth nested_outer reference: classes at the last completed snapshot
\newcommand{\NumDeepClassesEndNestedOuter}{11}
% source: e11_blast_depth nested_outer reference: blast of the last completed snapshot
\newcommand{\NumDeepClassesEndBlastNestedOuter}{15}
% source: e11_blast_depth nested_outer reference: active classes there
\newcommand{\NumDeepActiveEndNestedOuter}{3}
% source: e11_blast_depth nested_outer reference: inert classes there
\newcommand{\NumDeepInertEndNestedOuter}{8}
% source: e11_blast_depth nested_outer reference: first blast whose canonical words changed
\newcommand{\NumDeepFirstStructureNestedOuter}{1}
% source: e11_blast_depth nested_outer reference: blasts whose canonical words changed
\newcommand{\NumDeepStructureChangesNestedOuter}{3}
% source: e11_blast_depth nested_outer reference: blast of the largest class-count jump
\newcommand{\NumDeepJumpBlastNestedOuter}{2}
% source: e11_blast_depth nested_outer reference: classes before that jump (blast 1)
\newcommand{\NumDeepJumpBeforeNestedOuter}{12}
% source: e11_blast_depth nested_outer reference: classes after it
\newcommand{\NumDeepJumpAfterNestedOuter}{13}
% source: e11_blast_depth nested_outer reference: blast of the largest class-count drop
\newcommand{\NumDeepDropBlastNestedOuter}{10}
% source: e11_blast_depth nested_outer reference: classes before that drop
\newcommand{\NumDeepDropBeforeNestedOuter}{13}
% source: e11_blast_depth nested_outer reference: classes after it
\newcommand{\NumDeepDropAfterNestedOuter}{11}
% source: e11_blast_depth nested_outer reference: geometric-mean crossings factor per blast over the last 5 blasts
\newcommand{\NumDeepCrossingGrowthNestedOuter}{1.66}
% source: e11_blast_depth nested_outer reference: geometric-mean bridge_points factor per blast over the last 5 blasts
\newcommand{\NumDeepPointGrowthNestedOuter}{1.68}
% source: e11_blast_depth nested_outer reference: geometric-mean blast_s factor per blast over the last 5 blasts
\newcommand{\NumDeepTimeGrowthNestedOuter}{1.77}
% source: e11_blast_depth nested_outer reference: time of the last blast (s)
\newcommand{\NumDeepLastBlastSNestedOuter}{76}
% source: e11_blast_depth nested_outer reference: peak RSS (GB)
\newcommand{\NumDeepPeakRssGbNestedOuter}{7.5}
% source: e11_blast_depth nested_outer reference: n_unresolved at the last completed snapshot
\newcommand{\NumDeepUnresolvedEndNestedOuter}{0}
% source: e11_blast_depth nested_outer reference: n_virtual at the last completed snapshot
\newcommand{\NumDeepVirtualEndNestedOuter}{0}
% source: e11_blast_depth nested_outer reference: n_ambiguous at the last completed snapshot
\newcommand{\NumDeepAmbiguousEndNestedOuter}{0}
% source: e11_blast_depth nested_outer reference: is_reliable at the last completed snapshot
\newcommand{\NumDeepReliableEndNestedOuter}{yes}
% source: e11_blast_depth nested_outer reference: provisional before blasting
\newcommand{\NumDeepProvisionalStartNestedOuter}{yes}
% source: e11_blast_depth nested_outer reference: first blast (> 0) whose symbolic result is provisional
\newcommand{\NumDeepFirstProvisionalNestedOuter}{n/a}
% source: e11_blast_depth nested_outer reference: snapshots whose topology raised
\newcommand{\NumDeepErrorsNestedOuter}{0}
% source: e11_blast_depth nested_outer reference: blasts whose topology raised
\newcommand{\NumDeepErrorBlastsNestedOuter}{none}
% source: e11_blast_depth nested_outer: configs with a topology error
\newcommand{\NumDeepErrorCellsNestedOuter}{0}
% source: e11_blast_depth nested_outer: configs
\newcommand{\NumDeepCellsNestedOuter}{9}
% source: e11_blast_depth nested_outer: separation cells identical to the reference at every blast (blasts, crossings, bridge points, topology outcome and exact topology hash)
\newcommand{\NumDeepSeparationIdenticalNestedOuter}{4/4}
% source: e11_blast_depth nested_outer cutoff 1e-08: first blast whose canonical_hash differs from the reference, compared over blasts where both snapshots completed; 'never' if none
\newcommand{\NumDeepCutoffEightWordsNestedOuter}{never}
% source: e11_blast_depth nested_outer cutoff 1e-08: first blast whose crossings differs from the reference, compared over every blast both ran; 'never' if none
\newcommand{\NumDeepCutoffEightCrossingsNestedOuter}{12}
% source: e11_blast_depth nested_outer cutoff 1e-08: first blast whose topology outcome (completed / raised) differs from the reference; 'never' if none
\newcommand{\NumDeepCutoffEightOutcomeNestedOuter}{never}
% source: e11_blast_depth nested_outer cutoff 1e-08: classes at its last completed snapshot
\newcommand{\NumDeepCutoffEightClassesNestedOuter}{11}
% source: e11_blast_depth nested_outer cutoff 1e-08: blasts done
\newcommand{\NumDeepCutoffEightBlastsNestedOuter}{14}
% source: e11_blast_depth nested_outer cutoff 1e-06: first blast whose canonical_hash differs from the reference, compared over blasts where both snapshots completed; 'never' if none
\newcommand{\NumDeepCutoffSixWordsNestedOuter}{9}
% source: e11_blast_depth nested_outer cutoff 1e-06: first blast whose crossings differs from the reference, compared over every blast both ran; 'never' if none
\newcommand{\NumDeepCutoffSixCrossingsNestedOuter}{10}
% source: e11_blast_depth nested_outer cutoff 1e-06: first blast whose topology outcome (completed / raised) differs from the reference; 'never' if none
\newcommand{\NumDeepCutoffSixOutcomeNestedOuter}{never}
% source: e11_blast_depth nested_outer cutoff 1e-06: classes at its last completed snapshot
\newcommand{\NumDeepCutoffSixClassesNestedOuter}{12}
% source: e11_blast_depth nested_outer cutoff 1e-06: blasts done
\newcommand{\NumDeepCutoffSixBlastsNestedOuter}{16}
% source: e11_blast_depth nested_outer cutoff 1e-05: first blast whose canonical_hash differs from the reference, compared over blasts where both snapshots completed; 'never' if none
\newcommand{\NumDeepCutoffFiveWordsNestedOuter}{8}
% source: e11_blast_depth nested_outer cutoff 1e-05: first blast whose crossings differs from the reference, compared over every blast both ran; 'never' if none
\newcommand{\NumDeepCutoffFiveCrossingsNestedOuter}{10}
% source: e11_blast_depth nested_outer cutoff 1e-05: first blast whose topology outcome (completed / raised) differs from the reference; 'never' if none
\newcommand{\NumDeepCutoffFiveOutcomeNestedOuter}{never}
% source: e11_blast_depth nested_outer cutoff 1e-05: classes at its last completed snapshot
\newcommand{\NumDeepCutoffFiveClassesNestedOuter}{12}
% source: e11_blast_depth nested_outer cutoff 1e-05: blasts done
\newcommand{\NumDeepCutoffFiveBlastsNestedOuter}{16}
% source: e11_blast_depth nested_outer 2 iterations: last blast n with crossings equal to single blast 2n for every n' <= n
\newcommand{\NumDeepDoubleCrossingsAgreeNestedOuter}{3}
% source: e11_blast_depth nested_outer 2 iterations: blasts n >= 1 whose canonical words equal single blast 2n / pairs n >= 1 both completed (n = 0 is the shared pre-blast build and is left out)
\newcommand{\NumDeepDoubleWordsSameNestedOuter}{6/7}
% source: e11_blast_depth nested_outer 2 iterations: blasts n >= 1 whose canonical words differ from single blast 2n
\newcommand{\NumDeepDoubleWordsDifferNestedOuter}{5}
% source: e11_blast_depth nested_outer 2 iterations: pairs (n, 2n) compared
\newcommand{\NumDeepDoubleComparedNestedOuter}{8}
% source: e11_blast_depth nested_outer 2 iterations: blasts done
\newcommand{\NumDeepDoubleBlastsNestedOuter}{11}
% source: e11_blast_depth nested_inner reference: blasts done
\newcommand{\NumDeepBlastsNestedInner}{20}
% source: e11_blast_depth nested_inner reference: stop reason
\newcommand{\NumDeepStopNestedInner}{all blasts}
% source: e11_blast_depth nested_inner reference: crossings before blasting
\newcommand{\NumDeepCrossingsStartNestedInner}{40}
% source: e11_blast_depth nested_inner reference: crossings after the last blast
\newcommand{\NumDeepCrossingsEndNestedInner}{208}
% source: e11_blast_depth nested_inner reference: bridge classes before blasting
\newcommand{\NumDeepClassesStartNestedInner}{13}
% source: e11_blast_depth nested_inner reference: classes at the last completed snapshot
\newcommand{\NumDeepClassesEndNestedInner}{55}
% source: e11_blast_depth nested_inner reference: blast of the last completed snapshot
\newcommand{\NumDeepClassesEndBlastNestedInner}{20}
% source: e11_blast_depth nested_inner reference: active classes there
\newcommand{\NumDeepActiveEndNestedInner}{41}
% source: e11_blast_depth nested_inner reference: inert classes there
\newcommand{\NumDeepInertEndNestedInner}{14}
% source: e11_blast_depth nested_inner reference: first blast whose canonical words changed
\newcommand{\NumDeepFirstStructureNestedInner}{6}
% source: e11_blast_depth nested_inner reference: blasts whose canonical words changed
\newcommand{\NumDeepStructureChangesNestedInner}{15}
% source: e11_blast_depth nested_inner reference: blast of the largest class-count jump
\newcommand{\NumDeepJumpBlastNestedInner}{6}
% source: e11_blast_depth nested_inner reference: classes before that jump (blast 5)
\newcommand{\NumDeepJumpBeforeNestedInner}{13}
% source: e11_blast_depth nested_inner reference: classes after it
\newcommand{\NumDeepJumpAfterNestedInner}{27}
% source: e11_blast_depth nested_inner reference: classes added per blast over the last 5 blasts (steady)
\newcommand{\NumDeepClassStepNestedInner}{2}
% source: e11_blast_depth nested_inner reference: geometric-mean crossings factor per blast over the last 5 blasts
\newcommand{\NumDeepCrossingGrowthNestedInner}{1.08}
% source: e11_blast_depth nested_inner reference: geometric-mean bridge_points factor per blast over the last 5 blasts
\newcommand{\NumDeepPointGrowthNestedInner}{1.02}
% source: e11_blast_depth nested_inner reference: geometric-mean blast_s factor per blast over the last 5 blasts
\newcommand{\NumDeepTimeGrowthNestedInner}{1.27}
% source: e11_blast_depth nested_inner reference: time of the last blast (s)
\newcommand{\NumDeepLastBlastSNestedInner}{0.06}
% source: e11_blast_depth nested_inner reference: peak RSS (GB)
\newcommand{\NumDeepPeakRssGbNestedInner}{0.4}
% source: e11_blast_depth nested_inner reference: n_unresolved at the last completed snapshot
\newcommand{\NumDeepUnresolvedEndNestedInner}{1}
% source: e11_blast_depth nested_inner reference: n_virtual at the last completed snapshot
\newcommand{\NumDeepVirtualEndNestedInner}{1}
% source: e11_blast_depth nested_inner reference: n_ambiguous at the last completed snapshot
\newcommand{\NumDeepAmbiguousEndNestedInner}{0}
% source: e11_blast_depth nested_inner reference: is_reliable at the last completed snapshot
\newcommand{\NumDeepReliableEndNestedInner}{no}
% source: e11_blast_depth nested_inner reference: provisional before blasting
\newcommand{\NumDeepProvisionalStartNestedInner}{no}
% source: e11_blast_depth nested_inner reference: first blast (> 0) whose symbolic result is provisional
\newcommand{\NumDeepFirstProvisionalNestedInner}{6}
% source: e11_blast_depth nested_inner reference: snapshots whose topology raised
\newcommand{\NumDeepErrorsNestedInner}{0}
% source: e11_blast_depth nested_inner reference: blasts whose topology raised
\newcommand{\NumDeepErrorBlastsNestedInner}{none}
% source: e11_blast_depth nested_inner: configs with a topology error
\newcommand{\NumDeepErrorCellsNestedInner}{2}
% source: e11_blast_depth nested_inner: configs
\newcommand{\NumDeepCellsNestedInner}{11}
% source: e11_blast_depth nested_inner: separation cells identical to the reference at every blast (blasts, crossings, bridge points, topology outcome and exact topology hash)
\newcommand{\NumDeepSeparationIdenticalNestedInner}{4/4}
% source: e11_blast_depth nested_inner: repeat cells identical to the reference at every blast (blasts, crossings, bridge points, topology outcome and exact topology hash)
\newcommand{\NumDeepRepeatIdenticalNestedInner}{2/2}
% source: e11_blast_depth nested_inner cutoff 1e-08: first blast whose canonical_hash differs from the reference, compared over blasts where both snapshots completed; 'never' if none
\newcommand{\NumDeepCutoffEightWordsNestedInner}{6}
% source: e11_blast_depth nested_inner cutoff 1e-08: first blast whose crossings differs from the reference, compared over every blast both ran; 'never' if none
\newcommand{\NumDeepCutoffEightCrossingsNestedInner}{9}
% source: e11_blast_depth nested_inner cutoff 1e-08: first blast whose topology outcome (completed / raised) differs from the reference; 'never' if none
\newcommand{\NumDeepCutoffEightOutcomeNestedInner}{never}
% source: e11_blast_depth nested_inner cutoff 1e-08: classes at its last completed snapshot
\newcommand{\NumDeepCutoffEightClassesNestedInner}{55}
% source: e11_blast_depth nested_inner cutoff 1e-08: blasts done
\newcommand{\NumDeepCutoffEightBlastsNestedInner}{20}
% source: e11_blast_depth nested_inner cutoff 1e-06: first blast whose canonical_hash differs from the reference, compared over blasts where both snapshots completed; 'never' if none
\newcommand{\NumDeepCutoffSixWordsNestedInner}{4}
% source: e11_blast_depth nested_inner cutoff 1e-06: first blast whose crossings differs from the reference, compared over every blast both ran; 'never' if none
\newcommand{\NumDeepCutoffSixCrossingsNestedInner}{7}
% source: e11_blast_depth nested_inner cutoff 1e-06: first blast whose topology outcome (completed / raised) differs from the reference; 'never' if none
\newcommand{\NumDeepCutoffSixOutcomeNestedInner}{never}
% source: e11_blast_depth nested_inner cutoff 1e-06: classes at its last completed snapshot
\newcommand{\NumDeepCutoffSixClassesNestedInner}{45}
% source: e11_blast_depth nested_inner cutoff 1e-06: blasts done
\newcommand{\NumDeepCutoffSixBlastsNestedInner}{20}
% source: e11_blast_depth nested_inner cutoff 1e-05: first blast whose canonical_hash differs from the reference, compared over blasts where both snapshots completed; 'never' if none
\newcommand{\NumDeepCutoffFiveWordsNestedInner}{3}
% source: e11_blast_depth nested_inner cutoff 1e-05: first blast whose crossings differs from the reference, compared over every blast both ran; 'never' if none
\newcommand{\NumDeepCutoffFiveCrossingsNestedInner}{6}
% source: e11_blast_depth nested_inner cutoff 1e-05: first blast whose topology outcome (completed / raised) differs from the reference; 'never' if none
\newcommand{\NumDeepCutoffFiveOutcomeNestedInner}{6}
% source: e11_blast_depth nested_inner cutoff 1e-05: classes at its last completed snapshot
\newcommand{\NumDeepCutoffFiveClassesNestedInner}{54}
% source: e11_blast_depth nested_inner cutoff 1e-05: blasts done
\newcommand{\NumDeepCutoffFiveBlastsNestedInner}{20}
% source: e11_blast_depth nested_inner 2 iterations: last blast n with crossings equal to single blast 2n for every n' <= n
\newcommand{\NumDeepDoubleCrossingsAgreeNestedInner}{5}
% source: e11_blast_depth nested_inner 2 iterations: blasts n >= 1 whose canonical words equal single blast 2n / pairs n >= 1 both completed (n = 0 is the shared pre-blast build and is left out)
\newcommand{\NumDeepDoubleWordsSameNestedInner}{10/10}
% source: e11_blast_depth nested_inner 2 iterations: blasts n >= 1 whose canonical words differ from single blast 2n
\newcommand{\NumDeepDoubleWordsDifferNestedInner}{none}
% source: e11_blast_depth nested_inner 2 iterations: pairs (n, 2n) compared
\newcommand{\NumDeepDoubleComparedNestedInner}{11}
% source: e11_blast_depth nested_inner 2 iterations: blasts done
\newcommand{\NumDeepDoubleBlastsNestedInner}{20}
% source: e11_blast_depth cutoff cells (every case, cutoffs 1e-5, 1e-6, 1e-8)
\newcommand{\NumDeepCutoffCells}{15}
% source: e11_blast_depth cutoff cells whose canonical words (where both completed) differ from the reference at an earlier blast than their crossing counts
\newcommand{\NumDeepCutoffWordsFirst}{12}
% source: e11_blast_depth unresolved-class entries with reason 'dual-graph walk unreachable', summed over reference snapshots after blast 0
\newcommand{\NumDeepUnresolvedUnreachable}{44}
% source: e11_blast_depth unresolved-class entries with reason 'a landing is a singleton ... no registered image chain', summed over reference snapshots after blast 0
\newcommand{\NumDeepUnresolvedSingleton}{2}
% source: e11_blast_depth unresolved-class entries with the singleton / no-registered-chain reason at blast 0, reference cells
\newcommand{\NumDeepUnresolvedSingletonStart}{3}
% source: e11_blast_depth unresolved-class entries with any other reason at blast 0, reference cells
\newcommand{\NumDeepUnresolvedOtherStart}{0}
% source: e11_blast_depth k28 reference: unresolved classes with the 'unreachable' reason at the last completed snapshot
\newcommand{\NumDeepUnresolvedUnreachableEndKTwoEight}{10}
% source: e11_blast_depth k28 reference: unresolved classes with the 'singleton' reason at the last completed snapshot
\newcommand{\NumDeepUnresolvedSingletonEndKTwoEight}{2}
% source: e11_blast_depth large-zone (k28, k32, nested_outer) cells at cutoff 1e-5 or 1e-6
\newcommand{\NumDeepCoarseCells}{6}
% source: e11_blast_depth of those, stopped by the time budget
\newcommand{\NumDeepCoarseTimeStops}{6}
% source: e11_blast_depth of those, blasts done (range)
\newcommand{\NumDeepCoarseBlasts}{16--17}
% source: e11_blast_depth of those, process peak RSS range (GB)
\newcommand{\NumDeepCoarseRss}{3.1--8.5}
% source: e11_blast_depth large-zone cells at cutoff 1e-7 with one iteration per blast
\newcommand{\NumDeepRefCutoffCells}{19}
% source: e11_blast_depth of those, stopped by the time budget (the rest by the memory budget)
\newcommand{\NumDeepRefCutoffTimeStops}{2}
% source: e11_blast_depth of those, stopped by the memory budget
\newcommand{\NumDeepRefCutoffMemoryStops}{17}
% source: e11_blast_depth p3 reference: mean time per blast (s)
\newcommand{\NumDeepMeanBlastSPThree}{0.011}
% source: e11_blast_depth p3 reference: process peak RSS from the watchdog (GiB)
\newcommand{\NumDeepProcRssGbPThree}{0.42}
% source: e11_blast_depth p3 reference: bridge points before blasting
\newcommand{\NumDeepPointsStartPThree}{908}
% source: e11_blast_depth p3 reference: bridge points after the last blast
\newcommand{\NumDeepPointsEndPThree}{4\,332}
% source: e11_blast_depth nested_inner reference: mean time per blast (s)
\newcommand{\NumDeepMeanBlastSNestedInner}{0.013}
% source: e11_blast_depth nested_inner reference: process peak RSS from the watchdog (GiB)
\newcommand{\NumDeepProcRssGbNestedInner}{0.47}
% source: e11_blast_depth nested_inner reference: bridge points before blasting
\newcommand{\NumDeepPointsStartNestedInner}{6\,122}
% source: e11_blast_depth nested_inner reference: bridge points after the last blast
\newcommand{\NumDeepPointsEndNestedInner}{9\,546}
% source: e11_blast_depth k28 E11 reference: crossings at blast 10 (E7's k28 blast count)
\newcommand{\NumDeepAtESevenCrossingsKTwoEight}{712}
% source: e11_blast_depth k28 E11 reference: topology outcome at blast 10
\newcommand{\NumDeepAtESevenOutcomeKTwoEight}{invariant\_assert}

% Generated by stress/numerics/report/make_numbers_extra.py -- do not edit.
% Computed from the full-run results/*.jsonl (and summary.json for e04 convergence).

% source: e01 sweep numpy: batch size N of the peak throughput
\newcommand{\NumMapPeakNumpyN}{\ensuremath{10^{4}}}
% source: e01 sweep numpy: throughput at the largest N (points/s)
\newcommand{\NumMapPlateauNumpy}{\ensuremath{1.7\times10^{8}}}
% source: e01 sweep numpy: peak throughput / throughput at the largest N
\newcommand{\NumMapPeakOverPlateauNumpy}{\ensuremath{5.4\times}}
% source: e01 sweep: map_batch GPU / map_batch CPU throughput at the largest N
\newcommand{\NumMapBatchGpuOverCpuMaxN}{\ensuremath{1.11\times}}
% source: e01 sweep: cupy end-to-end / numpy throughput at the largest N
\newcommand{\NumRawGpuOverCpuMaxN}{\ensuremath{0.87\times}}
% source: e02 profile k28 cpu: refinement share of tottime
\newcommand{\NumProfRefinementKTwoEight}{46\%}
% source: e02 profile k28 cpu: linked-list share of tottime
\newcommand{\NumProfLinkedListKTwoEight}{38\%}
% source: e02 profile k28 cpu: other share of tottime
\newcommand{\NumProfOtherKTwoEight}{17\%}
% source: e02 profile k28 cpu: total profiled seconds
\newcommand{\NumProfTotalSKTwoEight}{83}
% source: e02 profile k28 cpu: ManifoldMachine._refine_layer tottime / total
\newcommand{\NumProfRefineLayerShareKTwoEight}{44\%}
% source: e02 profile k28 cpu: Point.__init__ calls
\newcommand{\NumProfPointInitCallsKTwoEight}{\ensuremath{8.7\times10^{6}}}
% source: e02 profile k28 cpu: Point.__init__ cumtime / total
\newcommand{\NumProfPointInitShareKTwoEight}{18\%}
% source: e02 profile k28 cpu: numpy.array calls
\newcommand{\NumProfNpArrayCallsKTwoEight}{\ensuremath{8.7\times10^{6}}}
% source: e02 profile k28 cpu: refinement + linked-list share of tottime
\newcommand{\NumProfRefPlusListKTwoEight}{83\%}
% source: e02 profile p3 cpu: refinement share of tottime
\newcommand{\NumProfRefinementPThree}{30\%}
% source: e02 profile p3 cpu: linked-list share of tottime
\newcommand{\NumProfLinkedListPThree}{45\%}
% source: e02 profile p3 cpu: other share of tottime
\newcommand{\NumProfOtherPThree}{25\%}
% source: e02 profile p3 cpu: refinement + linked-list share of tottime
\newcommand{\NumProfRefPlusListPThree}{75\%}
% source: e02 profile k10 cpu: refinement share of tottime
\newcommand{\NumProfRefinementKTen}{42\%}
% source: e02 profile k10 cpu: linked-list share of tottime
\newcommand{\NumProfLinkedListKTen}{33\%}
% source: e02 profile k10 cpu: other share of tottime
\newcommand{\NumProfOtherKTen}{24\%}
% source: e02 profile k10 cpu: refinement + linked-list share of tottime
\newcommand{\NumProfRefPlusListKTen}{76\%}
% source: e02 verify k28: CPU unstable points at the last step (12)
\newcommand{\NumGpuVerifyPointsCpuKTwoEight}{4\,348\,557}
% source: e02 verify k28: GPU unstable points at the last step (12)
\newcommand{\NumGpuVerifyPointsGpuKTwoEight}{4\,348\,603}
% source: e02 verify k28: max |coord| at the last step
\newcommand{\NumGpuVerifyMaxCoordKTwoEight}{\ensuremath{5.5\times10^{8}}}
% source: e02 verify k28: first step with a non-zero CPU-GPU deviation
\newcommand{\NumGpuFirstDivergentStepKTwoEight}{8}
% source: e02 verify k28: max |CPU-GPU| coordinate at step 11
\newcommand{\NumGpuVerifyDevStepElevenKTwoEight}{\ensuremath{3.6\times10^{-11}}}
% source: e02 verify k28: max |CPU-GPU| / max|coord| at step 11
\newcommand{\NumGpuVerifyRelDevStepElevenKTwoEight}{\ensuremath{1.5\times10^{-15}}}
% source: e02 timed k28: unstable points after the last growth step
\newcommand{\NumGpuTimedMaxPointsKTwoEight}{\ensuremath{4.3\times10^{6}}}
% source: e03 growth k=2.8 cutoff 1e-07: point growth factor of the last completed step
\newcommand{\NumLastGrowthFactorKTwoEight}{24.0}
% source: e03 growth k=2.8: unstable eigenvalue lambda_u
\newcommand{\NumLambdaUKTwoEight}{5.72}
% source: e03 growth k=2.8 cutoff 1e-07: max|coord| after the last completed step
\newcommand{\NumLastMaxCoordKTwoEight}{\ensuremath{1.3\times10^{6}}}
% source: e03 growth k=2.8 cutoff 1e-07: unstable points after the last completed step
\newcommand{\NumLastPointsKTwoEight}{407\,709}
% source: e03 growth k=2.8 cutoff 1e-07: completed growth steps
\newcommand{\NumGrowthStepsKTwoEight}{12}
% source: e03 growth k=2.8 cutoff 1e-05: point growth factor of the last completed step
\newcommand{\NumLastGrowthFactorKTwoEightCoarse}{23.4}
% source: e03 growth k=2.8: unstable eigenvalue lambda_u
\newcommand{\NumLambdaUKTwoEightCoarse}{5.72}
% source: e03 growth k=2.8 cutoff 1e-05: max|coord| after the last completed step
\newcommand{\NumLastMaxCoordKTwoEightCoarse}{\ensuremath{1.3\times10^{6}}}
% source: e03 growth k=2.8 cutoff 1e-05: unstable points after the last completed step
\newcommand{\NumLastPointsKTwoEightCoarse}{86\,168}
% source: e03 growth k=2.8 cutoff 1e-05: completed growth steps
\newcommand{\NumGrowthStepsKTwoEightCoarse}{12}
% source: e03 growth k=4.0 cutoff 1e-07: point growth factor of the last completed step
\newcommand{\NumLastGrowthFactorKFour}{12.9}
% source: e03 growth k=4.0: unstable eigenvalue lambda_u
\newcommand{\NumLambdaUKFour}{6.31}
% source: e03 growth k=4.0 cutoff 1e-07: max|coord| after the last completed step
\newcommand{\NumLastMaxCoordKFour}{\ensuremath{1.6\times10^{5}}}
% source: e03 growth k=4.0 cutoff 1e-07: unstable points after the last completed step
\newcommand{\NumLastPointsKFour}{89\,077}
% source: e03 growth k=4.0 cutoff 1e-07: completed growth steps
\newcommand{\NumGrowthStepsKFour}{11}
% source: e03 growth k=6.0 cutoff 1e-07: point growth factor of the last completed step
\newcommand{\NumLastGrowthFactorKSix}{61.8}
% source: e03 growth k=6.0: unstable eigenvalue lambda_u
\newcommand{\NumLambdaUKSix}{7.15}
% source: e03 growth k=6.0 cutoff 1e-07: max|coord| after the last completed step
\newcommand{\NumLastMaxCoordKSix}{\ensuremath{3.5\times10^{7}}}
% source: e03 growth k=6.0 cutoff 1e-07: unstable points after the last completed step
\newcommand{\NumLastPointsKSix}{2\,194\,304}
% source: e03 growth k=6.0 cutoff 1e-07: completed growth steps
\newcommand{\NumGrowthStepsKSix}{11}
% source: e03 growth k=10.0 cutoff 1e-07: point growth factor of the last completed step
\newcommand{\NumLastGrowthFactorKTen}{38.6}
% source: e03 growth k=10.0: unstable eigenvalue lambda_u
\newcommand{\NumLambdaUKTen}{8.52}
% source: e03 growth k=10.0 cutoff 1e-07: max|coord| after the last completed step
\newcommand{\NumLastMaxCoordKTen}{\ensuremath{6.6\times10^{6}}}
% source: e03 growth k=10.0 cutoff 1e-07: unstable points after the last completed step
\newcommand{\NumLastPointsKTen}{520\,502}
% source: e03 growth k=10.0 cutoff 1e-07: completed growth steps
\newcommand{\NumGrowthStepsKTen}{10}
% source: e03 growth: configs killed at the 24 GB RSS cap
\newcommand{\NumScalingOomConfigs}{3}
% source: e03 growth oom: fewest unstable points held after the last completed step
\newcommand{\NumScalingOomMinPointsHeld}{86\,168}
% source: e03 growth oom: most unstable points held after the last completed step
\newcommand{\NumScalingOomMaxPointsHeld}{407\,709}
% source: e03 growth oom: process wall minus checkpointed wall = time spent in the killed step (min, s)
\newcommand{\NumScalingOomKilledStepMinS}{190}
% source: e03 growth oom: time spent in the killed step (max, s)
\newcommand{\NumScalingOomKilledStepMaxS}{208}
% source: derived: 24 GiB / median RSS slope 1057.7 B/point (summary.json BytesPerPoint)
\newcommand{\NumPointCeilingTwentyFourGB}{\ensuremath{2.4\times10^{7}}}
% source: e04 ('k28', 10): log-log slope of the p50 invariance error vs cutoff, cutoffs 1e-9..1e-6
\newcommand{\NumInvSlopeKTwoEightTen}{0.70}
% source: e04 ('k28', 10): crossings at cutoff 1e-9
\newcommand{\NumCrossingsKTwoEightTenRef}{5}
% source: e04 ('k28', 10): crossings at cutoff 1e-3
\newcommand{\NumCrossingsKTwoEightTenCutoffThree}{5}
% source: e04 ('k28', 10): crossings at cutoff 1e-4
\newcommand{\NumCrossingsKTwoEightTenCutoffFour}{5}
% source: e04 ('k28', 10): crossings at cutoff 1e-5
\newcommand{\NumCrossingsKTwoEightTenCutoffFive}{5}
% source: e04 ('k28', 12): log-log slope of the p50 invariance error vs cutoff, cutoffs 1e-9..1e-6
\newcommand{\NumInvSlopeKTwoEightTwelve}{0.70}
% source: e04 ('k28', 12): crossings at cutoff 1e-9
\newcommand{\NumCrossingsKTwoEightTwelveRef}{17}
% source: e04 ('k28', 12): crossings at cutoff 1e-3
\newcommand{\NumCrossingsKTwoEightTwelveCutoffThree}{17}
% source: e04 ('k28', 12): crossings at cutoff 1e-4
\newcommand{\NumCrossingsKTwoEightTwelveCutoffFour}{17}
% source: e04 ('k28', 12): crossings at cutoff 1e-5
\newcommand{\NumCrossingsKTwoEightTwelveCutoffFive}{17}
% source: e04 ('p3', 13): log-log slope of the p50 invariance error vs cutoff, cutoffs 1e-9..1e-6
\newcommand{\NumInvSlopePThreeThirteen}{0.64}
% source: e04 ('p3', 13): crossings at cutoff 1e-9
\newcommand{\NumCrossingsPThreeThirteenRef}{27}
% source: e04 ('p3', 13): crossings at cutoff 1e-3
\newcommand{\NumCrossingsPThreeThirteenCutoffThree}{14}
% source: e04 ('p3', 13): crossings at cutoff 1e-4
\newcommand{\NumCrossingsPThreeThirteenCutoffFour}{37}
% source: e04 ('p3', 13): crossings at cutoff 1e-5
\newcommand{\NumCrossingsPThreeThirteenCutoffFive}{27}
% source: e04 ('p3', 17): log-log slope of the p50 invariance error vs cutoff, cutoffs 1e-9..1e-6
\newcommand{\NumInvSlopePThreeSeventeen}{0.66}
% source: e04 ('p3', 17): crossings at cutoff 1e-9
\newcommand{\NumCrossingsPThreeSeventeenRef}{203}
% source: e04 ('p3', 17): crossings at cutoff 1e-4
\newcommand{\NumCrossingsPThreeSeventeenCutoffFour}{249}
% source: e04 ('p3', 17): crossings at cutoff 1e-5
\newcommand{\NumCrossingsPThreeSeventeenCutoffFive}{214}
% source: e04 k28/10: log-log slope of the median crossing distance to the 1e-9 run vs cutoff, cutoffs <= 1e-5
\newcommand{\NumCrossDevSlopeKTwoEightTen}{0.73}
% source: e04 k28/12: log-log slope of the median crossing distance to the 1e-9 run vs cutoff, cutoffs <= 1e-5
\newcommand{\NumCrossDevSlopeKTwoEightTwelve}{0.70}
% source: e04 p3/13: log-log slope of the median crossing distance to the 1e-9 run vs cutoff, cutoffs <= 1e-5
\newcommand{\NumCrossDevSlopePThreeThirteen}{0.66}
% source: e04 p3/17: log-log slope of the median crossing distance to the 1e-9 run vs cutoff, cutoffs <= 1e-5
\newcommand{\NumCrossDevSlopePThreeSeventeen}{0.58}
% source: e05 k10: ||lu ls| - 1|
\newcommand{\NumDetDefectKTen}{\ensuremath{1.2\times10^{-11}}}
% source: e05 k10: max |stored - analytic| Jacobian entry (the k10 recipe passes no analytic Jacobian)
\newcommand{\NumJacDevKTen}{\ensuremath{1.4\times10^{-12}}}
% source: e05 k28 two blasts: |signed area| of the turnstile-sized lobes
\newcommand{\NumLobeAreaKTwoEight}{7.004}
% source: e05 k28 two blasts: lobes of |area| within 0.1 of 7
\newcommand{\NumLobeAreaKTwoEightCount}{2}
% source: e05 k28 two blasts: max |A_image - A| over those lobes
\newcommand{\NumLobeAreaAbsErrKTwoEight}{\ensuremath{5.3\times10^{-6}}}
% source: e05 p3: median |signed lobe area|
\newcommand{\NumLobeAreaPThreeMedian}{\ensuremath{1.1\times10^{-4}}}
% source: e05 p3: max |A_image - A| over the lobes
\newcommand{\NumLobeAbsErrPThreeMax}{\ensuremath{5.4\times10^{-7}}}
% source: e05 p3: lobe pairs compared
\newcommand{\NumLobePairsPThree}{22}
% source: e05: max ||lu ls| - 1| over every case except k10 (analytic Jacobian)
\newcommand{\NumDetDefectMaxAnalytic}{\ensuremath{4.4\times10^{-16}}}
% source: e06 b=1 saddle: 'Max iterations reached' exceptions
\newcommand{\NumSweepMaxIterFails}{10}
% source: e06: smallest k of those
\newcommand{\NumSweepMaxIterKMin}{0.30}
% source: e06: largest k of those
\newcommand{\NumSweepMaxIterKMax}{0.545}
% source: e06 b=1 saddle: configs ok through trellis_pips
\newcommand{\NumSweepSaddleOk}{50}
% source: e06 b=1 saddle ok: configs with 0 non-anchor crossings
\newcommand{\NumSweepZeroCrossings}{26}
% source: e06 b=1 saddle ok: largest k with 0 crossings
\newcommand{\NumSweepZeroCrossingsKMax}{5.93}
% source: e06 b=1 saddle ok: smallest k with >= 1 crossing
\newcommand{\NumSweepFirstCrossingK}{2.50}
% source: e06 b=1 saddle ok: most crossings in any config
\newcommand{\NumSweepMaxCrossings}{3}
% source: e06 b=1 other point: 'not a saddle' exceptions
\newcommand{\NumSweepOtherElliptic}{6}
% source: e06: largest k of those
\newcommand{\NumSweepOtherEllipticKMax}{2.63}
% source: e06 b=1 other point: oom at the 4 GB cap
\newcommand{\NumSweepOtherOom}{3}
% source: e06: smallest k of those
\newcommand{\NumSweepOtherOomKMin}{10.2}
% source: e06 b=-1: exceptions
\newcommand{\NumSweepBMinusOneFails}{8}
% source: e06 b in {0.3,0.5,0.9}: ok / configs
\newcommand{\NumSweepDissipativeOk}{15/15}
% source: e07 blasts p3: crossings before blasting
\newcommand{\NumBlastCrossingsStartPThree}{30}
% source: e07 blasts p3: crossings after the last blast
\newcommand{\NumBlastCrossingsEndPThree}{75}
% source: e07 blasts p3: time of the last blast (s)
\newcommand{\NumBlastLastSPThree}{0.0093}
% source: e07 blasts p3: bridge classes after the last blast
\newcommand{\NumBlastClassesPThree}{31}
% source: e07 blasts p3: unresolved classes
\newcommand{\NumBlastUnresolvedPThree}{1}
% source: e07 blasts p3: virtual classes
\newcommand{\NumBlastVirtualPThree}{1}
% source: e07 blasts nested: crossings before blasting
\newcommand{\NumBlastCrossingsStartNested}{34}
% source: e07 blasts nested: crossings after the last blast
\newcommand{\NumBlastCrossingsEndNested}{276}
% source: e07 blasts nested: time of the last blast (s)
\newcommand{\NumBlastLastSNested}{1.5}
% source: e07 blasts nested: bridge classes after the last blast
\newcommand{\NumBlastClassesNested}{13}
% source: e07 blasts nested: unresolved classes
\newcommand{\NumBlastUnresolvedNested}{0}
% source: e07 blasts nested: virtual classes
\newcommand{\NumBlastVirtualNested}{0}
% source: e07 blasts k28: crossings before blasting
\newcommand{\NumBlastCrossingsStartKTwoEight}{4}
% source: e07 blasts k28: crossings after the last blast
\newcommand{\NumBlastCrossingsEndKTwoEight}{712}
% source: e07 blasts k28: time of the last blast (s)
\newcommand{\NumBlastLastSKTwoEight}{5.6}
% source: e07 blasts k28: bridge classes after the last blast
\newcommand{\NumBlastClassesKTwoEight}{32}
% source: e07 blasts k28: unresolved classes
\newcommand{\NumBlastUnresolvedKTwoEight}{7}
% source: e07 blasts k28: virtual classes
\newcommand{\NumBlastVirtualKTwoEight}{2}
% source: e07 escape fixed: process wall minus checkpointed wall (time inside the killed step, s)
\newcommand{\NumEscapeKilledStepS}{175}
% source: e07 depth p3 18: unstable points after step 17, before the kill
\newcommand{\NumPThreePointsSeventeen}{584\,747}
% source: e07 depth p3 18: time of growth step 17 (s)
\newcommand{\NumPThreeStepSeventeenS}{7.2}
% source: e07 depth inversion 7: unstable points after step 6, before the kill
\newcommand{\NumInversionPointsSixSteps}{49\,417}
% source: e07 depth p3: deepest run with is_reliable = True
\newcommand{\NumDepthReliableMaxPThree}{14}
% source: e07 depth p3: shallowest run with is_reliable = False
\newcommand{\NumDepthUnreliableMinPThree}{15}
% source: e07 depth p3: unstable steps whose run hit the RSS cap
\newcommand{\NumDepthOomPThree}{18}
% source: e07 depth inversion: unstable steps whose topology stack raised an invariant AssertionError
\newcommand{\NumInversionAssertSteps}{4 and 5}
% source: e07 depth inversion: unstable steps whose run hit the RSS cap
\newcommand{\NumInversionOomSteps}{7 and 8}
% source: e07 tangency: configs (24 k values x 3 cutoffs)
\newcommand{\NumTangencyConfigs}{72}
% source: e07 tangency k+1=1e-4 cutoff 1e-07: non-anchor crossings
\newcommand{\NumTangencyCrossingsNearCutoffSeven}{25}
% source: e07 tangency cutoff 1e-07: most non-anchor crossings over k
\newcommand{\NumTangencyCrossingsPeakCutoffSeven}{149}
% source: e07 tangency k+1=1e-4 cutoff 1e-10: non-anchor crossings
\newcommand{\NumTangencyCrossingsNearCutoffTen}{153}
% source: e07 tangency cutoff 1e-10: most non-anchor crossings over k
\newcommand{\NumTangencyCrossingsPeakCutoffTen}{1\,607}
% source: e07 tangency k+1=1e-4 cutoff 1e-12: non-anchor crossings
\newcommand{\NumTangencyCrossingsNearCutoffTwelve}{499}
% source: e07 tangency cutoff 1e-12: most non-anchor crossings over k
\newcommand{\NumTangencyCrossingsPeakCutoffTwelve}{4\,401}
% source: e07 tangency: smallest k+1 from which all three cutoffs give the same crossing count
\newcommand{\NumTangencyAgreeKPlusOne}{0.16}
% source: e07 tangency k+1 = 1: non-anchor crossings (every cutoff)
\newcommand{\NumTangencyCrossingsFar}{7}
% source: e07 tangency k+1=1e-4: heuristic splitting reference exp(-pi^2/ln lambda)
\newcommand{\NumTangencyRefNear}{\ensuremath{4.6\times10^{-31}}}
% source: e07 tangency k+1=0.025: heuristic splitting reference
\newcommand{\NumTangencyRefMid}{\ensuremath{1.9\times10^{-8}}}
% source: e07 tangency cutoff 1e-12, k+1 <= 0.025: smallest per-config min |sin|
\newcommand{\NumTangencyMinSinFineLo}{\ensuremath{2.0\times10^{-6}}}
% source: e07 tangency cutoff 1e-12, k+1 <= 0.025: largest per-config min |sin|
\newcommand{\NumTangencyMinSinFineHi}{\ensuremath{4.9\times10^{-5}}}
% source: e07 tangency cutoff 1e-12, k+1 <= 0.025: median over configs of the median |sin|
\newcommand{\NumTangencyMedianSinFine}{\ensuremath{4.8\times10^{-4}}}
% source: e07 tangency: WARNING records of any kind, all configs
\newcommand{\NumTangencyWarningsTotal}{0}
% source: e08: solves with a converged orbit (minimal period measured)
\newcommand{\NumSolverConverged}{8\,450}
% source: e08: lower-period solves at requested period 2 with k < 3 (no period-2 orbit)
\newcommand{\NumSolverLowerNoOrbit}{1\,198}
% source: e08: those / converged
\newcommand{\NumSolverLowerNoOrbitRate}{14.2\%}
% source: e08: lower-period solves where the requested period has orbits
\newcommand{\NumSolverLowerSlide}{1\,744}
% source: e08: those / converged
\newcommand{\NumSolverLowerSlideRate}{20.6\%}
% source: e08 k = 10: lower-period / converged
\newcommand{\NumSolverLowerSlideRateKTen}{20.1\%}
% source: e08: lower-period solves that construct_fixed_point accepted as a saddle of the requested period
\newcommand{\NumSolverLowerAccepted}{1\,424}
% source: e08: median solve time (ms)
\newcommand{\NumSolverMedianMs}{0.67}
% source: e08: p99 solve time (ms)
\newcommand{\NumSolverPNinetyNineMs}{7.2}
% source: e09: configs where the brute force beat the library
\newcommand{\NumIsectBruteFasterConfigs}{23}
% source: e09: segments of the largest config
\newcommand{\NumIsectLargestSegments}{\ensuremath{4.1\times10^{5}}}
% source: e09 largest config: library time (s)
\newcommand{\NumIsectLargestLibraryS}{6.6}
% source: e09 largest config: brute-force time (s)
\newcommand{\NumIsectLargestBruteS}{9.0}
% source: e09: configs with at least one u x u crossing
\newcommand{\NumIsectUUConfigs}{2}
% source: e09: area_cutoff of every config with u x u crossings
\newcommand{\NumIsectUUCutoff}{\ensuremath{10^{-5}}}
% source: e10: median summed stage time, CPU (s)
\newcommand{\NumDetTotalCpuS}{0.112}
% source: e10: median summed stage time, GPU (s)
\newcommand{\NumDetTotalGpuS}{0.132}
% source: e10: CPU / GPU median summed stage time
\newcommand{\NumDetGpuSpeedup}{\ensuremath{0.85\times}}
% source: e01_map_throughput.jsonl: records
\newcommand{\NumConfigsEOne}{8}
% source: e02_gpu_end_to_end.jsonl: records
\newcommand{\NumConfigsETwo}{33}
% source: e03_scaling.jsonl: records
\newcommand{\NumConfigsEThree}{22}
% source: e04_area_cutoff.jsonl: records
\newcommand{\NumConfigsEFour}{28}
% source: e05_invariants.jsonl: records
\newcommand{\NumConfigsEFive}{8}
% source: e06_parameter_sweep.jsonl: records
\newcommand{\NumConfigsESix}{98}
% source: e07_breaking.jsonl: records
\newcommand{\NumConfigsESeven}{92}
% source: e08_fixed_point_solver.jsonl: records
\newcommand{\NumConfigsEEight}{24}
% source: e09_intersections.jsonl: records
\newcommand{\NumConfigsENine}{26}
% source: e10_determinism.jsonl: records
\newcommand{\NumConfigsETen}{10}
% source: e11_blast_depth.jsonl: records
\newcommand{\NumConfigsEEleven}{53}
% source: all full-run JSONL files (e01-e11): records
\newcommand{\NumConfigsTotal}{402}
% source: e01-e10 full-run JSONL files: records of the scripted campaign (e11 ran on its own)
\newcommand{\NumConfigsCampaign}{349}
% source: progress.log: wall-clock of the last --quick campaign (min)
\newcommand{\NumCampaignQuickMin}{1.2}
% source: progress.log: wall-clock of the full campaign e01-e10 (min; e11 ran on its own, see NumDeepWallMin)
\newcommand{\NumCampaignFullMin}{59.8}
% source: all full-run JSONL files: outcome ok
\newcommand{\NumOutcomeTotalOk}{367}
% source: all full-run JSONL files: outcome exception
\newcommand{\NumOutcomeTotalException}{24}
% source: all full-run JSONL files: outcome oom
\newcommand{\NumOutcomeTotalOom}{11}
% source: all full-run JSONL files: outcome timeout
\newcommand{\NumOutcomeTotalTimeout}{0}
% source: all full-run JSONL files: outcome invariant_assert
\newcommand{\NumOutcomeTotalInvariantAssert}{0}


\title{\textbf{Stress-testing the \texttt{tanglepack} numerics layer}\\[0.6em]
\large What the manifold-growth engine can and cannot do, with evidence}
\author{tanglepack numerics stress campaign}
\date{2026-10-07}

\begin{document}

\pagenumbering{roman}
\hypersetup{pageanchor=false}
\begin{titlepage}
\maketitle
\thispagestyle{empty}
\vfill
\begin{center}
\begin{minipage}{0.82\textwidth}
\small
\noindent This report presents a \NumCampaignFullMin-minute stress campaign of
\NumConfigsCampaign{} isolated runs (ten experiments, E1--E10), followed by a
\NumDeepWallMin-minute deep-blasting study of \NumConfigsEEleven{} runs (E11),
\NumConfigsTotal{} runs in all, against
\code{tanglepack.numerics}, the layer that grows the stable and unstable
manifolds of saddle orbits of area-preserving maps, detects their crossings
and builds bridges. Each section follows the same pattern: a claim, the figure
or table that supports it, an interpretation, and a caveat. Negative results
count as findings. Every measured number in the text is generated from the raw
result files (\code{numbers.tex}, \code{numbers\_extra.tex}); none is typed by
hand.

\medskip
\noindent Branch \code{numerics-stress}, library commit \code{e963893}. The
library itself was not modified.
\end{minipage}
\end{center}
\vfill
\end{titlepage}
\hypersetup{pageanchor=true}

\tableofcontents
\clearpage
\pagenumbering{arabic}

%======================================================================
\section{Executive summary}
\label{sec:summary}

\paragraph{Six headline findings.}
\begin{enumerate}[leftmargin=*]
\item \textbf{The GPU does not help, and cannot help in the current design.}
  The end-to-end speed-up is \NumGpuTotalSpeedupMedianKTwoEight{} (k28),
  \NumGpuTotalSpeedupMedianPThree{} (p3) and \NumGpuTotalSpeedupMedianKTen{}
  (k10). The map takes at most \NumMapFractionMax{} of growth time, so
  Amdahl's law caps the gain at \NumAmdahlBoundMax{}
  (\cref{sec:gpu}).
\item \textbf{Cost is linear in the point count, but the point count is not
  set by the dynamics.} Step time scales as $N^{\NumStepTimeExponent}$ and
  memory as \NumBytesPerPoint{}\,B/point. However, once the unstable arm
  escapes to infinity a single step multiplies the points by up to
  \NumLastGrowthFactorKSix{}, against $\lambda_u = \NumLambdaUKSix$
  (\cref{sec:scaling}).
\item \textbf{\code{area\_cutoff} is a well-behaved accuracy knob.} The
  invariance error falls as $\text{cutoff}^{\NumInvSlopeKTwoEightTen}$ and the
  point count grows as $\text{cutoff}^{\NumPointsCutoffExponentKTwoEight}$
  (k28), close to the $2/3$ and $-1/3$ that chord geometry predicts. At $10^{-7}$ the median invariance error on k28 is
  \NumInvPFiftyKTwoEightCutoffSeven{} for \NumTimeKTwoEightCutoffSeven\,s of
  pipeline time. The library default, $10^{-4}$, is too coarse for the
  period-3 case: it invents and drops crossings (\cref{sec:accuracy}).
\item \textbf{Transverse crossings are found exactly, and the physical
  invariants hold.} An independent brute force matches all
  \NumIsectMatchedTotal{}/\NumIsectBruteTotal{} crossings with
  \NumIsectMissedTotal{} missed and \NumIsectExtraTotal{} extra. Lobe areas
  are preserved to \NumLobeRatioMax{}, and \NumCdistDecreasingManifolds{}
  cdist decreases occur in \NumCdistStepsManifolds{} steps
  (\cref{sec:invariants,sec:intersections}).
\item \textbf{Two failures are silent.} (a) Near the saddle-centre
  bifurcation, at fixed $k+1=10^{-4}$ and a fixed manifold length, the number
  of registered crossings grows from
  \NumTangencyCrossingsNearCutoffSeven{} to
  \NumTangencyCrossingsNearCutoffTwelve{} as the resolution is refined, with
  \NumTangencyWarningsTotal{} warnings. These are artefacts of the polygons,
  not dynamics. (b) The fixed-point solver hands back lower-period orbits:
  \NumSolverLowerAccepted{} of them were accepted as saddles of the requested
  period (\cref{sec:breaking,sec:solver}).
\item \textbf{Deep blasting keeps uncovering structure, no crossing or cdist
  law is broken, and the topology layer has four genuine errors to triage.}
  Over \NumDeepBlasts{} blasts in \NumDeepConfigs{} configurations there were
  \NumDeepLawViolations{} same-stability crossings or cdist decreases; the
  lobe-area check is open in the two-iteration cells. The canonical words changed in
  \NumDeepStructureSteps{} snapshots (k28: \NumDeepClassesStartKTwoEight{}
  $\to$ \NumDeepClassesEndKTwoEight{} classes), and memory limits the large
  zones to \NumDeepBlastsKThreeTwo--\NumDeepBlastsKTwoEight{} blasts at
  $\times\NumDeepPointGrowthKTwoEight$ points per blast (k28). The topology
  raised in \NumDeepErrorSnapshots{} of \NumDeepSnapshots{} snapshots
  (\NumDeepErrorConfigs{} configurations), with \NumDeepErrorClasses{}
  distinct errors, the most frequent being a zero-width iterated element
  (\cref{sec:blastdepth}).
\end{enumerate}

\noindent \Cref{fig:overview} puts one panel per theme side by side, and
\cref{tab:scorecard} grades each capability with a pointer to its
evidence. \Cref{sec:taxonomy} catalogues the failures, and
\cref{sec:recs} expands the recommendations below.

\begin{figure}[H]
  \centering
  \includegraphics[width=0.8\textwidth]{\fig{overview}}
  \caption{Campaign overview. (a)~Map throughput: NumPy beats the end-to-end
  CuPy path at every batch size, and the \code{map\_batch} GPU path only
  catches up at the largest batches (\cref{sec:gpu}). (b)~End-to-end speed-up per case:
  every bar is at or below 1. (c)~Invariance error falls steadily as
  \code{area\_cutoff} is refined; solid is the median, dotted the 99th
  percentile. (d)~Pass fractions of the main checks. The
  short bars are the solver's requested-period rate and the CPU/GPU hash
  agreement; both are explained in \cref{sec:solver,sec:repro}.}
  \label{fig:overview}
\end{figure}

\begin{table}[H]
\centering\small
\caption{Scorecard. Status: \works{} = does what it should over the tested
range; \limited{} = works, with a documented restriction or with genuine
errors~(iii) in a stated corner; \fails{} = does not deliver. Topology
change~(i) and provisional flags~(ii) never count against a capability
(\cref{sec:setup}).}
\label{tab:scorecard}
\begin{tabularx}{\textwidth}{@{}l l Y@{}}
\toprule
capability & status & evidence \\
\midrule
Manifold growth, stop-condition driven & \works & escape case stops at \NumEscapeStopArclengthPoints{} points (\cref{sec:escape}) \\
Manifold growth, fixed step count & \fails & killed at the 20\,GB cap after \NumEscapeStepsDone{} steps (\cref{sec:escape}) \\
Cost scaling & \works & time $\propto N^{\NumStepTimeExponent}$, \NumBytesPerPoint{}\,B/point (\cref{sec:scaling}) \\
Accuracy control (\code{area\_cutoff}) & \works & error $\propto$ cutoff$^{\NumInvSlopeKTwoEightTen}$ (\cref{sec:accuracy}) \\
Transverse crossing detection & \works & \NumIsectMatchedTotal/\NumIsectBruteTotal{} matched, \NumIsectExtraTotal{} extra (\cref{sec:intersections}) \\
Crossings near tangency & \fails & count tracks resolution, \NumTangencyWarningsTotal{} warnings (\cref{sec:tangency}) \\
Self-crossing detection & \limited & \NumIsectUUTotal{} u$\times$u crossings exist, none reported (\cref{sec:intersections}) \\
Area preservation, invariants & \works & lobe $|A'/A-1|\le\NumLobeRatioMax$ (\cref{sec:invariants}) \\
Arrangement \code{image\_of} & -- & \NumArrImageFound{} images found, \NumArrImageMissing{} missing; open question, not graded (\cref{sec:invariants}) \\
Blasting (numerics) & \works & \NumDeepBlasts{} blasts in \NumDeepConfigs{} configs, \NumDeepLawViolations{} crossing/cdist law violations (area probe open in the two-iteration cells), deterministic to the hash; crossings \NumDeepCrossingsStartKTwoEight{}$\to$\NumDeepCrossingsEndKTwoEight{} and classes \NumDeepClassesStartKTwoEight{}$\to$\NumDeepClassesEndKTwoEight{} (k28) (\cref{sec:blasts,sec:blastdepth}) \\
Blast depth, large zones & \limited & points $\times\NumDeepPointGrowthKTwoEight$ per blast; at cutoff $10^{-7}$ the memory budget ends k28 at \NumDeepBlastsKTwoEight{} blasts (\NumDeepPeakRssGbKTwoEight\,GB), at coarser cutoffs the time budget; p3 runs all \NumDeepMaxBlasts{} in \NumDeepProcRssGbPThree\,GB (\cref{sec:blastdepth-cost}) \\
Topology under deep blasting & \limited & \NumDeepErrorSnapshots/\NumDeepSnapshots{} snapshots raised (\NumDeepErrorClasses{} genuine errors, \NumDeepErrorConfigs/\NumDeepConfigs{} configs); the other snapshots completed, some with provisional symbolic results (open, needs more growth) (\cref{sec:blastdepth-errors}) \\
Deep growth (p3) & \limited & works through \NumDepthNumericsMaxPThree{} steps, the tangle growing with each step; OOM (genuine error) at \NumDepthOomPThree{} (\cref{sec:depth}) \\
Inversion saddle & \limited & topology \code{AssertionError} at \NumInversionAssertSteps{} steps and OOM at \NumInversionOomSteps{} (genuine errors); at \NumDepthNumericsMaxInversion{} steps it completes with a provisional symbolic result (\cref{sec:depth}) \\
Parameter robustness (b = 1 saddle) & \limited & \NumSaddleSuccessFraction{} complete; \NumSweepZeroCrossings{} trivial tangles (\cref{sec:robustness}) \\
Fixed-point solver & \limited & converges \NumSolverConvergedRate{}, but accepts lower periods (\cref{sec:solver}) \\
GPU acceleration & \fails & \NumGpuTotalSpeedupMedianKTen{}--\NumGpuTotalSpeedupMedianKTwoEight{} end to end (\cref{sec:gpu}) \\
Determinism & \works & CPU--CPU and GPU--GPU bit-identical (\cref{sec:repro}) \\
\bottomrule
\end{tabularx}
\end{table}

\paragraph{What I would do next, in priority order.}
\begin{enumerate}[leftmargin=*]
\item \textbf{Guard the crossing registration near tangency.} Warn when a
  crossing's $|\sin\theta|$ is not large compared to the local invariance
  error, and add a cheap u$\times$u self-crossing check as a resolution
  diagnostic. Both failures are currently silent.
\item \textbf{Make the solver check the minimal period} before
  \code{construct\_fixed\_point} accepts an orbit.
\item \textbf{Put a point/memory budget into growth.} Predict the next step's
  points from the last growth factor, as the E3 harness already does, and
  make stop-condition growth the documented default.
\item \textbf{Raise the default \code{area\_cutoff} to $10^{-7}$}, which is
  what the k28, p3 and nested test cases already use.
\item \textbf{Triage the four E11 topology errors} (\cref{sec:blastdepth-errors}),
  starting with the zero-width iterated element, which hits every k28 and
  k32 cell at cutoff $10^{-7}$.
\item \textbf{Leave the GPU off.} Revisit it only if refinement and the
  linked list move to arrays that live on the device.
\end{enumerate}

%======================================================================
\section{Setup and methods}
\label{sec:setup}

\paragraph{The map and the pipeline.} All runs use the Hénon family
(parameters $k$, $b$; $b=1$ is area preserving) defined in the module
\code{tanglepack.examples.henon}, via the frozen recipes of \code{tests/cases.py}. These
recipes are: k10, the $k=10$ binary horseshoe at the default cutoff
$10^{-4}$; k28, $k=2.8$ at $10^{-7}$, optionally blasted; p3, the period-3
orbit of the nested map; nested, period 1 plus period 3; and inversion, the
other saddle of $k=10$, with both eigenvalues negative. The pipeline under
test is \code{TangleWorkbench}: find the fixed point, then initialise the
fundamental segments, grow the manifolds (\code{grow\_n\_times},
\code{grow\_until\_*}), and finally run \code{compute\_intersections}, trim,
\code{create\_bridges} and the iterate table. Where an experiment touches
topology (E5, E7), it is only as a downstream consumer of the numerics.

\paragraph{Topology change, provisional results and errors.} Blasting a zone
or growing the manifolds further is \emph{meant} to change the topology:
new crossings are registered, new bridges, holes and bridge classes appear,
elements split, and words, letters and transition matrices change. This is
new information about the tangle, and the report never counts it as a
failure. The report keeps three things apart.
(i)~\emph{Topology change}: the expected outcome of more growth or more
blasts. (ii)~\emph{Provisional symbolic results}: the library's own flags
(\code{is\_reliable=False}, ambiguous walks, unresolved classes, virtual
classes \code{new1}, \ldots). They mean the symbolic computation at that
stage is incomplete: some class images are not registered yet, or a
dual-graph walk cannot yet be completed, and they call for more growth or
blasting. They are open questions, not
numerical errors. (iii)~\emph{Errors}: exceptions, invariant
\code{AssertionError}s, memory or time exhaustion, and numerics that violate
a physical law (a u$\times$u or s$\times$s crossing, area not preserved,
cdist decreasing). Topology change~(i) and provisional flags~(ii) never
count against a capability; only~(iii) does. The scorecard grade also
reflects whether the capability delivers (e.g.\ the GPU). A crossing count that
depends on \code{area\_cutoff} at \emph{fixed} depth is a resolution error
(\cref{sec:accuracy,sec:tangency}); a count that grows with depth or blasts
is~(i).

\paragraph{Harness.} \code{stress/numerics/harness.py} runs every
configuration in a fresh subprocess with a hard timeout. A psutil watchdog
samples the worker's resident set every 50\,ms and kills the worker above an
RSS cap. \code{RLIMIT\_AS} is not used because CUDA reserves huge virtual
ranges. Each configuration writes one JSON line with its parameters,
per-stage wall times (with a device synchronise before the clock on the
GPU), per-step checkpoints, metrics, captured \code{logging} warnings by
template, peak RSS and git SHA. The outcome is one of \code{ok},
\code{exception}, \code{invariant\_assert}, \code{oom}, \code{timeout} or
\code{crash}. \emph{Failures are data}: a sweep never stops on one, and the
per-step checkpoints keep everything up to the failure. The per-experiment
caps are in \cref{tab:counts}. Timing experiments (E1--E4, E10) ran serially
with the GPU never shared; the robustness sweeps (E5--E9) ran in worker
pools.

\paragraph{Hardware and software.} One host (\code{pop-os}): Intel Core
i9-14900K (24 cores, 8P + 16E, 32 threads, 32\,MB L2, 36\,MB L3), 125\,GB
RAM; NVIDIA GeForce RTX 5060 Ti, 16\,GB, driver 580.173.02; Linux 7.1.1.
Python 3.10.12, NumPy 2.2.5, SciPy 1.15.1, CuPy 14.1.1, rtree 1.4.0. Every
record carries library commit \code{e963893}. The full campaign took
\NumCampaignFullMin\,min of wall-clock, after a \NumCampaignQuickMin-minute
\code{--quick} pilot; E11 ran on its own afterwards in \NumDeepWallMin\,min.

\paragraph{Experiments.}
\begin{description}[leftmargin=1.2em,style=nextline]
\item[E1 Raw map throughput] NumPy against CuPy (end to end, compute only,
  transfer only) and the real \code{map\_batch} path, $N=10^2$--$10^7$.
\item[E2 End-to-end GPU] k28, p3 and k10, CPU against GPU, three repeats;
  also a cProfile breakdown, a sweep of \code{min\_batch\_points} and a
  lockstep CPU/GPU comparison.
\item[E3 Scaling] Growth step by step to $\sim2\times10^6$ points for
  $k\in\{2.8,4,6,10\}$, and downstream stages as a function of the segment
  count.
\item[E4 Accuracy against cost] \code{area\_cutoff} $10^{-3}$--$10^{-9}$ on
  k28 and p3 at two depths, against the finest run.
\item[E5 Invariants] Fixed point, manifold invariance, iterate links, lobe
  areas, cdist monotonicity and \code{image\_of}, on every frozen case.
\item[E6 Robustness] One fixed recipe over $k\in[0.3,15]$ for $b=1$ (both
  period-1 points), $b=-1$ and dissipative $b$.
\item[E7 Limits and deep exploration] Escape to infinity and near-tangency
  at $k\to-1$, where the numerics are expected to struggle; 0--10 blasts and
  deep growth, where the tangle is expected to reveal new structure.
\item[E8 Solver] 200 random seeds $\times$ 2 guess styles, periods 1--6, four
  maps.
\item[E9 Crossing correctness] The library's crossings against an exact
  O($NM$) brute force.
\item[E10 Determinism] Five fresh-process repeats on the CPU and five on the
  GPU.
\item[E11 Deep blasting] Five zones blasted up to \NumDeepMaxBlasts{} times,
  with the full topology stack snapshotted after every blast, over
  \code{min\_separation}, blasting \code{area\_cutoff}, iterations per blast
  and repeats (\cref{sec:blastdepth}).
\end{description}

%======================================================================
\section{Performance and scaling (E3)}
\label{sec:scaling}

\claim Growth, crossing detection and bridge construction all cost
essentially linear time in the number of points, and memory costs about
\NumBytesPerPoint\,B per point. The binding constraint is not the algorithm but the number
of points the escaping unstable arm generates.

\evidence \Cref{fig:e03cost}(a) shows the time of one growth step against
the points it produces, for four values of $k$ and two cutoffs. Above
$10^3$ points the fitted slopes run from \NumStepTimeExponentKSix{} to
\NumStepTimeExponentKTwoEightCoarse{} (median \NumStepTimeExponent). The
downstream stages in \cref{fig:e03cost}(b) scale with the segment count with
slopes \NumIntersectionsExponent{} (\code{compute\_intersections}) and
\NumTrimBridgesExponent{} (trim and bridges). The iterate table is flat
(\NumIterateTableExponent) because it depends on the crossings, not on the
points. Resident memory (\cref{fig:e03mem}) grows by \NumBytesPerPoint\,B per
point (median over runs; maximum \NumBytesPerPointMax).

\begin{figure}[htbp]
  \centering
  \includegraphics[width=0.88\textwidth]{\fig{e03_cost}}
  \caption{E3 cost. (a)~Growth-step time against unstable points: all $k$
  and both cutoffs collapse onto one line with slope near 1 (legend), so cost
  depends on the point count alone. (b)~Downstream stages against segments:
  intersections and trim/bridges are linear; the iterate table does not
  depend on the point count.}
  \label{fig:e03cost}
\end{figure}

\begin{figure}[htbp]
  \centering
  \includegraphics[width=0.88\textwidth]{\fig{e03_memory}}
  \caption{E3 memory. (a)~RSS above the interpreter baseline: linear over the
  top two decades, with nearly the same slope for every $k$ (legend).
  (b)~Bytes per point: fixed interpreter overhead dominates small manifolds;
  the curve flattens onto the per-point slope for large ones.}
  \label{fig:e03mem}
\end{figure}

\Cref{fig:e03growth} explains why the point count is the real problem. For
the first steps the points grow roughly like $|\lambda_u|^n$. Then, at step
\NumEscapeStepKTen{} ($k=10$) to \NumEscapeStepKTwoEight{} ($k=2.8$), the
largest coordinate passes 100 and the growth factor departs from
$\lambda_u$. On its last completed step the $k=6$ manifold grew by a factor
\NumLastGrowthFactorKSix{} ($\lambda_u=\NumLambdaUKSix$), $k=10$ by
\NumLastGrowthFactorKTen{} ($\lambda_u=\NumLambdaUKTen$) and $k=2.8$ by
\NumLastGrowthFactorKTwoEight{} ($\lambda_u=\NumLambdaUKTwoEight$). Meanwhile
the largest coordinate reached \NumLastMaxCoordKSix{} and
\NumLastMaxCoordKTen.

\begin{figure}[htbp]
  \centering
  \includegraphics[width=0.88\textwidth]{\fig{e03_growth}}
  \caption{E3 growth. (a)~Unstable points per step against the
  $|\lambda_u|^n$ guide (dotted). The curves follow it for a few steps and
  then shoot above it; crosses mark runs killed at the 24\,GB cap. (b)~The
  largest coordinate: the take-off in (a) coincides with the arm leaving for
  infinity.}
  \label{fig:e03growth}
\end{figure}

\reading The refinement places points to keep a fixed chord area, so the
point count follows the arclength and curvature of the image, not the
eigenvalue. Once the tip of the arm is at $|x|\sim10^3$, one step of a
quadratic map sends it to $\sim10^6$, and the refinement faithfully resolves
an arc that is irrelevant to the tangle. ``More steps'' therefore mostly
buys points far outside the tangle. \NumScalingOomConfigs{} growth runs
were killed at the 24\,GB cap. At that moment they held only
\NumScalingOomMinPointsHeld--\NumScalingOomMaxPointsHeld{} points after their
last completed step, but they had spent
\NumScalingOomKilledStepMinS--\NumScalingOomKilledStepMaxS\,s inside the
step that was killed.

\caveat The \NumBytesPerPoint\,B/point slope gives a ceiling of roughly
\NumPointCeilingTwentyFourGB{} points in a 24\,GB budget. This is an
\emph{estimate}: it extrapolates the steady-state slope, and a step's
transient working set comes on top of it. The exponents are fitted over
$\le$3 decades on one machine.

%======================================================================
\section{GPU acceleration (E1, E2)}
\label{sec:gpu}

\claim Enabling the GPU (\code{enable\_gpu}) does not speed up growth.
Today it is neutral at best and slower at worst, and it cannot pay off while
the map is a negligible share of the work.

\subsection{Raw throughput}
\evidence \Cref{fig:e01}(a): NumPy peaks at \NumMapPeakNumpy\,points/s near
$N=\NumMapPeakNumpyN$, then drops by \NumMapPeakOverPlateauNumpy{} to
\NumMapPlateauNumpy\,points/s. This is consistent with a cache effect: the
working set leaves the per-core L2. CuPy on data already on the device
reaches \NumMapPeakCupyCompute\,points/s, but the end-to-end path (host
$\to$ device $\to$ host, which is what \code{gpu.py} does) peaks at
\NumMapPeakCupyEndToEnd. At $N=10^7$, transfer is
\NumGpuTransferFractionAtMaxN{} of the GPU time (\cref{fig:e01}(b)). The raw
crossover: \NumGpuCrossoverNRaw{} (end-to-end CuPy reaches
\NumRawGpuOverCpuMaxN{} of NumPy at the largest $N$). Through
\code{map\_batch}, whose CPU path pays for transposes and copies, the GPU
breaks even somewhere between \NumGpuCrossoverNSweepMapBatch{} (interpolated
from the sweep) and \NumGpuCrossoverNMapBatch{} (bisection). At $10^7$ it is
ahead by only \NumMapBatchGpuOverCpuMaxN.

\begin{figure}[htbp]
  \centering
  \includegraphics[width=0.88\textwidth]{\fig{e01_map_throughput}}
  \caption{E1. (a)~Throughput against batch size. Device-only compute (dotted
  triangles) is fastest, but every path that includes the transfer sits on
  the transfer-only line (grey crosses). The dashed line marks the bisected
  \code{map\_batch} crossover; read it as the upper edge of a band, not a
  sharp threshold. (b)~Share of GPU end-to-end time: transfer dominates
  from mid-size batches on.}
  \label{fig:e01}
\end{figure}

\subsection{End to end}
\evidence \Cref{fig:e02} and \cref{tab:e02}: the summed pipeline time on
the GPU is \NumGpuTotalSpeedupMedianKTwoEight{} (k28),
\NumGpuTotalSpeedupMedianPThree{} (p3) and \NumGpuTotalSpeedupMedianKTen{}
(k10) of the CPU speed (median of three repeats; best repeats
\NumGpuTotalSpeedupBestKTwoEight, \NumGpuTotalSpeedupBestPThree,
\NumGpuTotalSpeedupBestKTen). The best single growth step, k28 at
\NumGpuTimedMaxPointsKTwoEight{} points, gains \NumGpuStepSpeedupBestKTwoEight.
On the E10 case the GPU is clearly slower: \NumDetGpuSpeedup{}
(\cref{sec:repro}).

\begin{figure}[htbp]
  \centering
  \includegraphics[width=0.88\textwidth]{\fig{e02_end_to_end}}
  \caption{E2. (a)~Per-step growth time; CPU (filled) and GPU (open) markers
  lie on top of each other up to \NumGpuTimedMaxPointsKTwoEight{} points. (b)~Per-stage
  speed-up. Everything sits within a few percent of 1 except tiny stages,
  where sub-millisecond noise dominates (p3 iterate table, k28 stable
  growth).}
  \label{fig:e02}
\end{figure}

\begin{table}[htbp]
\centering\footnotesize
\caption{E2: median CPU time / median GPU time per stage, the CPU map share
$f_{\rm map}$ (cProfile) and the Amdahl bound $1/(1-f_{\rm map})$.}
\label{tab:e02}
\setlength{\tabcolsep}{4pt}
\begin{tabular}{lrrrrrrrr}
\tablebody tables/e02_speedup.tex 
\end{tabular}
\end{table}

\subsection{Why: Amdahl's law}
\evidence \Cref{fig:e02prof}: in the cProfile of unstable growth the map is
\NumMapFractionKTwoEight{} (k28), \NumMapFractionPThree{} (p3) and
\NumMapFractionKTen{} (k10) of the CPU time. The bounds are
\NumAmdahlBoundKTwoEight, \NumAmdahlBoundPThree{} and \NumAmdahlBoundKTen.
Refinement plus linked-list work is \NumProfRefPlusListKTwoEight{} (k28),
\NumProfRefPlusListPThree{} (p3) and \NumProfRefPlusListKTen{} (k10). In the
k28 profile (\NumProfTotalSKTwoEight\,s), \code{\_refine\_layer} alone takes
\NumProfRefineLayerShareKTwoEight{} of the self time. \code{Point.\_\_init\_\_}
runs \NumProfPointInitCallsKTwoEight{} times (\NumProfPointInitShareKTwoEight{}
cumulative), each paired with a \code{numpy.array} call
(\NumProfNpArrayCallsKTwoEight). The \code{min\_batch\_points} sweep
(\cref{fig:e02mbp}) is flat: the best setting (\NumMinBatchBest) is
\NumMinBatchBestGain{} faster than ``never'', which is inside the run-to-run
noise of \cref{sec:repro}.

\begin{figure}[htbp]
  \centering
  \includegraphics[width=0.92\textwidth]{\fig{e02_profile}}
  \caption{E2 cProfile breakdown of unstable growth (self time by category).
  The map (orange) is invisible at this scale; refinement and the
  linked list (Python objects, pointer walks) take
  \NumProfRefPlusListPThree--\NumProfRefPlusListKTwoEight{} of it.
  The Amdahl bound on the right is the most any map accelerator could give.}
  \label{fig:e02prof}
\end{figure}

\begin{figure}[htbp]
  \centering
  \includegraphics[width=0.6\textwidth]{\fig{e02_min_batch}}
  \caption{E2 \code{min\_batch\_points} sweep on k28 (12 unstable steps). The
  GPU session is flat across the threshold and matches the CPU median
  (dashed): the threshold has nothing to tune.}
  \label{fig:e02mbp}
\end{figure}

\reading The GPU path accelerates the only part of growth that was already
fast. Growth is a Python object graph: one \code{Point} per vertex, two
doubly linked lists, and refinement that walks them layer by layer. For the
GPU to pay, refinement and the linked list must become arrays (coordinates,
cdists, iterate indices), with the whole growth loop resident on the device.
Only then would $f_{\rm map}$ become the GPU-friendly bulk of the work.
Short of that, the measured ceiling is \NumAmdahlBoundMax.

\caveat cProfile inflates the per-call overhead of Python functions relative
to NumPy kernels, so the true $f_{\rm map}$ is probably a little larger than
measured. Even ten times the measured share would keep the bound close to 1.
Three repeats per configuration resolve differences of a few percent, not
less.

%======================================================================
\section{Accuracy versus cost (E4)}
\label{sec:accuracy}

\claim \code{area\_cutoff} trades accuracy for points along a predictable
power law. The default $10^{-4}$ is adequate for the period-1 cases but not
for p3.

\evidence \Cref{fig:e04pareto}(a) plots invariance error against pipeline
time. The invariance error is the distance from the image of a point inside
a polyline segment to the polyline of the image branch; the probe samples
segment interiors because a vertex's image is, by construction, often a
vertex. On k28 (10 steps), going from $10^{-4}$ to $10^{-7}$ moves the median
error from \NumInvPFiftyKTwoEightCutoffFour{} to
\NumInvPFiftyKTwoEightCutoffSeven{} and the 99th percentile from
\NumInvPNinetyNineKTwoEightCutoffFour{} to
\NumInvPNinetyNineKTwoEightCutoffSeven. The cost goes from
\NumTimeKTwoEightCutoffFour\,s to \NumTimeKTwoEightCutoffSeven\,s and from
\NumPointsKTwoEightCutoffFour{} to \NumPointsKTwoEightCutoffSeven{} points.
Over cutoffs $10^{-6}$--$10^{-9}$ the median error has slopes
\NumInvSlopeKTwoEightTen{} and \NumInvSlopeKTwoEightTwelve{} (k28, 10 and
12 steps), and \NumInvSlopePThreeThirteen{} and
\NumInvSlopePThreeSeventeen{} (p3, 13 and 17 steps). The point count has
slope \NumPointsCutoffExponentKTwoEight{} (k28) and
\NumPointsCutoffExponentPThree{} (p3) (\cref{fig:e04conv}(b)). Crossing
positions converge to the finest run with slopes \NumCrossDevSlopeKTwoEightTen,
\NumCrossDevSlopeKTwoEightTwelve, \NumCrossDevSlopePThreeThirteen{} and
\NumCrossDevSlopePThreeSeventeen{} (\cref{fig:e04conv}(a)). The full grid is
in \cref{tab:e04}.

\begin{figure}[htbp]
  \centering
  \includegraphics[width=0.88\textwidth]{\fig{e04_pareto}}
  \caption{E4. (a)~Accuracy against cost; labels are $\log_{10}$ of the
  cutoff. Every case traces a straight Pareto front, so there is no knee and
  no cheap win. (b)~Registered crossings per cutoff. k28 is stable at every
  cutoff; p3 needs $\le10^{-5}$ (13 steps) or $\le10^{-7}$ (17 steps) before
  the count settles. Counts are compared at fixed depth; more steps register
  more crossings (k28: \NumCrossingsKTwoEightTenRef{} at 10 steps,
  \NumCrossingsKTwoEightTwelveRef{} at 12), which is new tangle, not error.}
  \label{fig:e04pareto}
\end{figure}

\begin{figure}[htbp]
  \centering
  \includegraphics[width=0.88\textwidth]{\fig{e04_convergence}}
  \caption{E4 convergence. (a)~Distance of each finest-run crossing to the
  nearest crossing of the coarser run: steady power-law convergence; p3 at
  17 steps is the slowest. (b)~Cost in points: slopes \NumPointsCutoffExponentKTwoEight{} and
  \NumPointsCutoffExponentPThree, close to the predicted $-1/3$.}
  \label{fig:e04conv}
\end{figure}

\reading The exponents are what chord geometry predicts. A chord of length
$\ell$ on a curve of curvature $\kappa$ deviates by $h\approx\kappa\ell^2/8$
and encloses an area $\propto\kappa\ell^3$. A fixed area cutoff $A$ therefore
gives $\ell\propto A^{1/3}$, so points $\propto A^{-1/3}$ and a geometric
error $h\propto A^{2/3}$. The crossing \emph{count} is what matters for
topology. For p3 at 13 steps the default $10^{-4}$ registers
\NumCrossingsPThreeThirteenCutoffFour{} crossings and $10^{-3}$ registers
\NumCrossingsPThreeThirteenCutoffThree, while every cutoff from $10^{-5}$ on
gives \NumCrossingsPThreeThirteenRef. At 17 steps, $10^{-4}$ gives
\NumCrossingsPThreeSeventeenCutoffFour{} and $10^{-5}$ gives
\NumCrossingsPThreeSeventeenCutoffFive, against \NumCrossingsPThreeSeventeenRef{}
from $10^{-7}$ on. A coarse polygon both misses true crossings and creates
spurious ones. A count that depends on the cutoff at fixed depth is a
resolution error; a count that grows with depth or blasts is the tangle
growing. $10^{-7}$ is the cheapest setting that is converged in every case
tested.

\caveat The reference is the finest run ($10^{-9}$), not an exact manifold,
so the convergence plots measure self-consistency. The invariance error is a
reference-free proxy, but it measures one-step consistency, not distance to
the true manifold. p3 at 17 steps and $10^{-3}$ ran out of memory (16\,GB
cap): a coarse polygon of the escaping arm is not cheaper.

%======================================================================
\section{Physical invariants (E5)}
\label{sec:invariants}

\claim On every frozen case the numerics respect what the dynamics requires:
fixed points to machine precision, monotone cdists, manifolds mapped onto
themselves to within the refinement tolerance, and lobe areas preserved.

\evidence \Cref{tab:e05} summarises the cases. The fixed-point residual is
at most \NumFpResidualMax{} and the eigenvector residual at most
\NumEvecResidualMax. Along the manifolds, \NumCdistStepsManifolds{} cdist
steps (and \NumCdistStepsBridges{} along bridges) contain
\NumCdistDecreasingManifolds{} decreases and \NumCdistTiesManifolds{} ties.
No sampled image lands beyond the tail of its image branch
(\NumInvarianceBeyond{} such images). The largest invariance error,
\NumInvarianceMax, belongs to k10, one of the two cases (with inversion) grown at the default
$10^{-4}$ (\cref{fig:e05inv}). The \NumIterLinks{} registered iterate links
agree with the map to a median \NumIterLinkErrMedian{} (maximum
\NumIterLinkErrMax{}, again k10), so the error tracks each case's cutoff
(\cref{fig:e05lobes}(a)). The \NumLobePairs{} lobe/image-lobe pairs have
$|A'/A-1|\le\NumLobeRatioMax$ with \NumLobeSignFlips{} sign flips
(\cref{fig:e05lobes}(b)). The k28 lobes measure $\pm\NumLobeAreaKTwoEight$,
and their images agree to \NumLobeAreaAbsErrKTwoEight{} in absolute terms.
The p3 lobes are tiny (median \NumLobeAreaPThreeMedian), so their larger
\emph{relative} error corresponds to an absolute error of at most
\NumLobeAbsErrPThreeMax, a few times the cutoff.

\begin{table}[htbp]
\centering\footnotesize
\caption{E5 per-case invariants (maxima over orbit points and branches).
``arr.\ image'' counts the arrangement regions whose \code{image\_of} was
found/missing.}
\label{tab:e05}
\setlength{\tabcolsep}{3pt}
\resizebox{\textwidth}{!}{%
\begin{tabular}{llrrrrrccc}
\tablebody tables/e05_invariants.tex 
\end{tabular}}
\end{table}

\begin{figure}[htbp]
  \centering
  \includegraphics[width=0.88\textwidth]{\fig{e05_invariance}}
  \caption{E5 invariance error, one cumulative distribution per branch. The
  curves to the right are k10 and the inversion saddle, the two cases grown at
  the default cutoff $10^{-4}$; the $10^{-7}$ cases sit well to the left.}
  \label{fig:e05inv}
\end{figure}

\begin{figure}[htbp]
  \centering
  \includegraphics[width=0.88\textwidth]{\fig{e05_iterates_lobes}}
  \caption{E5. (a)~Error of each registered iterate link: the links of the
  default-cutoff cases (k10, inversion) sit highest, those of the $10^{-7}$
  cases lower. (b)~Relative area change of each lobe under one map step
  (filled) and of each arrangement region with a found image (open). All are
  $\le\NumLobeRatioMax$; the largest values belong to the tiny p3 lobes. (The
  legend overlaps a few points; it hides no outlier.)}
  \label{fig:e05lobes}
\end{figure}

\reading The area check is the strongest of these results. It compares
polygons built from independently refined arcs, so a systematic drift
would show up as a ratio that grows with depth, and none does.

\caveat (i) $||\lambda_u\lambda_s|-1|$ is \emph{weak} evidence: the
eigenvalues come from \code{np.linalg.eig} of the analytic Jacobian, whose
determinant is exactly 1, so the defect is
\NumDetDefectMaxAnalytic{} by construction. The exception confirms the
point: k10 passes no analytic Jacobian, its stored Jacobian deviates from
the analytic one by \NumJacDevKTen, and its defect is \NumDetDefectKTen.
(ii) The arrangement's \code{image\_of} found an image for
\NumArrImageFound{} regions and none for \NumArrImageMissing. It found them
for most of p3 and nested, for one region of each blasted k28 case, and for
none of k10, unblasted k28 or inversion (\cref{tab:e05}). The data does not
say whether this is a gap in \code{image\_of} or an expected consequence of
image crossings that are not registered. I list it as an open question
(\cref{sec:recs}), not as a bug.

%======================================================================
\section{Correctness of crossing detection (E9)}
\label{sec:intersections}

\claim For transverse crossings the R-tree path finds exactly the crossings
of the polygons, no more and no less.

\evidence Over \NumIsectConfigs{} configurations (\cref{tab:e09},
\cref{fig:e09}), an exact orientation-test brute force found
\NumIsectBruteTotal{} proper u$\times$s crossings. The registry matched
\NumIsectMatchedTotal{}, with \NumIsectMissedTotal{} missed (none needed the
near-parallel, zigzag or collision explanations) and \NumIsectExtraTotal{}
extra.

\begin{table}[htbp]
\centering\footnotesize
\caption{E9: crossings against brute force, summed per case. $^*$Some
configurations were too large for the u$\times$u check.}
\label{tab:e09}
\setlength{\tabcolsep}{4pt}
\begin{tabular}{lrrrrrrrrrr}
\tablebody tables/e09_intersections.tex 
\end{tabular}
\end{table}

\begin{figure}[htbp]
  \centering
  \includegraphics[width=0.88\textwidth]{\fig{e09_intersections}}
  \caption{E9. (a)~Detection time against segments. The library is
  linear; the vectorised brute force is steeper but faster on all but the
  largest configurations. (b)~Library against brute-force counts: every point
  lies on the diagonal.}
  \label{fig:e09}
\end{figure}

\reading Speed is not where the R-tree path shines at these sizes: its
slope is \NumIsectLibraryExponent{} against \NumIsectBruteExponent{} for the
brute force, but the brute force was faster in
\NumIsectBruteFasterConfigs{} of \NumIsectConfigs{} configurations. On the
largest (\NumIsectLargestSegments{} segments) the library took
\NumIsectLargestLibraryS\,s and the brute force \NumIsectLargestBruteS\,s
(the largest brute/library ratio was \NumIsectBruteOverLibraryMax). This is
not a problem: the stable side is short, so the brute force is effectively
linear here, and the library time includes registration and bookkeeping.
The asymptotic advantage remains with the library.

\caveat Two gaps. First, the polygons \emph{can} cross themselves:
\NumIsectUUTotal{} proper u$\times$u crossings appeared in
\NumIsectUUConfigs{} configurations, all at cutoff \NumIsectUUCutoff, and
\NumIsectUUSkipped{} large configurations were not checked. A manifold
cannot cross itself (CLAUDE.md), so these are discretisation artefacts. The
library never tests u$\times$u pairs, so nobody is told. A cheap self-crossing
check would make a good resolution diagnostic. Second, an API trap:
\code{compute\_intersections} (\code{reset=True}) \emph{replaces} the
\code{IntersectionRegistry} object, so a caller that cached
\code{workbench.intersection\_registry} beforehand holds a stale, empty
registry. This is documented in the docstring, but easy to trip over (the E6
script carries a comment about it).

%======================================================================
\section{Robustness across parameters (E6)}
\label{sec:robustness}

\claim One fixed recipe runs end to end over most of the area-preserving
$b=1$ family. It fails cleanly where it should, but at small $k$ it produces
tangles that are numerically fine and topologically empty.

\evidence \Cref{fig:e06out} and \cref{tab:e06}: \NumSweepComplete{} of
\NumSweepConfigs{} configurations (\NumSweepSuccessFraction) reached the
trellis, and \NumSaddleSuccessFraction{} of the $b=1$ saddles did. The
failure classes are:
\begin{itemize}[leftmargin=*]
\item \textbf{Small $k$:} \NumSweepMaxIterFails{} saddles with
  $k\in[\NumSweepMaxIterKMin,\NumSweepMaxIterKMax]$ grew the unstable side
  and then raised ``Max iterations reached'' in \code{grow\_until\_turnaround}
  on the stable side. The smallest completed $k$ is \NumSaddleMinKComplete.
\item \textbf{Elliptic point:} \NumSweepOtherElliptic{} configurations
  ($k\le\NumSweepOtherEllipticKMax$) were rejected as ``not a saddle'', by
  design.
\item \textbf{Inversion saddle:} \NumSweepOtherOom{} runs with
  $k\ge\NumSweepOtherOomKMin$ hit the 4\,GB cap during growth, the same
  blow-up as in \cref{sec:depth}.
\item \textbf{Orientation reversing $b=-1$:} all \NumSweepBMinusOneFails{}
  were rejected (``eigenvalues of disagreeing sign''), by design.
\item \textbf{Dissipative $b$:} \NumSweepDissipativeOk{} completed. Nothing
  in the pipeline checks area preservation, so completion here says nothing
  about correctness.
\end{itemize}

\begin{figure}[htbp]
  \centering
  \includegraphics[width=0.88\textwidth]{\fig{e06_sweep_outcomes}}
  \caption{E6. Last stage completed per configuration (height) and outcome
  (colour). Read the $b=1$ saddle panel left to right: a block of small-$k$
  exceptions, then complete runs. The other point fails as elliptic below
  $k=3$ and runs out of memory as the inversion saddle at large $k$.}
  \label{fig:e06out}
\end{figure}

\begin{table}[htbp]
\centering\footnotesize
\caption{E6: configurations per family by last stage completed.}
\label{tab:e06}
\setlength{\tabcolsep}{3.5pt}
\resizebox{\textwidth}{!}{%
\begin{tabular}{lrrrrrrrrrr}
\tablebody tables/e06_stages.tex 
\end{tabular}}
\end{table}

\reading ``Complete'' is not ``useful''. Of the \NumSweepSaddleOk{}
completed $b=1$ saddles, \NumSweepZeroCrossings{} have \emph{no} crossing
at all (up to $k=\NumSweepZeroCrossingsKMax$). The first crossing appears at
$k=\NumSweepFirstCrossingK$, and no configuration has more than
\NumSweepMaxCrossings{} (\cref{fig:e06saddle}). With a fixed 8-step recipe
the manifolds at small $k$ are simply too short to meet, because
$\lambda_u$ is small. A recipe has to scale its depth with
$1/\ln\lambda_u$, as E7's tangency study does.

\begin{figure}[htbp]
  \centering
  \includegraphics[width=0.88\textwidth]{\fig{e06_saddle_sweep}}
  \caption{E6 $b=1$ saddle. (a)~Crossings and bridges found by the fixed
  recipe: none below $k=\NumSweepFirstCrossingK$ and intermittent up to
  $k=\NumSweepZeroCrossingsKMax$. (b)~Stage times across $k$: cheap
  everywhere; cost is not what limits this sweep.}
  \label{fig:e06saddle}
\end{figure}

\caveat The ``Max iterations'' failures are a limit of the stable-side stop
condition with this recipe, not necessarily of the map. A larger cap or a
different stop condition might rescue them; I did not try either. Note also
that the 4\,GB cap of E6 is lower than elsewhere.

%======================================================================
\section{Limits and deep exploration (E7)}
\label{sec:breaking}

Escape and tangency probe failure modes of the numerics; depth and blasts
probe how far the pipeline carries a growing tangle, whose topology is
expected to change along the way (\cref{sec:setup}).

\subsection{Escape to infinity}
\label{sec:escape}
\claim The documented escape reproduces, and the documented mitigation
works.

\evidence \Cref{fig:e07escape}: on the $k=2.1$ period-3 orbit with a fixed
step count, the run completed \NumEscapeStepsDone{} steps holding
\NumEscapePeakPoints{} points, with $\max|x|=\NumEscapeMaxCoord$. It was then
killed at the 20\,GB cap (peak \NumEscapePeakRssMb\,MB) after
\NumEscapeKilledStepS\,s inside the next step. The same arm grown with a stop
condition (arclength 0.5, or $\max|x|>2.5$) stopped safely at
\NumEscapeStopArclengthPoints{} points; arclength 3 stopped at
\NumEscapeStopArclengthThreePoints.

\begin{figure}[htbp]
  \centering
  \includegraphics[width=0.88\textwidth]{\fig{e07_escape}}
  \caption{E7 escape ($k=2.1$, period 3). (a)~Points per step: the fixed
  count dies one step after the curve turns up; the stop conditions
  (open markers) end before the take-off. (b)~The largest coordinate jumps
  in the last completed step, which is what the $\max|x|$ stop condition
  catches.}
  \label{fig:e07escape}
\end{figure}

\reading \code{grow\_n\_times} has no notion of budget. A single step can go
from harmless to tens of gigabytes, so stop-condition growth should be the
recommended interface.

\subsection{Depth}
\label{sec:depth}
\claim p3 is numerically sound through \NumDepthNumericsMaxPThree{} unstable
steps, and each extra step uncovers more of the tangle. From
\NumDepthUnreliableMinPThree{} steps on, the symbolic layer flags its result
as provisional. The inversion saddle hits two genuine errors: a topology
invariant assertion when shallow and memory exhaustion when deep.

\evidence \Cref{fig:e07depth}: p3 completes the numerics through
\NumDepthNumericsMaxPThree{} steps (\NumPThreePointsSeventeen{} points);
\NumDepthOomPThree{} steps hit the 20\,GB cap. The topology stack completes
at every depth that finishes the numerics. The tangle it finds grows with
depth: the bridge-class count is constant through
\NumDepthReliableMaxPThree{} steps and then rises with every extra step.
From \NumDepthUnreliableMinPThree{} steps on, unresolved and virtual classes
appear (\code{is\_reliable=False}): some classes cannot be resolved yet
(an image not registered, or a dual-graph walk that does not complete) and
need more growth or blasting. It does not mean the numerics are wrong. The
inversion saddle completes the numerics only up to
\NumDepthNumericsMaxInversion{} steps. At \NumInversionOomSteps{} steps it is
killed at the cap after reaching \NumInversionPointsSixSteps{} points at
step \NumDepthNumericsMaxInversion. So the documented ``7 steps does not finish'' is a memory blow-up,
not a hang. At \NumInversionAssertSteps{} steps the numerics succeed, but the
topology stack trips its own invariant: ``crossing 0 at stable cdist 0.0
\ldots{} is owned by 0 partition elements''. This is a genuine error (an
\code{AssertionError}), unlike the provisional flag at
\NumDepthNumericsMaxInversion{} steps, where the stack completes with a
virtual class.

\begin{figure}[htbp]
  \centering
  \includegraphics[width=0.88\textwidth]{\fig{e07_depth}}
  \caption{E7 depth. (a)~Pipeline time against unstable steps; p3 climbs steeply
  at the deepest steps. (b)~Outcome by layer. The numerics fail only by memory,
  and the inversion topology raises an assertion at shallow depth. Open
  circles are topology runs that completed with a provisional symbolic
  result (\code{is\_reliable=False}: some classes still need more growth).
  They are not failures.}
  \label{fig:e07depth}
\end{figure}

\caveat The inversion assertion lives in the topology layer and is recorded
here because the stress test reached it. Whether the shallow inversion
trellis is too thin to partition or the assertion is too strict is for the
topology owner to decide. Depth itself is not a fault: a deeper trellis
has more crossings and classes by design.

\subsection{Blasts}
\label{sec:blasts}
\evidence \Cref{fig:e07blasts}: k28 and p3 ran
\NumMaxBlastsKTwoEight{} and \NumMaxBlastsPThree{} blasts and nested ran
\NumMaxBlastsNested{}, all without numerical error. k28's crossings grow
from \NumBlastCrossingsStartKTwoEight{} to \NumBlastCrossingsEndKTwoEight, and
its last blast took \NumBlastLastSKTwoEight\,s. Each blast uncovers new
structure, which is the purpose of blasting. After the last blast, k28 has
\NumBlastClassesKTwoEight{} bridge classes, \NumBlastUnresolvedKTwoEight{} of
them still unresolved and \NumBlastVirtualKTwoEight{} virtual, so the
symbolic layer marks its result provisional. p3 has
\NumBlastClassesPThree{} classes (\NumBlastUnresolvedPThree{} unresolved,
\NumBlastVirtualPThree{} virtual). Nested has \NumBlastClassesNested{}
classes, all resolved. The topology stack completed in all three cases. No
exception, assertion or memory failure occurred.

\begin{figure}[htbp]
  \centering
  \includegraphics[width=0.88\textwidth]{\fig{e07_blasts}}
  \caption{E7 blasts. (a)~Time per blast grows geometrically for k28 and
  nested and stays at milliseconds for p3. (b,c)~Crossings and bridges grow
  geometrically. No numerical error within \NumMaxBlastsNested--\NumMaxBlastsKTwoEight{} blasts; the growing counts
  are the new tangle structure each blast reveals.}
  \label{fig:e07blasts}
\end{figure}

\reading Blasting is numerically sound to the depth tested and does what it
is for: it reveals new crossings, bridges and classes. Unresolved and
virtual classes after many blasts are the symbolic layer saying that some
classes cannot be resolved yet (E11 records why, \cref{sec:blastdepth-provisional}). They are open questions for more
growth or blasting, not errors.

\caveat The topology was probed only on the final state, so this
experiment does not show how the classes and words changed blast by blast
(see \cref{sec:blastdepth}).

\subsection{Near tangency: artefact crossings}
\label{sec:tangency}
\claim This is the most important limitation found. As $k\to-1$ (the
saddle-centre bifurcation) the true separatrix splitting becomes
exponentially small. The numerics then register polygon-on-polygon crossings
as if they were transverse, their number at fixed $k$ and fixed manifold
length is set by the resolution, and no warning fires.

\evidence \Cref{fig:e07tang}: \NumTangencyConfigs{} configurations, with $k+1$
log-spaced from $10^{-4}$ to 2.5 and three cutoffs. At $k+1=10^{-4}$ the
registered crossings are \NumTangencyCrossingsNearCutoffSeven,
\NumTangencyCrossingsNearCutoffTen{} and \NumTangencyCrossingsNearCutoffTwelve{}
for cutoffs $10^{-7}$, $10^{-10}$ and $10^{-12}$. The peaks over $k$ are
\NumTangencyCrossingsPeakCutoffSeven, \NumTangencyCrossingsPeakCutoffTen{}
and \NumTangencyCrossingsPeakCutoffTwelve. From $k+1\ge\NumTangencyAgreeKPlusOne$
on, all three cutoffs agree (e.g.\ \NumTangencyCrossingsFar{} crossings at
$k+1=1$). At the finest cutoff and $k+1\le0.025$, the smallest $|\sin|$ per
configuration lies between \NumTangencyMinSinFineLo{} and
\NumTangencyMinSinFineHi{} (median of the medians \NumTangencyMedianSinFine).
The heuristic splitting $e^{-\pi^2/\ln\lambda}$ over the same range is
\NumTangencyRefNear{} to \NumTangencyRefMid. Across all of this, warnings of
any kind: \NumTangencyWarningsTotal. The smallest angle anywhere,
\NumTangencyMinSin, sits at the transition, not near $k=-1$.

\begin{figure}[htbp]
  \centering
  \includegraphics[width=0.88\textwidth]{\fig{e07_tangency}}
  \caption{E7 tangency. (a)~Crossings against $k+1$: below $k+1\approx\NumTangencyAgreeKPlusOne$
  the three cutoffs disagree strongly. That is the
  signature of artefacts, since a true crossing cannot depend on the
  polygon's resolution. (b)~Smallest crossing angle against the heuristic
  splitting (dashed): the registered angles are many decades larger than the
  true splitting, i.e.\ they measure polygon noise.}
  \label{fig:e07tang}
\end{figure}

\reading When the true separation of the manifolds is far below the
polygon's own geometric error (\cref{sec:accuracy}), the two polygons
interleave at random. Each interleaving is a clean, well-conditioned segment
crossing (median $|\sin|$ \NumTangencyMedianSinFine), far from the library's relative near-parallel
threshold of $10^{-12}$, so it is registered as transverse. The crossing
tests are correct about the \emph{polygons}; the polygons are not accurate
enough to resolve the \emph{manifolds}.

\caveat The splitting reference is a heuristic order of magnitude, not a
computed splitting. Some of the near-tangent crossings may be real; the
claim is that their \emph{count} is not trustworthy. A guard that compares
each crossing's angle and position with the local invariance error would
flag this regime.

%======================================================================
\section{Deep blasting: how the topology evolves (E11)}
\label{sec:blastdepth}

\paragraph{Why this experiment.} A blast iterates the bridges inside a
resonance zone, which registers crossings that the original growth did not
reach. Its purpose is to change the topology: new crossings, bridges and
holes, split elements, new bridge classes, new letters, new words and new
transition matrices. In this section topology change is the quantity being
measured, not a fault (\cref{sec:setup}). E7 blasted three zones and probed
the topology once, after the last blast (\cref{sec:blasts}). E11 blasts five
zones up to \NumDeepMaxBlasts{} times and snapshots the whole topology stack
(trellis, partitions, minimal trellis, dual graph, symbolic dynamics) after
\emph{every} blast. It asks five questions: how the symbolic description
evolves, what each blast costs, which parameters change what is uncovered,
how the library's provisional flags behave as the tangle grows, and whether
any genuine error appears.

\subsection{Design of the sweep}
\label{sec:blastdepth-design}

\paragraph{Cases.} k28 and p3 as in E7; k32, the $k=3.2$ saddle built by the
k28 recipe, as a second single saddle; nested\_outer, the period-1 zone of
the nested map; and nested\_inner, its period-3 zone, blasted after the outer
zone has been blasted twice (the \code{build\_nested} recipe). Every case is
built at \code{area\_cutoff} $10^{-7}$, so all cells of one case start from
the same tangle; the swept \code{area\_cutoff} is the refinement threshold
used \emph{while blasting}.

\paragraph{Grid.} The \emph{reference cell} of a case is cutoff $10^{-7}$, the
recipe's \code{min\_separation} ($10^{-5}$, or $10^{-4}$ for the nested cases)
and one iteration per blast. Around it one parameter moves at a time:
\code{min\_separation} $\in\{\text{none},10^{-6},10^{-5},10^{-4},10^{-3}\}$
(plus $10^{-2}$ for k28 and k32), \code{area\_cutoff}
$\in\{10^{-5},10^{-6},10^{-8}\}$, \code{num\_iterations} $=2$, and two repeats
of the reference cell for k28, p3 and nested\_inner. That makes
\NumDeepConfigs{} configurations, each in a fresh process (900\,s timeout,
12\,GB RSS cap, 8 workers), which took \NumDeepWallMin\,min of wall-clock.
The full cutoff $\times$ separation product was dropped for time
(\cref{sec:threats}).

\paragraph{Stopping.} A configuration ends after \NumDeepMaxBlasts{} blasts
(\NumDeepStopMaxBlasts{} configs) or when its own guard declines the next
blast: the time budget (timeout minus 120\,s, \NumDeepStopBudget{} configs)
or the memory budget ($0.9\times$ the RSS cap, predicted from the last
blast's growth, \NumDeepStopMemoryBudget{} configs). These budget stops are
the experiment's resource guard. They keep the last snapshot intact and are
not failures; all \NumDeepConfigsOk/\NumDeepConfigs{} configurations ended
\code{ok} at the harness level. Topology errors raised inside them are
recorded per blast (\cref{sec:blastdepth-errors}).

\paragraph{What is recorded.} After each blast: the numerics counts
(crossings, bridges, points, time, RSS), the two law checks (same-stability
crossings; decreasing cdist along a new bridge), a lobe-area probe, the
topology counts (classes, active, inert, holes, elements, unresolved,
virtual, ambiguous, \code{is\_reliable}) and the topology outcome. A
topology exception or \code{AssertionError} is recorded for that blast and
the blasting goes on; such a snapshot has no counts, and the next one is
compared with the last snapshot that completed.

\paragraph{Canonical words.} The session alphabet never reuses a letter, so
a class whose elements a blast has renumbered returns under a new letter
even when its dynamics did not change. To tell a relabelling from new
symbolic structure, E11 also records \emph{canonical words}: the unrefined
words with every letter renamed by its class's position in the class table
(active classes $A_0, A_1,\ldots$, inert ones $I_0, I_1,\ldots$). A snapshot
is \emph{new structure} when its canonical words differ from the previous
completed snapshot, a \emph{relabelling} when change flags are set but the
canonical words are equal, and \emph{no change} otherwise.

\subsection{How the topology evolves}
\label{sec:blastdepth-evolution}

\claim Blasting keeps uncovering new symbolic structure. On k28 and the
period-3 zones it arrives in bursts; k32 changes its words at every
completed blast. Nearly every change is a change of the words themselves,
not a relabelling.

\evidence \Cref{tab:e11summary,fig:e11evolution}. Over all configurations,
\NumDeepStructureSteps{} post-blast snapshots changed the canonical words,
only \NumDeepRelabelSteps{} were pure relabellings, and
\NumDeepNoChangeSteps{} changed nothing.
\begin{itemize}[leftmargin=*]
\item \emph{k28} goes from \NumDeepClassesStartKTwoEight{} to
  \NumDeepClassesEndKTwoEight{} classes (\NumDeepActiveEndKTwoEight{} active,
  \NumDeepInertEndKTwoEight{} inert) and from
  \NumDeepCrossingsStartKTwoEight{} to \NumDeepCrossingsEndKTwoEight{}
  crossings. Its largest step is at blast \NumDeepJumpBlastKTwoEight, from
  \NumDeepJumpBeforeKTwoEight{} to \NumDeepJumpAfterKTwoEight{} classes. The
  first new words appear at blast \NumDeepFirstStructureKTwoEight.
\item \emph{k32} changes its canonical words at
  \NumDeepStructureChangesKThreeTwo{} completed blasts and reaches
  \NumDeepClassesEndKThreeTwo{} classes at blast
  \NumDeepClassesEndBlastKThreeTwo{} (its last completed snapshot), after a
  step from \NumDeepJumpBeforeKThreeTwo{} to \NumDeepJumpAfterKThreeTwo{} at
  blast \NumDeepJumpBlastKThreeTwo.
\item \emph{p3 and nested\_inner} do not change at all until their first
  new words (p3: blast \NumDeepFirstStructurePThree; nested\_inner: blast
  \NumDeepFirstStructureNestedInner). There the
  classes jump from \NumDeepJumpBeforePThree{} to \NumDeepJumpAfterPThree{}
  and from \NumDeepJumpBeforeNestedInner{} to \NumDeepJumpAfterNestedInner,
  and then grow by exactly \NumDeepClassStepPThree{} classes per blast to
  \NumDeepClassesEndPThree{} and \NumDeepClassesEndNestedInner. From the
  burst on, every blast changes the words (\NumDeepStructureChangesPThree{}
  and \NumDeepStructureChangesNestedInner{} new-structure blasts).
\item \emph{nested\_outer} rises from \NumDeepClassesStartNestedOuter{} to
  \NumDeepJumpAfterNestedOuter{} classes at blast
  \NumDeepJumpBlastNestedOuter, then falls from
  \NumDeepDropBeforeNestedOuter{} to \NumDeepDropAfterNestedOuter{} at blast
  \NumDeepDropBlastNestedOuter, while its crossings grow from
  \NumDeepCrossingsStartNestedOuter{} to \NumDeepCrossingsEndNestedOuter.
\end{itemize}

\begin{table}[htbp]
\centering
\caption{E11 per case at the reference cell. ``first new words'' is the first
blast whose canonical words changed; ``new-word blasts'' counts them;
``prov.'' marks a provisional symbolic result (\code{is\_reliable=False}).
The last two columns cover every configuration of the case.}
\label{tab:e11summary}
\scriptsize\setlength{\tabcolsep}{2.5pt}
\begin{adjustbox}{max width=\textwidth}
\begin{tabular}{lllllrlllr}
\tablebody tables/e11_summary.tex 
\end{tabular}
\end{adjustbox}
\end{table}

\begin{figure}[htbp]
  \centering
  \includegraphics[width=\textwidth]{\fig{e11_evolution}}
  \caption{E11 reference cells. (a)~Active and (b)~inert bridge classes, and
  (c)~registered crossings (log scale), against blasts done. Filled markers:
  the canonical words changed at that blast (new structure). Red $\times$:
  the topology raised at that blast, so there is no snapshot. The class
  counts move in bursts for k28 (blast \NumDeepJumpBlastKTwoEight) and for
  p3 and nested\_inner (blast \NumDeepJumpBlastPThree); k32's words change
  at every completed blast, with its largest class step at blast \NumDeepJumpBlastKThreeTwo; the crossings
  grow geometrically for the period-1 zones and slowly for the period-3
  ones.}
  \label{fig:e11evolution}
\end{figure}

\reading This is blasting doing its job. Each burst is a blast after which
the class table is rebuilt over a stable manifold with new holes and cuts,
so the class count is recomputed, not appended to; this is also why
nested\_outer can \emph{lose} classes. Why the bursts fall at those
particular blasts (presumably the first blast whose images reach a new part
of the stable manifold) I have not checked. The period-3 zones show a
regular pattern after their burst, a fixed number of new classes per blast.
The period-1 zones have no such pattern, and their crossings grow
geometrically rather than slowly.

\caveat Class counts at a given blast depend on the blasting cutoff
(\cref{sec:blastdepth-params}), so these are counts at cutoff $10^{-7}$.
Snapshots that raised leave gaps: k28 and k32 have none between blasts
\NumDeepErrorBlastsKTwoEight, which is why k32's last count is at blast
\NumDeepClassesEndBlastKThreeTwo{} and k28's jump to
\NumDeepClassesEndKTwoEight{} is not resolved blast by blast.

\subsection{Cost of blasting: the practical depth limit}
\label{sec:blastdepth-cost}

\claim For the large zones (k28, k32, nested\_outer) every blast costs a
roughly constant factor more than the one before. At cutoff $10^{-7}$ memory
sets how deep one can blast; at coarser cutoffs the time budget binds first.
The period-3 zones cost almost nothing.

\evidence \Cref{fig:e11cost}. Per blast, as geometric means over the last
five blasts, k28 multiplies its crossings by \NumDeepCrossingGrowthKTwoEight,
its bridge points by \NumDeepPointGrowthKTwoEight{} and its blast time by
\NumDeepTimeGrowthKTwoEight; k32 by \NumDeepCrossingGrowthKThreeTwo,
\NumDeepPointGrowthKThreeTwo{} and \NumDeepTimeGrowthKThreeTwo; nested\_outer
by \NumDeepCrossingGrowthNestedOuter, \NumDeepPointGrowthNestedOuter{} and
\NumDeepTimeGrowthNestedOuter. Their last blasts took
\NumDeepLastBlastSKTwoEight, \NumDeepLastBlastSKThreeTwo{} and
\NumDeepLastBlastSNestedOuter\,s and peaked at
\NumDeepPeakRssGbKTwoEight, \NumDeepPeakRssGbKThreeTwo{} and
\NumDeepPeakRssGbNestedOuter\,GB, and the memory budget stopped them after
\NumDeepBlastsKTwoEight, \NumDeepBlastsKThreeTwo{} and
\NumDeepBlastsNestedOuter{} blasts. Of the \NumDeepRefCutoffCells{}
large-zone cells at cutoff $10^{-7}$ with one iteration per blast,
\NumDeepRefCutoffMemoryStops{} were stopped by the memory budget and
\NumDeepRefCutoffTimeStops{} by the time budget at the same blast count, so
which guard binds first there depends on machine load. At the coarser
blasting cutoffs $10^{-5}$ and $10^{-6}$ the time budget stopped all
\NumDeepCoarseTimeStops/\NumDeepCoarseCells{} large-zone cells, after
\NumDeepCoarseBlasts{} blasts at \NumDeepCoarseRss\,GB. p3 and nested\_inner
grow their crossings by only \NumDeepCrossingGrowthPThree{} and
\NumDeepCrossingGrowthNestedInner{} per blast, finish all
\NumDeepMaxBlasts{} blasts at a mean of \NumDeepMeanBlastSPThree{} and
\NumDeepMeanBlastSNestedInner\,s per blast, and peak at
\NumDeepProcRssGbPThree{} and \NumDeepProcRssGbNestedInner\,GB.

\begin{figure}[htbp]
  \centering
  \includegraphics[width=\textwidth]{\fig{e11_cost}}
  \caption{E11 cost at the reference cells. (a)~Points on all bridges and
  (b)~time per blast (log scales): straight lines for the period-1 zones,
  rising only slowly for the period-3 zones (p3's bridge points go from
  \NumDeepPointsStartPThree{} to \NumDeepPointsEndPThree). (c)~Peak RSS against the memory budget
  (dashed). Squares mark configurations that their own budget stopped
  because the next blast would not fit. That is the experiment's guard,
  not a failure.}
  \label{fig:e11cost}
\end{figure}

\reading Each blast iterates every interior bridge, and the bridges of a
large zone are stretched by the saddle at each step, so points, crossings
and time all grow by about the same factor. Because each further blast
needs that factor more memory, more hardware buys only a blast or two. This
is the same geometric memory growth as unbudgeted growth (\cref{sec:escape});
here the harness's budget stopped it before any OOM. The same cure applies: predict the next blast's cost from the last factor, as
the harness does here, and refuse or warn above a budget. In the period-3
zones the bridge points grow only by \NumDeepPointGrowthPThree{} (p3) and
\NumDeepPointGrowthNestedInner{} (nested\_inner) per blast, so deep
blasting is cheap.

\caveat The growth factors are fitted on five blasts. p3's time factor
(\NumDeepTimeGrowthPThree) is below one because its blasts take milliseconds
and the timing is noise.

\subsection{Dependence on parameters, and determinism}
\label{sec:blastdepth-params}

\claim \code{min\_separation} has no effect at any value tested, and the
runs are deterministic. The blasting \code{area\_cutoff} changes what the
blasts uncover. Two iterations per blast give the same canonical words as
two single blasts wherever both completed, in every case but one
nested\_outer pair.

\evidence \Cref{fig:e11events,fig:e11parameters,fig:e11agreement}.
\emph{Separation and repeats:} every separation cell matches its reference
exactly at every blast, in crossings, bridge points, topology outcome and
the exact topology hash: \NumDeepSeparationIdenticalKTwoEight{} (k28),
\NumDeepSeparationIdenticalKThreeTwo{} (k32),
\NumDeepSeparationIdenticalPThree{} (p3),
\NumDeepSeparationIdenticalNestedOuter{} (nested\_outer) and
\NumDeepSeparationIdenticalNestedInner{} (nested\_inner). So do the repeats
(\NumDeepRepeatIdenticalKTwoEight, \NumDeepRepeatIdenticalPThree,
\NumDeepRepeatIdenticalNestedInner).
\emph{Cutoff:} words are compared only at blasts where both snapshots
completed; a blast where one cell raised and the other did not is an
outcome difference, counted apart. In \NumDeepCutoffWordsFirst{} of the
\NumDeepCutoffCells{} cutoff cells the canonical words leave the reference
before the crossing counts do. At $10^{-5}$ p3's words differ from blast
\NumDeepCutoffFiveWordsPThree{} and its crossings from blast
\NumDeepCutoffFiveCrossingsPThree; k28's from blasts
\NumDeepCutoffFiveWordsKTwoEight{} and \NumDeepCutoffFiveCrossingsKTwoEight{}
(its topology outcome first differs at blast
\NumDeepCutoffFiveOutcomeKTwoEight, where the reference raises). The three
exceptions: k28 at $10^{-6}$ (crossings from blast
\NumDeepCutoffSixCrossingsKTwoEight, words from blast
\NumDeepCutoffSixWordsKTwoEight), k28 at $10^{-8}$ (crossings from blast
\NumDeepCutoffEightCrossingsKTwoEight, words \NumDeepCutoffEightWordsKTwoEight)
and nested\_outer at $10^{-8}$ (crossings from blast
\NumDeepCutoffEightCrossingsNestedOuter, words
\NumDeepCutoffEightWordsNestedOuter). After
\NumDeepBlastsPThree{} blasts, p3 holds \NumDeepCutoffFiveClassesPThree,
\NumDeepCutoffSixClassesPThree, \NumDeepClassesEndPThree{} and
\NumDeepCutoffEightClassesPThree{} classes at cutoffs $10^{-5}$, $10^{-6}$,
$10^{-7}$ and $10^{-8}$.
\emph{Two iterations:} for $n\ge1$ (blast 0 is the shared pre-blast build),
blast $n$ with two iterations has the same canonical
words as single blast $2n$ in \NumDeepDoubleWordsSameKTwoEight{} (k28),
\NumDeepDoubleWordsSameKThreeTwo{} (k32), \NumDeepDoubleWordsSamePThree{}
(p3) and \NumDeepDoubleWordsSameNestedInner{} (nested\_inner) of the pairs
where both snapshots completed, and \NumDeepDoubleWordsSameNestedOuter{}
for nested\_outer (they differ at $n=\NumDeepDoubleWordsDifferNestedOuter$). The crossing counts agree only up to
$n=\NumDeepDoubleCrossingsAgreeKTwoEight$ (k28),
\NumDeepDoubleCrossingsAgreeKThreeTwo{} (k32),
\NumDeepDoubleCrossingsAgreeNestedOuter{} (nested\_outer) and
\NumDeepDoubleCrossingsAgreeNestedInner{} (nested\_inner); for p3 they agree
up to $n=\NumDeepDoubleCrossingsAgreePThree$, every pair compared.

\begin{figure}[htbp]
  \centering
  \includegraphics[width=0.95\textwidth]{\fig{e11_events}}
  \caption{E11: what each blast changed, one row per configuration, grouped
  by case. Dark blue: new canonical words (new structure). Mid-blue: change
  flags set but canonical words unchanged (relabelling). Grey: no change.
  Red: the topology raised (genuine error; $\times$ marks the one
  exception, the rest are assertions). White: not run, because the
  configuration's budget stopped it. A white dot marks a provisional
  symbolic result (\code{is\_reliable=False}). The separation and repeat
  rows equal their reference row; the cutoff rows do not.}
  \label{fig:e11events}
\end{figure}

\begin{figure}[htbp]
  \centering
  \includegraphics[width=\textwidth]{\fig{e11_parameters}}
  \caption{E11 parameters. (a)~First blast at which the canonical words /
  the crossing counts differ from the reference cell, words compared only
  where both snapshots completed (``='': identical at every blast; ``--'':
  never). (b,c)~Bridge classes against blasts for each
  blasting cutoff, k28 and p3; red $\times$: snapshots where the topology
  raised (any cutoff).}
  \label{fig:e11parameters}
\end{figure}

\begin{figure}[htbp]
  \centering
  \includegraphics[width=\textwidth]{\fig{e11_agreement}}
  \caption{E11 agreement. (a)~Blast $n$ with two iterations per blast
  against single-iteration blast $2n$: crossing counts and canonical words.
  (b)~Separation and repeat cells identical to the reference at every
  blast.}
  \label{fig:e11agreement}
\end{figure}

\reading In these zones \code{min\_separation}, even at $10^{-2}$,
evidently never removes a blast child, so it is not a knob that matters
here, and the whole pipeline (blast, repin, topology) is reproducible to the
hash. The cutoff does matter: blasted bridges are refined at that cutoff, so
a different cutoff can move a crossing from one element to another (new
words, same count) before it adds or removes one. As in E4 and the tangency
study (\cref{sec:accuracy,sec:tangency}), a count that depends on the
cutoff at a \emph{fixed} blast reflects resolution, while growth with
blasts at a fixed cutoff is discovery. On the two-iteration side, the
numbers of crossings drift apart after a few blasts, probably because the
two routes refine and repin at different moments, but the symbolic words
agree wherever both completed, except in one nested\_outer pair. Along this
route the words are more robust than the raw counts.

\caveat No cutoff is a ground truth here: the finer $10^{-8}$ runs stopped
earlier, so the cutoff spread at deep blasts is not settled. The
nested\_outer two-iteration pair whose words differ has not been examined.

\subsection{Provisional symbolic results as the tangle grows}
\label{sec:blastdepth-provisional}

\claim The library's provisional flags come and go as blasting proceeds. A
blast can resolve them, and a burst of new classes opens new ones. They are
open questions that call for more growth or blasting, not errors.

\evidence \Cref{fig:e11provisional}. k28, k32 and nested\_outer start
provisional (an unresolved class before any blast), and the first blast
clears it. The result turns provisional again at blast
\NumDeepFirstProvisionalKTwoEight{} (k28), \NumDeepFirstProvisionalKThreeTwo{}
(k32), \NumDeepFirstProvisionalPThree{} (p3) and
\NumDeepFirstProvisionalNestedInner{} (nested\_inner): at or before each
case's class burst. At the end, k28 has \NumDeepUnresolvedEndKTwoEight{}
unresolved, \NumDeepVirtualEndKTwoEight{} virtual and
\NumDeepAmbiguousEndKTwoEight{} ambiguous classes; k32
\NumDeepUnresolvedEndKThreeTwo, \NumDeepVirtualEndKThreeTwo{} and
\NumDeepAmbiguousEndKThreeTwo; p3 \NumDeepUnresolvedEndPThree,
\NumDeepVirtualEndPThree{} and \NumDeepAmbiguousEndPThree; nested\_inner
\NumDeepUnresolvedEndNestedInner, \NumDeepVirtualEndNestedInner{} and
\NumDeepAmbiguousEndNestedInner. nested\_outer ends with
\NumDeepUnresolvedEndNestedOuter{} of each and a reliable result. In p3 and
nested\_inner the counts stay flat from the burst to the last blast.

\begin{figure}[htbp]
  \centering
  \includegraphics[width=\textwidth]{\fig{e11_provisional}}
  \caption{E11 provisional flags at the reference cells: (a)~unresolved
  classes, (b)~virtual classes (\code{new1}, \ldots), (c)~ambiguous walks
  against blasts. Open markers: provisional symbolic result
  (\code{is\_reliable=False}); red $\times$: no snapshot (topology error).
  Cases are offset vertically by $\pm0.16$ so that they do not overlap.}
  \label{fig:e11provisional}
\end{figure}

\reading The library records why each class is unresolved, and the
reasons differ before and after the bursts. Before any blast, the
\NumDeepUnresolvedSingletonStart{} unresolved classes of the reference cells
are all ``a landing is a singleton element and no member has a registered
image chain; grow or blast'', and the first blast registers those images,
which is exactly the library's ``grow or blast until the images are
registered''. After the bursts the reason is mostly different: of the
unresolved entries in the reference snapshots after blast 0,
\NumDeepUnresolvedUnreachable{} are ``dual-graph walk unreachable'' and only
\NumDeepUnresolvedSingleton{} the singleton reason (both in k28's last
snapshot, next to \NumDeepUnresolvedUnreachableEndKTwoEight{} unreachable
walks). In p3 and nested\_inner the one unresolved class is an unreachable
walk at every blast after the burst. So the burst-era flags are mostly not
missing images but walks the dual graph cannot complete. That is still a
provisional result for the topology owner (an open question that more
growth, or a further rule, may settle), not a numerics fault.

\caveat The flags are counted per snapshot; I did not check whether the
unresolved class is the same class from blast to blast.

\subsection{Genuine errors}
\label{sec:blastdepth-errors}

\claim No same-stability crossing and no decreasing cdist appeared at any
blast. The area check is inconclusive in the two-iteration cells (see
below). The topology
layer raised in \NumDeepErrorSnapshots{} of \NumDeepSnapshots{} snapshots,
in \NumDeepErrorConfigs{} of \NumDeepConfigs{} configurations, with
\NumDeepErrorClasses{} distinct errors. These are genuine errors for the
topology owner to triage. The blasting itself never failed.

\evidence Over \NumDeepBlasts{} blasts there were
\NumDeepSameStability{} same-stability crossings and
\NumDeepCdistDecreasing{} decreasing cdist steps
(\NumDeepLawViolations{} crossing/cdist law violations). The lobe-area probe
compared \NumDeepLobeCompared{} lobe pairs; the median of the per-snapshot
median $|A'/A-1|$ is \NumDeepLobeRatioMedian{} (per-pair ratios are not
stored), and the reference cells peak at \NumDeepLobeRatioMaxReference.
In the two-iteration cells (\NumDeepLobeRatioMaxWhere) $|A'/A-1|$ reaches
\NumDeepLobeRatioMax{} with \NumDeepLobeSignFlips{} sign flips. The probe
compares polygons, which a fold can make self-intersecting, so this is
either a self-intersecting probe polygon or a real area defect; I have not
decided which, and the area check stays open there.
\Cref{tab:e11errors} lists the topology errors
(\NumDeepErrorAsserts{} assertions, \NumDeepErrorExceptions{} exception):
\begin{enumerate}[leftmargin=*]
\item \textbf{Zero-width iterated element ``not inside its parent''}
  (\code{PartitionFamily.\allowbreak from\_\allowbreak minimal}, assertion):
  \NumDeepErrZeroWidthSnapshots{} snapshots in
  \NumDeepErrZeroWidthConfigs{} configurations, blasts
  \NumDeepErrZeroWidthBlasts. It is raised at blasts \NumDeepErrorBlastsKTwoEight{} in every k28
  and k32 cell at cutoff $10^{-7}$ (all separations and repeats), earlier in blast count (though later in iterations) with two iterations per blast, and
  in k32 at $10^{-8}$. On k28 the snapshot at blast
  \NumDeepClassesEndBlastKTwoEight{} completes again, and k28 at the other
  cutoffs never raises it.
\item \textbf{Holes disagree on \code{bridge\_side}}
  (\code{check\_holes\_\allowbreak share\_\allowbreak bridge\_side}, assertion):
  \NumDeepErrHoleSideSnapshots{} snapshots in
  \NumDeepErrHoleSideConfigs{} configurations (k28 at cutoffs $10^{-5}$ and
  $10^{-6}$), blasts \NumDeepErrHoleSideBlasts.
\item \textbf{Several outer faces} (\code{DualGraph}, a \code{ValueError}
  ``component \ldots{} faces of positive traversal area''):
  \NumDeepErrOuterFacesSnapshots{} snapshot, k28 at cutoff $10^{-6}$, blast
  \NumDeepErrOuterFacesBlasts.
\item \textbf{Bridge approaches its two crossings from opposite sides}
  (\code{check\_bridge\_\allowbreak rows\_\allowbreak consistent}, assertion):
  \NumDeepErrBridgeRowsSnapshots{} snapshots in
  \NumDeepErrBridgeRowsConfigs{} configurations, blasts
  \NumDeepErrBridgeRowsBlasts: p3 and nested\_inner at cutoff $10^{-5}$ from
  blast \NumDeepCutoffFiveOutcomePThree{} on (the reference cell's burst;
  these cells' own words had changed earlier), and with two iterations per
  blast near the end of the run.
\end{enumerate}

\begin{table}[htbp]
\centering\footnotesize
\caption{E11 topology errors, per error and case. These are genuine errors;
the blasting continued after each, and no crossing or cdist law was
violated. ``snapshots''
counts the blasts at which the error was raised.}
\label{tab:e11errors}
\setlength{\tabcolsep}{3pt}
\begin{adjustbox}{max width=\textwidth}
\begin{tabular}{lllllrr}
\tablebody tables/e11_errors.tex 
\end{tabular}
\end{adjustbox}
\end{table}

\reading All four errors live in the topology layer, appear only after
several blasts, and none comes with a crossing or cdist violation, so they
are not caused by the numerics producing impossible crossings. They depend on
the blasting cutoff, but that alone does not show they are geometric. E7's
k28 run, built by the same recipe, reached \NumBlastCrossingsEndKTwoEight{}
crossings after \NumMaxBlastsKTwoEight{} blasts and its topology completed
(\NumBlastClassesKTwoEight{} classes, provisional). The E11 k28 reference at
the same blast also has \NumDeepAtESevenCrossingsKTwoEight{} crossings, yet
it raises the zero-width assertion, as do all k28 cells at cutoff
$10^{-7}$. So the zero-width error may depend on the E11 snapshot path (a
topology stack after every blast, session state left by an assertion
mid-stack, or the extra \code{bridge\_classes}/\code{minimal\_trellis}/\allowbreak
\code{iterated\_partition} calls) rather than on the trellis; it has to be
reproduced on a fresh session at that blast before it is attributed to
geometry. My working hypotheses, to be checked by the topology owner:
(1)~an image bridge's endpoint lands exactly on an existing element
boundary, giving an element $[x,x]$, and the parent containment test does
not accept a degenerate element at a boundary;
(2)~the propagated-hole side rule, fixed on 2026-10-02 for the period-3
case, has another corner case on deeply blasted trellises;
(3)~the arrangement's outer-face detection meets a component with
near-degenerate faces;
(4)~the row check trips where two crossing signs come out equal, which the
row rule reads as a bridge crossing its stable branch. At a coarse cutoff
in a folded region that may be a resolution effect. In the two-iteration
cells the lobe-area sign flips fall exactly on the first blasts at which
this error is raised, which points at the geometry of those bridges rather
than at the row rule.

\caveat Each error's traceback is stored only at its first occurrence, and
a snapshot that raised has no counts. Whether an assertion is a bug or an
over-strict invariant is not decided here.

\paragraph{What to investigate next.}
\begin{itemize}[leftmargin=*]
\item Reproduce the zero-width element at k28's first failing blast from the
  stored configuration, on a fresh session that runs the topology stack only
  once (E7's k28 run completes at the same crossing count); decide whether degenerate iterated elements are legal (as
  singletons) and make the containment test closed at boundaries.
\item Re-run the bridge-row and hole-side cells with the crossing signs and
  the offending bridge dumped, to separate a resolution effect from a rule
  defect.
\item Make blasting memory-aware in the library, the way the harness is.
\item Follow p3's persistent unresolved class: grow the stable manifold
  instead of blasting again, and see whether it resolves.
\end{itemize}

%======================================================================
\section{Fixed-point solver (E8)}
\label{sec:solver}

\claim The solver is fast and accurate when it finds the orbit you asked
for. It finds the full set of short orbits of the $k=10$ horseshoe. It does
\emph{not} check the minimal period, so it regularly hands back a
lower-period orbit as if it were period $p$.

\evidence \NumSolverSolves{} solves (\cref{fig:e08out}, \cref{tab:e08} in
the appendix): \NumSolverConvergedRate{} converged and
\NumSolverNoConvergenceRate{} did not. \NumSolverSaddleRate{} of all solves
gave a saddle of the requested minimal period. The maximum orbit residual is
\NumSolverMaxResidual{}, and the median solve takes \NumSolverMedianMs\,ms
(99th percentile \NumSolverPNinetyNineMs\,ms, \cref{fig:e08time}). At $k=10$
the distinct orbits found equal the full 2-shift count for
\NumHorseshoeMatched{} periods (\NumHorseshoeFound{} of
\NumHorseshoeExpected{} orbits, \cref{fig:e08orb}(a)).

\begin{figure}[htbp]
  \centering
  \includegraphics[width=0.88\textwidth]{\fig{e08_solver_outcomes}}
  \caption{E8 outcomes per map and requested period. Orange (``converged to
  a lower period'') is the trap; the all-orange period-2 bars at $k<3$ are
  maps with no period-2 orbit at all.}
  \label{fig:e08out}
\end{figure}

\begin{figure}[htbp]
  \centering
  \includegraphics[width=0.88\textwidth]{\fig{e08_orbits}}
  \caption{E8. (a)~Distinct orbits found: at $k=10$ the curve lies exactly
  on the 2-shift count. (b)~Fraction of converged solves with a lower
  minimal period, per requested period.}
  \label{fig:e08orb}
\end{figure}

\reading The headline wrong-period rate, \NumSolverWrongPeriodRate{} of
converged solves, mixes two different things. \NumSolverLowerNoOrbit{}
solves (\NumSolverLowerNoOrbitRate) asked for period 2 at $k<3$. There is
no period-2 orbit there: none was found from any seed, consistent with its
birth in the period doubling of the elliptic point at $k=3$. So falling
back to a fixed point is the only thing a shooting solver can do. The other
\NumSolverLowerSlide{} solves (\NumSolverLowerSlideRate) are genuine
slides: the requested period exists, and the solver found a fixed point or
a divisor-period orbit instead. At $k=10$, where every period exists, the
rate is \NumSolverLowerSlideRateKTen. The real problem is downstream:
\NumSolverLowerAccepted{} lower-period solves passed
\code{construct\_fixed\_point}, which built a saddle with the
\emph{requested} period (e.g.\ a period-1 point labelled period 2 with
$k_{\rm value}=2$). The lower-period solves that were rejected were rejected
for being elliptic, not for their period.

\begin{figure}[htbp]
  \centering
  \includegraphics[width=0.88\textwidth]{\fig{e08_solve_time}}
  \caption{E8 solve time. (a)~Distribution: sub-millisecond for most
  solves. (b)~Median and IQR by period: cost grows mildly with the period.}
  \label{fig:e08time}
\end{figure}

\caveat Seeds are uniform in a box sized to each map. Basins elsewhere are
not sampled, so the orbit counts at $k<3$ are lower bounds.

%======================================================================
\section{Reproducibility (E10)}
\label{sec:repro}

\claim Runs are deterministic on a given device, and CPU and GPU agree to
floating-point rounding.

\evidence \Cref{tab:e10} and \cref{fig:e10}: across \NumRepeatsCpu{}
fresh-process repeats per device, CPU--CPU is \NumBitIdentCpuCpu{} and
GPU--GPU is \NumBitIdentGpuGpu{} (sha256 of every point array). CPU--GPU
\NumBitIdentCpuGpu, and still \NumBitIdentCpuGpuRounded{} after rounding to
$10^{-12}$. The crossing coordinates agree to \NumCrossingDevCpuGpu. The CV
of the summed stage time is \NumTimingCvCpuTotal{} (CPU) and
\NumTimingCvGpuTotal{} (GPU); per stage it is at most \NumTimingCvCpuMax{}
and \NumTimingCvGpuMax. In this case the GPU run is slower, at
\NumDetGpuSpeedup{} of the CPU speed (\NumDetTotalCpuS\,s against
\NumDetTotalGpuS\,s).

\begin{table}[htbp]
\centering\footnotesize
\caption{E10: do the sha256 hashes agree across repeats (exact bytes, and
rounded to $10^{-12}$)?}
\label{tab:e10}
\resizebox{\textwidth}{!}{%
\begin{tabular}{lcccccc}
\tablebody tables/e10_identity.tex 
\end{tabular}}
\end{table}

\begin{figure}[htbp]
  \centering
  \includegraphics[width=0.88\textwidth]{\fig{e10_determinism}}
  \caption{E10. (a)~Stage times over five repeats, labelled with their CV:
  the spread is a few percent. (b)~Total time per repeat: no drift, and
  the GPU is consistently slower on this mid-size case.}
  \label{fig:e10}
\end{figure}

\reading The CPU--GPU differences are rounding differences in the map that
scale with the coordinate magnitude. In the lockstep comparison
(\cref{fig:e02dev}) they first appear at step
\NumGpuFirstDivergentStepKTwoEight. At step 11 of k28 they reach
\NumGpuVerifyDevStepElevenKTwoEight{} in absolute terms but only
\NumGpuVerifyRelDevStepElevenKTwoEight{} relative to the largest coordinate.
At step 12, with coordinates near \NumGpuVerifyMaxCoordKTwoEight, they flip
refinement decisions: the CPU holds \NumGpuVerifyPointsCpuKTwoEight{} points
and the GPU \NumGpuVerifyPointsGpuKTwoEight. The crossings, which live near
the saddle, are unaffected.

\begin{figure}[htbp]
  \centering
  \includegraphics[width=0.88\textwidth]{\fig{e02_deviation}}
  \caption{E2 lockstep CPU/GPU deviation of the unstable manifolds.
  (a)~Absolute: it grows once the arm escapes. (b)~Relative to the largest
  coordinate: it stays at machine $\epsilon$. The cross marks the step where
  the point counts differ.}
  \label{fig:e02dev}
\end{figure}

\caveat Five repeats bound the CV loosely. A timing difference below
$\sim$5\% between two configurations should not be read as real.

%======================================================================
\section{Failure taxonomy}
\label{sec:taxonomy}

\Cref{tab:taxonomy} collects every failure class the campaign hit.
``By design'' means the library refuses the input on purpose; ``limitation''
means correct behaviour within an unstated envelope; ``bug'' means behaviour
the library's own contract rules out. Topology changing under blasting or
deeper growth is not a failure and does not appear here. Neither do the
symbolic layer's provisional flags (\cref{sec:setup}); both are listed
after the table.

\begin{table}[p]
\centering\footnotesize
\caption{Failure taxonomy.}
\label{tab:taxonomy}
\setlength{\tabcolsep}{3pt}
\begin{tabularx}{\textwidth}{@{}>{\raggedright\arraybackslash}p{2.3cm} Y Y >{\raggedright\arraybackslash}p{1.6cm} Y@{}}
\toprule
failure class & trigger & symptom & kind & suggested fix \\
\midrule
Escape OOM & fixed step count past the tangle (\cref{sec:escape,sec:scaling}) & unbudgeted growth: one step goes from MB to the RSS cap; OOM kill & genuine error (OOM) & point/memory budget in \code{grow\_n\_times}; stop-condition growth by default \\
Inversion blow-up & inversion saddle at \NumInversionOomSteps{} unstable steps; $k\ge\NumSweepOtherOomKMin$ in E6 & unbudgeted growth: OOM, documented as ``does not finish'' & genuine error (OOM) & same budget guard \\
Artefact crossings & near tangency, $k+1<\NumTangencyAgreeKPlusOne$ (\cref{sec:tangency}) & crossing count set by cutoff; no warning & limitation (silent) & angle/invariance-error guard; resolution warning \\
Self-crossing polygon & coarse cutoff, e.g.\ \NumIsectUUCutoff{} (\cref{sec:intersections}) & u$\times$u crossings undetected & limitation (silent) & cheap u$\times$u diagnostic \\
Coarse-cutoff crossing errors & default $10^{-4}$ on p3 (\cref{sec:accuracy}) & crossings invented and missed & limitation & default $10^{-7}$ \\
Wrong period & \code{compute\_fixed\_point} from an arbitrary seed (\cref{sec:solver}) & lower-period orbit accepted as period $p$ & bug & minimal-period check in \code{construct\_fixed\_point} \\
Max iterations & $k\le\NumSweepMaxIterKMax$, \code{grow\_until\_turnaround} (\cref{sec:robustness}) & \code{ValueError} & limitation & raise cap or change stop condition at small $\lambda_u$ \\
Trivial tangle & fixed recipe at small $k$ (\cref{sec:robustness}) & completes with 0 crossings & limitation & scale depth with $1/\ln\lambda_u$ \\
Not a saddle & elliptic point & \code{ValueError} & by design & none \\
Orientation reversing & $b=-1$ & \code{ValueError} & by design & none (out of scope) \\
Partition assert & inversion at \NumInversionAssertSteps{} steps (\cref{sec:depth}) & \code{AssertionError} ``owned by 0 partition elements'' & genuine error; bug or too strict (topology) & topology owner to triage \\
Zero-width iterated element & k28, k32 at blasts \NumDeepErrZeroWidthBlasts, cutoff dependent (\cref{sec:blastdepth-errors}) & \code{AssertionError} ``not inside its parent'' in \code{PartitionFamily.\allowbreak from\_\allowbreak minimal} & genuine error (topology) & decide whether $[x,x]$ elements are legal; closed containment test \\
Hole side & k28 at cutoffs $10^{-5}$, $10^{-6}$, blasts \NumDeepErrHoleSideBlasts{} & \code{AssertionError} ``disagree on bridge\_side'' & genuine error (topology) & dump the hole pair; recheck the propagated-side rule \\
Bridge rows & p3, nested\_inner at cutoff $10^{-5}$ or two iterations, blasts \NumDeepErrBridgeRowsBlasts{} & \code{AssertionError} ``opposite sides'' (\code{check\_bridge\_\allowbreak rows\_\allowbreak consistent}) & genuine error (topology) & dump the crossing signs; separate resolution effect from rule defect \\
Several outer faces & k28 at cutoff $10^{-6}$, blast \NumDeepErrOuterFacesBlasts{} & \code{ValueError} in \code{DualGraph} & genuine error (topology) & inspect the arrangement component \\
Stale registry & caching \code{intersection\_registry} across \code{compute\_intersections} & empty/stale object & API trap & return or re-fetch; warn in docs \\
\bottomrule
\end{tabularx}
\end{table}

\paragraph{Expected behaviour (not failures).} Blasting and deeper growth
change the topology: k28's crossings went from
\NumBlastCrossingsStartKTwoEight{} to \NumBlastCrossingsEndKTwoEight{} over
\NumMaxBlastsKTwoEight{} blasts, and p3's bridge classes rise with every
step past \NumDepthReliableMaxPThree{} (\cref{sec:depth,sec:blasts}). In
E11, \NumDeepStructureSteps{} snapshots brought new canonical words, and
nested\_outer's classes even fell from \NumDeepDropBeforeNestedOuter{} to
\NumDeepDropAfterNestedOuter{} (\cref{sec:blastdepth-evolution}). That is
the tangle being uncovered. The memory and time budget stops of E11 are the
experiment's own guard, not failures.

\paragraph{Provisional symbolic results (open, not errors).} The
library's \code{is\_reliable} flag is False for p3 from
\NumDepthUnreliableMinPThree{} steps, for the inversion saddle at
\NumDepthNumericsMaxInversion{} steps, and after the last blast on k28
(\NumBlastUnresolvedKTwoEight{} unresolved, \NumBlastVirtualKTwoEight{}
virtual classes) and p3 (\NumBlastUnresolvedPThree{} unresolved,
\NumBlastVirtualPThree{} virtual). In E11 the flags come and go: the first
blast clears them on k28, k32 and nested\_outer, and they return at the
class burst on p3 and nested\_inner and at or before it on k28 and k32;
nested\_outer's never returns (\cref{sec:blastdepth-provisional}). Before
blasting the unresolved classes lack a registered image; after the bursts
they are mostly dual-graph walks that cannot complete. Both call for more
growth or a further rule. The arrangement's \code{image\_of} misses
(\NumArrImageMissing{} of the regions tried, \cref{sec:invariants}) are an
open question of the same kind, not a failure.

%======================================================================
\section{Recommendations}
\label{sec:recs}

\paragraph{Defaults.}
\begin{itemize}[leftmargin=*]
\item \code{area\_cutoff}: $10^{-7}$. It is converged in every case tested,
  and its cost over $10^{-4}$ is the point ratio
  \NumPointsKTwoEightCutoffFour{} $\to$ \NumPointsKTwoEightCutoffSeven{} on
  k28 (\cref{sec:accuracy}).
\item Growth: stop-condition growth (\code{grow\_until\_*}) as the
  documented default. \code{grow\_n\_times} only with a budget.
\item GPU: off. The measured ceiling is \NumAmdahlBoundMax{}
  (\cref{sec:gpu}).
\end{itemize}

\paragraph{Guards to add (cheap, numerics layer).}
\begin{itemize}[leftmargin=*]
\item \textbf{Self-crossing diagnostic}: test non-adjacent u$\times$u (and
  s$\times$s) segment pairs through the existing R-tree after growth, and
  warn with the cutoff that would remove them.
\item \textbf{Angle/resolution warning}: for each registered crossing,
  compare $|\sin\theta|$ times the local segment length with the local
  invariance error, and warn when the crossing cannot be distinguished from
  polygon noise.
\item \textbf{Minimal-period check}: in \code{construct\_fixed\_point},
  reject (or relabel) an orbit whose minimal period is a proper divisor of
  the requested one.
\item \textbf{Memory guard}: predict the next step's points from the last
  growth factor, and refuse or warn above a budget. The E3 harness already
  does this in a few lines.
\item \textbf{Memory-aware blasting}: the same guard for \code{blast\_zone},
  since a large zone multiplies its points by about
  \NumDeepPointGrowthKTwoEight{} per blast (\cref{sec:blastdepth-cost}). The
  E11 harness's memory budget is a template.
\end{itemize}

\paragraph{Guards to add (topology layer, for its owner).}
\begin{itemize}[leftmargin=*]
\item \textbf{Zero-width iterated elements}: either admit a degenerate
  element $[x,x]$ as a singleton or merge it into a neighbour before the
  parent-containment assertion in \code{PartitionFamily.\allowbreak from\_\allowbreak minimal}
  (\cref{sec:blastdepth-errors}).
\item \textbf{Diagnostics on assertion}: have the row and hole-side checks
  report the offending crossings and their signs, so that a resolution
  effect can be told from a rule defect.
\end{itemize}

\paragraph{Open questions for the advisor.}
\begin{enumerate}[leftmargin=*]
\item Near tangency, should we accept that crossings below a resolution
  threshold are undecidable with polygons, or invest in a validated method
  (interval arithmetic, or parameterisation-method manifolds)?
\item Is \code{image\_of} expected to fail for regions whose image crossings
  are not registered (\NumArrImageMissing{} of the regions tried), or is
  this a gap?
\item Should dissipative maps ($b\ne\pm1$) be in scope? They currently run
  without complaint and without validation.
\item Is the inversion partition assertion at shallow depth a real defect or
  an over-strict invariant?
\item Do the unresolved and virtual classes after many blasts resolve with
  further growth? In p3 and nested\_inner they persist through every blast
  after the burst (\cref{sec:blastdepth-provisional}), so further blasting of
  the same zone does not resolve them; stable growth is untested.
\item Which blasting cutoff should define the symbolic result? In most
  cutoff cells the words diverge before the crossing counts do
  (\cref{sec:blastdepth-params}).
\item Is the lobe-area excess in the two-iteration E11 cells
  ($|A'/A-1|$ up to \NumDeepLobeRatioMax, \NumDeepLobeSignFlips{} sign
  flips) a self-intersecting probe polygon or a real area defect
  (\cref{sec:blastdepth-errors})?
\end{enumerate}

%======================================================================
\section{Threats to validity}
\label{sec:threats}

\begin{itemize}[leftmargin=*]
\item \textbf{One machine.} All timings are from a single host and one GPU.
  The GPU conclusion rests on Amdahl's law, which transfers to other
  hardware; the absolute crossover sizes do not.
\item \textbf{Hénon only.} Every result is for the Hénon family. Maps with
  more expensive evaluations would raise $f_{\rm map}$ and could change the
  GPU verdict.
\item \textbf{Quick calibration of grids.} The sweep grids and depths were
  calibrated in a short pilot to fit the budget. Exponents rest on
  $\le$3 decades, and some limits (escape and depth OOM, the tangency
  transition) are only known to the resolution of the grid. The E7 blast
  study ran one cutoff per case, with the topology probed only after the
  last blast.
\item \textbf{E11 grid thinned and budgets.} E11 moved one parameter at a
  time around a reference cell; the full cutoff $\times$ separation product
  and repeats of the heavy k32 and nested\_outer cells were dropped. Its
  heavy cells were stopped by their own memory or time budget
  (\NumDeepStopMemoryBudget{} and \NumDeepStopBudget{} configs), so the deep
  behaviour of the large zones beyond \NumDeepBlastsKTwoEight{} blasts is unknown, and
  cutoffs are compared at different final depths. Snapshots that raised
  leave gaps in the series.
\item \textbf{cProfile overhead.} Profiling inflates Python call costs
  relative to NumPy kernels, which biases $f_{\rm map}$ low (see the caveat
  in \cref{sec:gpu}).
\item \textbf{RSS sampling.} The watchdog samples every 50\,ms, so peak RSS
  can be missed by a short spike, and OOM kills happen up to one sample
  after the cap. Caps differ per experiment (\cref{tab:counts}).
\item \textbf{Invariance probe design.} The probe maps points sampled inside
  segments, because a vertex's image is often exactly a vertex, and measures
  one-step consistency. It is not a distance to the true manifold.
\item \textbf{References are finest runs.} Convergence in E4 is measured
  against $10^{-9}$, and E9's ``truth'' is a brute force on the \emph{same}
  polygons. E9 therefore validates the intersection code, not the manifold
  approximation.
\end{itemize}

%======================================================================
\clearpage
\appendix
\section{Tables}
\label{app:tables}

\begin{table}[H]
\centering\footnotesize
\caption{E4: cost and accuracy per cutoff. ``crossing dev.'' is the median
distance of the finest run's crossings to this run's.}
\label{tab:e04}
\setlength{\tabcolsep}{3.5pt}
\begin{tabular}{llrlrrrrrr}
\tablebody tables/e04_cutoff.tex 
\end{tabular}
\end{table}

\begin{table}[H]
\centering\footnotesize
\caption{E8: solver outcomes per map and requested period. ``orbits found''
counts distinct orbits of exactly that minimal period over all solves (for
$k=10$: found / full 2-shift count).}
\label{tab:e08}
\begin{tabular}{lrrrrrr}
\tablebody tables/e08_solver.tex 
\end{tabular}
\end{table}

\section{Reproduction}
\label{app:repro}

From the repository root, with the project's environment:
\begin{verbatim}
env/bin/python stress/numerics/run_campaign.py --quick   # pilot
env/bin/python stress/numerics/run_campaign.py           # full run
env/bin/python stress/numerics/report/make_figures.py
env/bin/python stress/numerics/report/make_numbers.py
env/bin/python stress/numerics/report/make_numbers_extra.py
cd stress/numerics/report && ~/.local/bin/tectonic report.tex
\end{verbatim}
Results land in \code{stress/numerics/results/*.jsonl} (one record per
configuration) with a running \code{progress.log}. A single experiment runs
on its own as
\begin{verbatim}
env/bin/python stress/numerics/eNN_*.py [--quick] [--only I,J] \
    [--rss-cap-gb G] [--timeout S] [--workers N]
\end{verbatim}
E11 is not part of \code{run\_campaign.py}; it runs on its own (8 workers,
up to 12\,GB each), before the report scripts above:
\begin{verbatim}
env/bin/python stress/numerics/e11_blast_depth.py           # full grid
env/bin/python stress/numerics/e11_blast_depth.py --quick   # smoke test
\end{verbatim}


\section{Configuration and outcome counts}
\label{app:counts}

\begin{table}[H]
\centering\footnotesize
\caption{Records and outcomes per experiment in the full run, with each
experiment's default per-configuration timeout and RSS cap (some E7 configs
override the timeout). ``assert'' counts \code{invariant\_assert} outcomes of
the numerics; the topology assertions of E7 and E11 are genuine errors
recorded inside \code{ok} records (E11's are in \cref{tab:e11errors}).}
\label{tab:counts}
\begin{tabular}{lrrrrrrrr}
\tablebody tables/campaign_counts.tex 
\end{tabular}
\end{table}

\end{document}
